\newMonth{October}
\setrunningtitles{October}{October}

% ---- martyrology/mart10/mart1001.htm
\needspace{10\baselineskip}
\begin{paracol}{2}
\selectlanguage{latin}
\begin{center}{\color{gregoriocolor} Kaléndis Octóbris. 
 Luna\dots\ }\end{center}
\switchcolumn
\selectlanguage{english}
\begin{center}{\color{gregoriocolor} The   1\textsuperscript{st} Day of 
 October. The\dots\ Day of the Moon.}\end{center}
\end{paracol}

\noindent\begin{tabularx}{\linewidth}{*{19}{>{\centering\arraybackslash}X}}
 \textcolor{gregoriocolor}{a} & \textcolor{gregoriocolor}{b} & \textcolor{gregoriocolor}{c} & \textcolor{gregoriocolor}{d} & \textcolor{gregoriocolor}{e} & \textcolor{gregoriocolor}{f} & \textcolor{gregoriocolor}{g} & \textcolor{gregoriocolor}{h} & \textcolor{gregoriocolor}{i} & \textcolor{gregoriocolor}{k} & \textcolor{gregoriocolor}{l} & \textcolor{gregoriocolor}{m} & \textcolor{gregoriocolor}{n} & \textcolor{gregoriocolor}{p} & \textcolor{gregoriocolor}{q} & \textcolor{gregoriocolor}{r} & \textcolor{gregoriocolor}{s} & \textcolor{gregoriocolor}{t} & \textcolor{gregoriocolor}{u} \\
 9 & 10 & 11 & 12 & 13 & 14 & 15 & 16 & 17 & 18 & 19 & 20 & 21 & 22 & 23 & 24 & 25 & 26 & 27 \\
\end{tabularx}
\vspace{0.5\baselineskip}
\noindent\begin{tabularx}{\linewidth}{*{12}{>{\centering\arraybackslash}X}}
 \textcolor{gregoriocolor}{A} & \textcolor{gregoriocolor}{B} & \textcolor{gregoriocolor}{C} & \textcolor{gregoriocolor}{D} & \textcolor{gregoriocolor}{E} & F & \textcolor{gregoriocolor}{F} & \textcolor{gregoriocolor}{G} & \textcolor{gregoriocolor}{H} & \textcolor{gregoriocolor}{M} & \textcolor{gregoriocolor}{N} & \textcolor{gregoriocolor}{P} \\
 28 & 29 & 1 & 2 & 3 & 4 & 3 & 4 & 5 & 6 & 7 & 8 \\
\end{tabularx}

\begin{paracol}{2}
\selectlanguage{latin}
\lettrine[lines=2]{S}{ancti} Remígii, 
 Epíscopi Rheménsis et Confessóris, qui Idibus Januárii obdormívit in Dómino, 
 sed hac die, ob Translatiónem córporis ejus, potíssimum cólitur.
\switchcolumn
\selectlanguage{english}
\lettrine[lines=2]{S}{t.} Remigius, bishop of Rheims and 
 confessor, who fell asleep in the Lord on the 13th of January, but is 
 commemorated on this day because of the translation of his body.
\switchcolumn*
\selectlanguage{latin}
Romæ beáti Arétæ 
 Mártyris, et aliórum quingentórum quátuor.
\switchcolumn
\selectlanguage{english}
At Rome, blessed Aretas and five 
 hundred and four other martyrs.
\switchcolumn*
\selectlanguage{latin}
Tornáci, in Gálliis, 
 sancti Piatónis, Presbyteri et Mártyris; qui, prædicatiónis causa, cum beáto 
 Quinctíno ejúsque Sóciis, ab urbe Roma in Gálliam perréxit, ac póstea, in 
 persecutióne Maximiáni, consummáto martyrio, migrávit ad Dóminum.
\switchcolumn
\selectlanguage{english}
At Tournai in France, St. Piaton, 
 priest and martyr, who went from Rome to France to preach, together with 
 blessed Quinctinus and his companions. Afterwards, his martyrdom was 
 completed in the persecution of Maximian and he passed from earth to heaven.
\switchcolumn*
\selectlanguage{latin}
Tomis, in Ponto, 
 sanctórum Mártyrum Prisci, Crescéntis et Evágrii.
\switchcolumn
\selectlanguage{english}
At Tomis in Pontus, the holy martyrs 
 Priscus, Crescens, and Evagrius.
\switchcolumn*
\selectlanguage{latin}
Ulyssipóne, in 
 Lusitánia, sanctórum Mártyrum Veríssimi, Máximæ et Júliæ, sorórum ejus; qui 
 in Diocletiáni Imperatóris persecutióne passi sunt.
\switchcolumn
\selectlanguage{english}
At Lisbon in Portugal, the holy 
 martyrs Verissimus, and his sisters Maxima and Julia, who suffered in the 
 persecution of Diocletian.
\switchcolumn*
\selectlanguage{latin}
Thessalonícæ sancti 
 Domníni Mártyris, sub Maximiáno Imperatóre.
\switchcolumn
\selectlanguage{english}
At Thessalonica, St. Domninus, 
 martyr, under Emperor Maximian.
\switchcolumn*
\selectlanguage{latin}
Urbe Véteri sancti 
 Sevéri, Presbyteri et Confessóris.
\switchcolumn
\selectlanguage{english}
At Orvieto, St. Severus, priest and 
 confessor.
\switchcolumn*
\selectlanguage{latin}
In Portu Gandæ sancti 
 Bavónis Confessóris.
\switchcolumn
\selectlanguage{english}
At the port of Ghent, St. Bavo, 
 confessor.
\switchcolumn*
\selectlanguage{latin}
\end{paracol}

% \continueMonth{Octobver}

% ---- martyrology/mart10/mart1002.htm
\needspace{10\baselineskip}
\begin{paracol}{2}
\selectlanguage{latin}
\begin{center}{\color{gregoriocolor} Sexto Nonas Octóbris. 
 Luna\dots\ }\end{center}
\switchcolumn
\selectlanguage{english}
\begin{center}{\color{gregoriocolor} The   2\textsuperscript{nd} Day of 
 October. The\dots\ Day of the Moon.}\end{center}
\end{paracol}

\noindent\begin{tabularx}{\linewidth}{*{19}{>{\centering\arraybackslash}X}}
 \textcolor{gregoriocolor}{a} & \textcolor{gregoriocolor}{b} & \textcolor{gregoriocolor}{c} & \textcolor{gregoriocolor}{d} & \textcolor{gregoriocolor}{e} & \textcolor{gregoriocolor}{f} & \textcolor{gregoriocolor}{g} & \textcolor{gregoriocolor}{h} & \textcolor{gregoriocolor}{i} & \textcolor{gregoriocolor}{k} & \textcolor{gregoriocolor}{l} & \textcolor{gregoriocolor}{m} & \textcolor{gregoriocolor}{n} & \textcolor{gregoriocolor}{p} & \textcolor{gregoriocolor}{q} & \textcolor{gregoriocolor}{r} & \textcolor{gregoriocolor}{s} & \textcolor{gregoriocolor}{t} & \textcolor{gregoriocolor}{u} \\
 10 & 11 & 12 & 13 & 14 & 15 & 16 & 17 & 18 & 19 & 20 & 21 & 22 & 23 & 24 & 25 & 26 & 27 & 28 \\
\end{tabularx}
\vspace{0.5\baselineskip}
\noindent\begin{tabularx}{\linewidth}{*{12}{>{\centering\arraybackslash}X}}
 \textcolor{gregoriocolor}{A} & \textcolor{gregoriocolor}{B} & \textcolor{gregoriocolor}{C} & \textcolor{gregoriocolor}{D} & \textcolor{gregoriocolor}{E} & F & \textcolor{gregoriocolor}{F} & \textcolor{gregoriocolor}{G} & \textcolor{gregoriocolor}{H} & \textcolor{gregoriocolor}{M} & \textcolor{gregoriocolor}{N} & \textcolor{gregoriocolor}{P} \\
 29 & 1 & 2 & 3 & 4 & 5 & 4 & 5 & 6 & 7 & 8 & 9 \\
\end{tabularx}

\begin{paracol}{2}
\selectlanguage{latin}
\lettrine[lines=1]{F}{estum} sanctórum Angelórum Custódum.
\switchcolumn
\selectlanguage{english}
\lettrine[lines=1]{T}{he} Feast of the holy Guardian Angels.
\switchcolumn*
\selectlanguage{latin}
Romæ pássio sancti 
 Modésti Sardi, Levítæ et Mártyris; qui, sub Diocletiáno Imperatóre, equúleo 
 tortus atque igne adústus est. Ipsíus vero corpus, Benevéntum póstea 
 translátum, in Ecclésia suo insigníta nómine collocátum fuit.
\switchcolumn
\selectlanguage{english}
At Rome, the martyrdom of St. 
 Modestus, a Sardinian, deacon and martyr, who was racked and burned with 
 fire by Emperor Diocletian. His holy body was afterwards translated to 
 Benevento and buried there in a church named after him.
\switchcolumn*
\selectlanguage{latin}
In território 
 Atrebaténsi item pássio beáti Leodegárii, Augustodunénsis Epíscopi; quem, 
 váriis injúriis et divérsis supplíciis pro veritáte afflíctum, Ebroínus, 
 Major domus régiæ Theodoríci, intérfici jussit.
\switchcolumn
\selectlanguage{english}
In the vicinity of Arras, the 
 martyrdom of blessed Leodegar, bishop of Autun. After being 
 subjected to various insults and torments for the truth, he was put to death 
 by Ebroin, chief minister of Theodoric.
\switchcolumn*
\selectlanguage{latin}
Nicomedíæ sancti 
 Eleuthérii, mílitis et Mártyris, cum áliis innúmeris; qui, cum Diocletiáni 
 régia incéndio conflagrásset et falso hujus críminis essent accusáti, omnes, 
 jubénte eódem sævíssimo Imperatóre, acervátim necáti sunt. Horum porro 
 álii gládiis obtruncabántur, álii cremabántur ígnibus, álii in mare 
 præcipitabántur; sed inter eos primus Eleuthérius, cum per síngula torménta, 
 diu cruciátus, valídior redderétur, martyrium victóriæ suæ, ígnibus 
 velut aurum examinátus, complévit.
\switchcolumn
\selectlanguage{english}
At Nicomedia, St. Eleutherius, 
 soldier and martyr, with innumerable others. They were falsely accused 
 of having set fire to the palace of Diocletian and, by order of this cruel 
 emperor, were barbarously massacred in groups. Some were put to the 
 sword, some consumed by fire, while others were cast into the sea. But 
 the principal one, Eleutherius, after long tortures, being found stronger 
 after each torment, completed his victorious martyrdom by fire, as 
 well-tried gold.
\switchcolumn*
\selectlanguage{latin}
Antiochíæ sanctórum 
 Mártyrum Primi, Cyrílli et Secundárii.
\switchcolumn
\selectlanguage{english}
At Antioch, the holy martyrs Primus, 
 Cyril, and Secundarius.
\switchcolumn*
\selectlanguage{latin}
Eódem die sancti Geríni 
 Mártyris, qui frater éxstitit beáti Leodegárii, Augustodunénsis Epíscopi, 
 et, jubénte ipso Ebroíno, lapídibus óbrutus est.
\switchcolumn
\selectlanguage{english}
On the same day, St. Gerinus, 
 martyr, brother of blessed Leodegar, bishop of Autun. He was stoned 
 to death by the same Ebroin.
\switchcolumn*
\selectlanguage{latin}
Constantinópoli sancti 
 Theóphili Mónachi, qui pro defensióne sanctárum Imáginum a Leóne Isáurico 
 sævíssime cæsus et in exsílium pulsus, migrávit ad Dóminum.
\switchcolumn
\selectlanguage{english}
At Constantinople, St. Theophilus, a 
 monk. He was cruelly scourged by Leo the Isaurian for his defense of 
 holy images, was driven into exile, and there went gloriously to heaven.
\switchcolumn*
\selectlanguage{latin}
\end{paracol}


% ---- martyrology/mart10/mart1003.htm
\needspace{10\baselineskip}
\begin{paracol}{2}
\selectlanguage{latin}
\begin{center}{\color{gregoriocolor} Quinto Nonas Octóbris. 
 Luna\dots\ }\end{center}
\switchcolumn
\selectlanguage{english}
\begin{center}{\color{gregoriocolor} The   3\textsuperscript{rd} Day of 
 October. The\dots\ Day of the Moon.}\end{center}
\end{paracol}

\noindent\begin{tabularx}{\linewidth}{*{19}{>{\centering\arraybackslash}X}}
 \textcolor{gregoriocolor}{a} & \textcolor{gregoriocolor}{b} & \textcolor{gregoriocolor}{c} & \textcolor{gregoriocolor}{d} & \textcolor{gregoriocolor}{e} & \textcolor{gregoriocolor}{f} & \textcolor{gregoriocolor}{g} & \textcolor{gregoriocolor}{h} & \textcolor{gregoriocolor}{i} & \textcolor{gregoriocolor}{k} & \textcolor{gregoriocolor}{l} & \textcolor{gregoriocolor}{m} & \textcolor{gregoriocolor}{n} & \textcolor{gregoriocolor}{p} & \textcolor{gregoriocolor}{q} & \textcolor{gregoriocolor}{r} & \textcolor{gregoriocolor}{s} & \textcolor{gregoriocolor}{t} & \textcolor{gregoriocolor}{u} \\
 11 & 12 & 13 & 14 & 15 & 16 & 17 & 18 & 19 & 20 & 21 & 22 & 23 & 24 & 25 & 26 & 27 & 28 & 29 \\
\end{tabularx}
\vspace{0.5\baselineskip}
\noindent\begin{tabularx}{\linewidth}{*{12}{>{\centering\arraybackslash}X}}
 \textcolor{gregoriocolor}{A} & \textcolor{gregoriocolor}{B} & \textcolor{gregoriocolor}{C} & \textcolor{gregoriocolor}{D} & \textcolor{gregoriocolor}{E} & F & \textcolor{gregoriocolor}{F} & \textcolor{gregoriocolor}{G} & \textcolor{gregoriocolor}{H} & \textcolor{gregoriocolor}{M} & \textcolor{gregoriocolor}{N} & \textcolor{gregoriocolor}{P} \\
 1 & 2 & 3 & 4 & 5 & 6 & 5 & 6 & 7 & 8 & 9 & 10 \\
\end{tabularx}

\begin{paracol}{2}
\selectlanguage{latin}
\lettrine[lines=2]{S}{anctæ} Terésiæ a Jesu 
 Infánte, ex Ordine Carmelitárum Excalceatórum, Vírginis, peculiáris ómnium 
 Missiónum Patrónæ; cujus dies natális prídie Kaléndas Octóbris recensétur.
\switchcolumn
\selectlanguage{english}
\lettrine[lines=2]{S}{t.} Theresa of the Child Jesus, 
 virgin of the Order of Discalced Carmelites, special patroness of all 
 missions. Her birthday is commemorated on the 30th day of September.
\switchcolumn*
\selectlanguage{latin}
Romæ, ad Ursum pileátum, 
 sancti Cándidi Mártyris.
\switchcolumn
\selectlanguage{english}
At Rome, near the place called Ursus 
 Pileatus, St. Candidus, martyr.
\switchcolumn*
\selectlanguage{latin}
Apud antíquos Sáxones 
 sanctórum Mártyrum duórum Ewaldórum, qui, cum essent Presbyteri et Christum 
 ibi prædicáre cœpíssent, comprehénsi sunt a Pagánis et occísi; ad quorum 
 córpora noctu lux multa, diu appárens, et ubi essent et cujus essent mériti, 
 declarávit.
\switchcolumn
\selectlanguage{english}
Among the ancient Saxons, two holy 
 martyrs of the name of Ewald, priests who had been preaching in that 
 country. They were seized by the pagans and put to death. During 
 the night, a great light shone over the bodies for a long time, pointing out 
 where they were and also how distinguished were their merits.
\switchcolumn*
\selectlanguage{latin}
Eódem die sanctórum 
 Mártyrum Dionysii, Fausti, Caji, Petri, Pauli et aliórum quátuor, qui primum, 
 sub Décio, multa passi sunt, ac demum, sub Valeriáno, ab Æmiliáno Præside 
 diu torméntis vexáti, martyrii palmam meruérunt.
\switchcolumn
\selectlanguage{english}
Also, the holy martyrs Denis, 
 Faustus, Caius, Peter, Paul, and four others who had suffered greatly under 
 Decius. In the time of Valerian, they were long subjected to torments 
 by the governor Aemilian, and merited the palm of martyrdom.
\switchcolumn*
\selectlanguage{latin}
In Africa sancti 
 Maximiáni, Epíscopi Bagajénsis, qui, a Donatístis íterum atque íterum 
 sævíssima perpéssus, ex alta dénique turri præcipitátus est, et pro mórtuo 
 derelíctus; sed, póstmodum a transeúntibus colléctus et pia curatióne 
 sanátus, cathólicam fidem propugnáre non déstitit, donec, glória 
 confessiónis clarus, quiévit in Dómino.
\switchcolumn
\selectlanguage{english}
In Africa, St. Maximian, bishop of 
 Bagaia. Again and again he suffered great cruelties from the Donatists, 
 was finally cast headlong from a high tower, and left for dead. He was 
 found by passers-by, and having been healed by their zealous care, he did 
 not cease to defend the Catholic faith until he rested in the Lord, renowned 
 for the glory of his witness to the faith.
\switchcolumn*
\selectlanguage{latin}
Legióne, in Hispánia, 
 sancti Froiláni, ejúsdem civitátis Epíscopi, monásticæ vitæ propagándæ 
 stúdio, beneficéntia in páuperes, ceterísque virtútibus et miráculis clari.
\switchcolumn
\selectlanguage{english}
At Leon in Spain, St. Froylan, 
 bishop of that city, noted for his zeal in spreading the monastic life, his 
 generosity to the poor and other virtues, and for his miracles.
\switchcolumn*
\selectlanguage{latin}
In diœcési Namurcénsi, 
 apud Belgas, sancti Gerárdi Abbátis.
\switchcolumn
\selectlanguage{english}
In Belgium, in the diocese of Namur, 
 St. Gerard, abbot.
\switchcolumn*
\selectlanguage{latin}
In Palæstína sancti 
 Hesychii Confessóris, qui fuit sancti Hilariónis discípulus et in 
 peregrinatióne sócius.
\switchcolumn
\selectlanguage{english}
In Palestine, St. Hesychius, 
 confessor, disciple of St. Hilarion, and the companion of his travels.
\switchcolumn*
\selectlanguage{latin}
Savónæ, in Ligúria, 
 sanctæ Maríæ Joséphæ Rosséllo, Institúti Filiárum Nostræ Dóminæ a 
 Misericórdia Fundatrícis, quam, caritátis opéribus præcláram, Pius Papa 
 Duodécimus sanctis Virgínibus adnumerávit.
\switchcolumn
\selectlanguage{english}
At Savona in Liguria, St. Maria 
 Giuseppe Rossello, foundress of the Daughters of our Lady of Mercy. 
 Renowned for her charitable works, Pope Pius XII placed her among the number 
 of holy virgins.
\switchcolumn*
\selectlanguage{latin}
\end{paracol}


% ---- martyrology/mart10/mart1004.htm
\needspace{10\baselineskip}
\begin{paracol}{2}
\selectlanguage{latin}
\begin{center}{\color{gregoriocolor} Quarto Nonas Octóbris. 
 Luna\dots\ }\end{center}
\switchcolumn
\selectlanguage{english}
\begin{center}{\color{gregoriocolor} The   4\textsuperscript{th} Day of 
 October. The\dots\ Day of the Moon.}\end{center}
\end{paracol}

\noindent\begin{tabularx}{\linewidth}{*{19}{>{\centering\arraybackslash}X}}
 \textcolor{gregoriocolor}{a} & \textcolor{gregoriocolor}{b} & \textcolor{gregoriocolor}{c} & \textcolor{gregoriocolor}{d} & \textcolor{gregoriocolor}{e} & \textcolor{gregoriocolor}{f} & \textcolor{gregoriocolor}{g} & \textcolor{gregoriocolor}{h} & \textcolor{gregoriocolor}{i} & \textcolor{gregoriocolor}{k} & \textcolor{gregoriocolor}{l} & \textcolor{gregoriocolor}{m} & \textcolor{gregoriocolor}{n} & \textcolor{gregoriocolor}{p} & \textcolor{gregoriocolor}{q} & \textcolor{gregoriocolor}{r} & \textcolor{gregoriocolor}{s} & \textcolor{gregoriocolor}{t} & \textcolor{gregoriocolor}{u} \\
 12 & 13 & 14 & 15 & 16 & 17 & 18 & 19 & 20 & 21 & 22 & 23 & 24 & 25 & 26 & 27 & 28 & 29 & 1 \\
\end{tabularx}
\vspace{0.5\baselineskip}
\noindent\begin{tabularx}{\linewidth}{*{12}{>{\centering\arraybackslash}X}}
 \textcolor{gregoriocolor}{A} & \textcolor{gregoriocolor}{B} & \textcolor{gregoriocolor}{C} & \textcolor{gregoriocolor}{D} & \textcolor{gregoriocolor}{E} & F & \textcolor{gregoriocolor}{F} & \textcolor{gregoriocolor}{G} & \textcolor{gregoriocolor}{H} & \textcolor{gregoriocolor}{M} & \textcolor{gregoriocolor}{N} & \textcolor{gregoriocolor}{P} \\
 2 & 3 & 4 & 5 & 6 & 7 & 6 & 7 & 8 & 9 & 10 & 11 \\
\end{tabularx}

\begin{paracol}{2}
\selectlanguage{latin}
\lettrine[lines=2]{A}{ssísii,} in Umbria, 
 natális sancti Francísci, Levítæ et Confessóris; qui trium Ordinum, 
 scílicet Fratrum Minórum, Páuperum Dominárum, ac Fratrum et Sorórum de 
 Pæniténtia Fundátor éxstitit. Ipsíus autem vitam, sanctitáte ac 
 miráculis plenam, sanctus Bonaventúra conscrípsit.
\switchcolumn
\selectlanguage{english}
\lettrine[lines=2]{A}{t} Assisi in Umbria, the birthday of 
 St. Francis, cleric and confessor, founder of three orders: the Friars 
 Minor, the Poor Clares, and the Brothers and Sisters of Penance. His 
 life, filled with holy deeds and miracles, were written by St. Bonaventure.
\switchcolumn*
\selectlanguage{latin}
Corínthi item natális 
 sanctórum Crispi et Caji, quorum méminit sanctus Paulus Apóstolus ad 
 Corínthios scribens.
\switchcolumn
\selectlanguage{english}
At Corinth, the birthday of the 
 Saints Crispus and Caius, who are mentioned by the apostle St. Paul in his 
 Epistle to the Corinthians.
\switchcolumn*
\selectlanguage{latin}
Athénis sancti 
 Hieróthei, qui fuit discípulus ipsíus beáti Pauli Apóstoli.
\switchcolumn
\selectlanguage{english}
At Athens, St. Hierotheos, disciple 
 of the blessed apostle Paul.
\switchcolumn*
\selectlanguage{latin}
Damásci sancti Petri, 
 Epíscopi et Mártyris; qui, accusátus apud Agarenórum Príncipem quod fidem 
 Christi docéret, ídeo, lingua et mánibus pedibúsque amputátis, cruci affíxus 
 martyrium consummávit.
\switchcolumn
\selectlanguage{english}
At Damascus, St. Peter, bishop and 
 martyr, who was accused before the king of the Agarenians of teaching the 
 faith of Christ. His tongue, hands, and feet were cut off, and being 
 fastened to a cross, his martyrdom was fulfilled.
\switchcolumn*
\selectlanguage{latin}
Alexandríæ sanctórum 
 Presbyterórum, et Diaconórum Caji, Fausti, Eusébii, Chærémonis, Lúcii et 
 Sociórum; ex quibus álii, in persecutióne Valeriáni, Mártyres facti sunt, 
 álii, Martyribus serviéntes, mercédem Mártyrum recepérunt.
\switchcolumn
\selectlanguage{english}
At Alexandria, the holy priests and 
 deacons Caius, Faustus, Eusebius, Chaeremon, Lucius, and their companions. 
 Some of them were martyred in the persecution of Valerian; others, for 
 serving the martyrs, received the reward of martyrs.
\switchcolumn*
\selectlanguage{latin}
In Ægypto sanctórum 
 Mártyrum Marci et Marciáni fratrum, et aliórum ferme innumerabílium 
 utriúsque sexus atque omnis ætátis; quorum álii post vérbera, álii post 
 divérsa géneris horríbiles cruciátus flammis tráditi, álii in mare 
 præcipitáti, nonnúlli cápite cæsi, plúrimi inédia consúmpti, álii partíbulis 
 affíxi, quidam étiam, cápite deórsum verso et pédibus in sublíme sublátis, 
 appénsi, beatíssimam martyrii corónam meruérunt.
\switchcolumn
\selectlanguage{english}
In Egypt, the holy martyrs Mark and 
 Marcian, brothers, and an almost countless number of both sexes and of all 
 ages, who merited the blessed crown of martyrdom, some after being scourged, 
 others when they had suffered horrible torment, and others after being 
 delivered to the flames. Some were cast into the sea, some others were 
 beheaded; many were starved to death; many were fastened to gibbets; and 
 others again were suspended by the feet with their heads downward.
\switchcolumn*
\selectlanguage{latin}
Bonóniæ sancti Petrónii, 
 Epíscopi et Confessóris; qui doctrína, miráculis et sanctitáte cláruit.
\switchcolumn
\selectlanguage{english}
At Bologna, St. Petronius, bishop 
 and confessor, celebrated for learning, miracles, and sanctity.
\switchcolumn*
\selectlanguage{latin}
Lutétiæ Parisiórum 
 sanctæ Aureæ Vírginis.
\switchcolumn
\selectlanguage{english}
At Paris, St. Aurea, virgin.
\switchcolumn*
\selectlanguage{latin}
\end{paracol}


% ---- martyrology/mart10/mart1005.htm
\needspace{10\baselineskip}
\begin{paracol}{2}
\selectlanguage{latin}
\begin{center}{\color{gregoriocolor} Tértio Nonas Octóbris. 
 Luna\dots\ }\end{center}
\switchcolumn
\selectlanguage{english}
\begin{center}{\color{gregoriocolor} The   5\textsuperscript{th} Day of 
 October. The\dots\ Day of the Moon.}\end{center}
\end{paracol}

\noindent\begin{tabularx}{\linewidth}{*{19}{>{\centering\arraybackslash}X}}
 \textcolor{gregoriocolor}{a} & \textcolor{gregoriocolor}{b} & \textcolor{gregoriocolor}{c} & \textcolor{gregoriocolor}{d} & \textcolor{gregoriocolor}{e} & \textcolor{gregoriocolor}{f} & \textcolor{gregoriocolor}{g} & \textcolor{gregoriocolor}{h} & \textcolor{gregoriocolor}{i} & \textcolor{gregoriocolor}{k} & \textcolor{gregoriocolor}{l} & \textcolor{gregoriocolor}{m} & \textcolor{gregoriocolor}{n} & \textcolor{gregoriocolor}{p} & \textcolor{gregoriocolor}{q} & \textcolor{gregoriocolor}{r} & \textcolor{gregoriocolor}{s} & \textcolor{gregoriocolor}{t} & \textcolor{gregoriocolor}{u} \\
 13 & 14 & 15 & 16 & 17 & 18 & 19 & 20 & 21 & 22 & 23 & 24 & 25 & 26 & 27 & 28 & 29 & 1 & 2 \\
\end{tabularx}
\vspace{0.5\baselineskip}
\noindent\begin{tabularx}{\linewidth}{*{12}{>{\centering\arraybackslash}X}}
 \textcolor{gregoriocolor}{A} & \textcolor{gregoriocolor}{B} & \textcolor{gregoriocolor}{C} & \textcolor{gregoriocolor}{D} & \textcolor{gregoriocolor}{E} & F & \textcolor{gregoriocolor}{F} & \textcolor{gregoriocolor}{G} & \textcolor{gregoriocolor}{H} & \textcolor{gregoriocolor}{M} & \textcolor{gregoriocolor}{N} & \textcolor{gregoriocolor}{P} \\
 3 & 4 & 5 & 6 & 7 & 8 & 7 & 8 & 9 & 10 & 11 & 12 \\
\end{tabularx}

\begin{paracol}{2}
\selectlanguage{latin}
\lettrine[lines=2]{M}{essánæ,} in Sicília, 
 natális sanctórum Mártyrum Plácidi Mónachi, e beáti Benedícti Abbátis 
 discípulis, et ejus fratrum Eutychii et Victoríni, ac soróris eórum Fláviæ 
 Vírginis, itémque Donáti, Firmáti Diáconi, Fausti et aliórum trigínta 
 Monachórum, qui omnes a Manúcha piráta, pro Christi fide, necáti sunt.
\switchcolumn
\selectlanguage{english}
\lettrine[lines=2]{A}{t} Messina in Sicily, the birthday 
 of the holy martyrs Placidus, a monk who was a disciple of the blessed Abbot 
 Benedict, and of his brothers Eutychius and Victorinus, and the virgin 
 Flavia, their sister; also of Donatus, Firmatus, a deacon, Faustus, and 
 thirty other monks, who were murdered for the faith of Christ by the pirate 
 Manuchas.
\switchcolumn*
\selectlanguage{latin}
Apud Smyrnam item 
 natális beáti Thraséæ, Epíscopi Euméniæ, martyrio consummáti.
\switchcolumn
\selectlanguage{english}
At Smyrna, the birthday of blessed 
 Thraseas, bishop of Eumenia, who ended his career through martyrdom.
\switchcolumn*
\selectlanguage{latin}
Antisiodóri deposítio 
 sanctórum germanórum Firmáti Diáconi, et Flaviánæ Vírginis.
\switchcolumn
\selectlanguage{english}
At Auxerre, the death of the saintly 
 deacon Firmatus and the virgin Flaviana, his sister.
\switchcolumn*
\selectlanguage{latin}
Tréviris sanctórum 
 Mártyrum Palmátii et Sociórum; qui, in persecutióne Diocletiáni, sub 
 Rictiováro Præside, martyrium subiérunt.
\switchcolumn
\selectlanguage{english}
At Treves, the holy martyrs 
 Palmatius and his companions, who suffered martyrdom in the persecution of 
 Diocletian, under the governor Rictiovarus.
\switchcolumn*
\selectlanguage{latin}
Eódem die pássio sanctæ 
 Charitínæ Vírginis, quæ, sub Diocletiáno Imperatóre et Domítio Consulári, 
 ígnibus est cruciáta et in mare projécta, et, cum inde incólumis evasísset, 
 tandem, mánibus et pédibus abscíssis dentibúsque convúlsis, in oratióne 
 spíritum emísit.
\switchcolumn
\selectlanguage{english}
Also, under Emperor Diocletian and 
 the proconsul Domitius, St. Charitina, virgin. She was exposed to the 
 fire and thrown into the sea, but escaping uninjured, her hands and feet 
 were cut off and her teeth torn out, and finally she yielded up her spirit 
 in prayer.
\switchcolumn*
\selectlanguage{latin}
Ravénnæ sancti 
 Marcellíni, Epíscopi et Confessóris.
\switchcolumn
\selectlanguage{english}
At Ravenna, St. Marcellinus, bishop 
 and confessor.
\switchcolumn*
\selectlanguage{latin}
Valéntiæ, in Gállia, 
 sancti Apollináris Epíscopi, cujus vita virtútibus fuit illústris, et mors 
 signis ac prodígiis decoráta.
\switchcolumn
\selectlanguage{english}
At Valence in France, St. 
 Apollinaris, a bishop, renowned in life for virtues and in death for 
 miracles and prodigies.
\switchcolumn*
\selectlanguage{latin}
Eódem die sancti 
 Attiláni, Epíscopi Zamorénsis, quem beátus Urbánus Papa Secúndus in 
 Sanctórum númerum rétulit.
\switchcolumn
\selectlanguage{english}
Also, St. Attilanus, bishop of 
 Zamora, who was ranked among the saints by Pope Urban II.
\switchcolumn*
\selectlanguage{latin}
Romæ sanctæ Gallæ Víduæ, 
 fíliæ Symmachi Cónsulis, quæ, viro suo defúncto, apud Ecclésiam beáti Petri 
 multis annis oratióni, eleemósynis, jejúniis aliísque sanctis opéribus 
 inténta permánsit; cujus felicíssimum tránsitum sanctus Gregórius Papa 
 descrípsit.
\switchcolumn
\selectlanguage{english}
At Rome, St. Galla, widow, daughter 
 of the consul Symmachus. After the death of her husband, she remained 
 for many years near the church of St. Peter, devoted to prayer, almsgiving, 
 fasting, and other pious works. Her most happy death has been 
 described by Pope St. Gregory.
\switchcolumn*
\selectlanguage{latin}
\end{paracol}


% ---- martyrology/mart10/mart1006.htm
\needspace{10\baselineskip}
\begin{paracol}{2}
\selectlanguage{latin}
\begin{center}{\color{gregoriocolor} Prídie Nonas Octóbris. 
 Luna\dots\ }\end{center}
\switchcolumn
\selectlanguage{english}
\begin{center}{\color{gregoriocolor} The   6\textsuperscript{th} Day of 
 October. The\dots\ Day of the Moon.}\end{center}
\end{paracol}

\noindent\begin{tabularx}{\linewidth}{*{19}{>{\centering\arraybackslash}X}}
 \textcolor{gregoriocolor}{a} & \textcolor{gregoriocolor}{b} & \textcolor{gregoriocolor}{c} & \textcolor{gregoriocolor}{d} & \textcolor{gregoriocolor}{e} & \textcolor{gregoriocolor}{f} & \textcolor{gregoriocolor}{g} & \textcolor{gregoriocolor}{h} & \textcolor{gregoriocolor}{i} & \textcolor{gregoriocolor}{k} & \textcolor{gregoriocolor}{l} & \textcolor{gregoriocolor}{m} & \textcolor{gregoriocolor}{n} & \textcolor{gregoriocolor}{p} & \textcolor{gregoriocolor}{q} & \textcolor{gregoriocolor}{r} & \textcolor{gregoriocolor}{s} & \textcolor{gregoriocolor}{t} & \textcolor{gregoriocolor}{u} \\
 14 & 15 & 16 & 17 & 18 & 19 & 20 & 21 & 22 & 23 & 24 & 25 & 26 & 27 & 28 & 29 & 1 & 2 & 3 \\
\end{tabularx}
\vspace{0.5\baselineskip}
\noindent\begin{tabularx}{\linewidth}{*{12}{>{\centering\arraybackslash}X}}
 \textcolor{gregoriocolor}{A} & \textcolor{gregoriocolor}{B} & \textcolor{gregoriocolor}{C} & \textcolor{gregoriocolor}{D} & \textcolor{gregoriocolor}{E} & F & \textcolor{gregoriocolor}{F} & \textcolor{gregoriocolor}{G} & \textcolor{gregoriocolor}{H} & \textcolor{gregoriocolor}{M} & \textcolor{gregoriocolor}{N} & \textcolor{gregoriocolor}{P} \\
 4 & 5 & 6 & 7 & 8 & 9 & 8 & 9 & 10 & 11 & 12 & 13 \\
\end{tabularx}

\begin{paracol}{2}
\selectlanguage{latin}
\lettrine[lines=2]{I}{n} monastério Turris, 
 diœcésis Squillacénsis, in Calábria, sancti Brunónis Confessóris, qui 
 Ordinis Carthusianórum fuit Institútor.
\switchcolumn
\selectlanguage{english}
\lettrine[lines=2]{I}{n} the Monastery De Torre, in the 
 diocese of Squillace in Calabria, St. Bruno, confessor, founder of the Order 
 of the Carthusians.
\switchcolumn*
\selectlanguage{latin}
Laodicéæ, in Phrygia, 
 beáti Ságaris, Epíscopi et Mártyris; qui éxstitit unus de antíquis Pauli 
 Apóstoli discípulis.
\switchcolumn
\selectlanguage{english}
At Laodicea, the blessed bishop and 
 martyr Sagar, one of the first disciples of the apostle Paul.
\switchcolumn*
\selectlanguage{latin}
Antisiodóri sancti 
 Románi Epíscopi et Mártyris.
\switchcolumn
\selectlanguage{english}
At Auxerre, St. Romanus, bishop and 
 martyr.
\switchcolumn*
\selectlanguage{latin}
Cápuæ natális sanctórum 
 Mártyrum Marcélli, Cásti, Æmílii et Saturníni.
\switchcolumn
\selectlanguage{english}
At Capua, the birthday of the holy 
 martyrs Marcellus, Castus, Aemílius, and Saturninus.
\switchcolumn*
\selectlanguage{latin}
Tréviris commemorátio 
 innumerabílium fere Mártyrum, qui, in persecutióne Diocletiáni, sub 
 Rictiováro Præside, ob Christi fidem, vário mortis génere necáti sunt.
\switchcolumn
\selectlanguage{english}
At Treves, the commemoration of 
 innumerable martyrs, who were put death for the faith in various manners, 
 under the governor Rictiovarus, in the persecution of Diocletian.
\switchcolumn*
\selectlanguage{latin}
Agénni, in Gállia, 
 natális sanctæ Fídei, Vírginis et Mártyris; cujus exémplo beátus Caprásius, 
 ad martyrium animátus, agónem suum tertiodécimo Kaléndas Novémbris felíciter 
 consummávit.
\switchcolumn
\selectlanguage{english}
At Agen in France, the birthday of 
 St. Faith, virgin and martyr, by whose example blessed Caprasius was aroused 
 to martyrdom, and by martyrdom happily fulfilled his own trial.
\switchcolumn*
\selectlanguage{latin}
Item sanctæ Erótidis 
 Mártyris, quæ, Christi amóre succénsa, ignis superávit incéndium.
\switchcolumn
\selectlanguage{english}
Also, St. Erotis martyr, who, aflame 
 with love for Christ, triumphed over the flames of fire.
\switchcolumn*
\selectlanguage{latin}
Opitérgii, in Venetórum 
 fínibus, sancti Magni Epíscopi, cujus corpus Venétiis requiéscit.
\switchcolumn
\selectlanguage{english}
At Oderzo, in the neighbourhood of 
 Venice, St. Magnus, bishop, whose body rests at Venice.
\switchcolumn*
\selectlanguage{latin}
Neápoli, in Campánia, 
 deposítio sanctæ Maríæ-Francíscæ a Quinque Vulnéribus Dómini nostri Jesu 
 Christi, Vírginis, ex tértio Ordine sancti Francísci, quæ, virtútibus et 
 miráculis clara, a Pio Papa Nono sanctis Virgínibus adscrípta fuit.
\switchcolumn
\selectlanguage{english}
At Naples in Campania, the death of 
 St. Mary Frances of the Five Wounds of Our Lord Jesus Christ, a nun of the 
 Third Order of St. Francis. Because of her reputation for virtues and 
 the working of miracles, she was placed among the holy virgins by Pope Pius 
 IX.
\switchcolumn*
\selectlanguage{latin}
\end{paracol}


% ---- martyrology/mart10/mart1007.htm
\needspace{10\baselineskip}
\begin{paracol}{2}
\selectlanguage{latin}
\begin{center}{\color{gregoriocolor} Nonis Octóbris. 
 Luna\dots\ }\end{center}
\switchcolumn
\selectlanguage{english}
\begin{center}{\color{gregoriocolor} The   7\textsuperscript{th} Day of 
 October. The\dots\ Day of the Moon.}\end{center}
\end{paracol}

\noindent\begin{tabularx}{\linewidth}{*{19}{>{\centering\arraybackslash}X}}
 \textcolor{gregoriocolor}{a} & \textcolor{gregoriocolor}{b} & \textcolor{gregoriocolor}{c} & \textcolor{gregoriocolor}{d} & \textcolor{gregoriocolor}{e} & \textcolor{gregoriocolor}{f} & \textcolor{gregoriocolor}{g} & \textcolor{gregoriocolor}{h} & \textcolor{gregoriocolor}{i} & \textcolor{gregoriocolor}{k} & \textcolor{gregoriocolor}{l} & \textcolor{gregoriocolor}{m} & \textcolor{gregoriocolor}{n} & \textcolor{gregoriocolor}{p} & \textcolor{gregoriocolor}{q} & \textcolor{gregoriocolor}{r} & \textcolor{gregoriocolor}{s} & \textcolor{gregoriocolor}{t} & \textcolor{gregoriocolor}{u} \\
 15 & 16 & 17 & 18 & 19 & 20 & 21 & 22 & 23 & 24 & 25 & 26 & 27 & 28 & 29 & 1 & 2 & 3 & 4 \\
\end{tabularx}
\vspace{0.5\baselineskip}
\noindent\begin{tabularx}{\linewidth}{*{12}{>{\centering\arraybackslash}X}}
 \textcolor{gregoriocolor}{A} & \textcolor{gregoriocolor}{B} & \textcolor{gregoriocolor}{C} & \textcolor{gregoriocolor}{D} & \textcolor{gregoriocolor}{E} & F & \textcolor{gregoriocolor}{F} & \textcolor{gregoriocolor}{G} & \textcolor{gregoriocolor}{H} & \textcolor{gregoriocolor}{M} & \textcolor{gregoriocolor}{N} & \textcolor{gregoriocolor}{P} \\
 5 & 6 & 7 & 8 & 9 & 10 & 9 & 10 & 11 & 12 & 13 & 14 \\
\end{tabularx}

\begin{paracol}{2}
\selectlanguage{latin}
\lettrine[lines=2]{F}{estum} sacratíssimi 
 Rosárii beátæ Maríæ Vírginis; itémque sanctæ Maríæ de Victória commemorátio, 
 quam sanctus Pius Quintus, Póntifex Máximus, ob insígnem victóriam a 
 Christiánis bello naváli, ejúsdem sanctíssimæ Dei Genitrícis auxílio, hac 
 ipsa die de Turcis reportátam, quotánnis fíeri instítuit.
\switchcolumn
\selectlanguage{english}
\lettrine[lines=2]{T}{he} Feast of the Most Holy Rosary of 
 the blessed Virgin Mary, and the commemoration of St. Mary of Victory, which 
 Pope Pius V instituted to be kept yearly in memory of the great victory 
 granted on this day in a naval battle to the Christians over the Turks, by 
 the help of the Mother of God.
\switchcolumn*
\selectlanguage{latin}
Romæ, via Ardeatína, 
 sancti Marci, Papæ et Confessóris.
\switchcolumn
\selectlanguage{english}
At Rome, on the Ardeatine Way, the 
 death of St. Mark, pope and confessor.
\switchcolumn*
\selectlanguage{latin}
In Província quæ 
 nuncupátur Augústa Euphratésia, sanctórum Mártyrum Sérgii et Bacchi, 
 nobílium Romanórum, sub Maximiáno Imperatóre. Ex his Bacchus támdiu 
 nervis crudis cæsus est, quoadúsque, toto córpore discíssus, in Christi 
 confessióne emítteret spíritum; Sérgius vero clavátis cothúrnis pedes 
 indútus, et, cum in fide fixus manéret, data senténtia, jussus est decollári. 
 Beáti autem Sérgii nómine locus ubi quiéscit, Sergiópolis appellátus est, 
 et, ob præclára mirácula, frequénti Christianórum concúrsu honorátur.
\switchcolumn
\selectlanguage{english}
In the province of the Euphrates, 
 the holy martyrs Sergius and Bacchus, noble Romans, in the time of Emperor 
 Maximian. Bacchus was scourged with rough sinews until his body was 
 completely mangled, and breathed his last in the confession of Christ. 
 Sergius had his feet forced into shoes full of sharp-pointed nails, but, 
 remaining unshaken in the faith, he was sentenced to be beheaded. The 
 place where he rests is called after him Sergiopolis, and, on account of the 
 frequent miracles wrought there, is honoured by large gatherings of 
 Christians.
\switchcolumn*
\selectlanguage{latin}
Romæ sanctórum Mártyrum 
 Marcélli et Apuléji, qui prius quidem Simóni mago adhæsérunt; sed, vidéntes 
 mirabília quæ per Apóstolum Petrum Dóminus operabátur, ambo, relícto Simóne, 
 se doctrínæ Apostólicæ tradidérunt, ac, post passiónem Apostolórum, sub 
 Aureliáno Consulári, corónam martyrii reportárunt, sepultíque sunt non longe 
 ab Urbe.
\switchcolumn
\selectlanguage{english}
At Rome, the holy martyrs Marcellus 
 and Apulcius, who at first were followers of Simon Magus, but seeing the 
 wonders which the Lord performed through the apostle Peter, they abandoned 
 Simon and embraced the apostolic doctrine. After the death of the 
 apostles, under the proconsul Aurelian, they won the crown of martyrdom and 
 were buried near the city.
\switchcolumn*
\selectlanguage{latin}
Item apud Augústam 
 Euphratésiam sanctæ Júliæ Vírginis, quæ, sub Marciáno Præside, martyrium 
 consummávit.
\switchcolumn
\selectlanguage{english}
Also in the province of the 
 Euphrates, St. Julia, virgin, who suffered martyrdom under the governor 
 Marcian.
\switchcolumn*
\selectlanguage{latin}
Patávii sanctæ Justínæ, 
 Vírginis et Mártyris; quæ, a beáto Prosdócimo, sancti Petri discípulo, 
 baptizáta, et, cum in fide Christi constánter persísteret, Máximi Præsidis 
 jussu, transverberáta gládio, migrávit ad Dóminum.
\switchcolumn
\selectlanguage{english}
At Padua, St. Justina, virgin and 
 martyr, who was baptized by blessed Prosdocimus, a disciple of St. Peter. 
 Because she remained firm in the faith of Christ, she was put to the sword 
 by order of the governor Maximus, and thus went to God.
\switchcolumn*
\selectlanguage{latin}
Apud Bitúricas, in 
 Aquitánia, sancti Augústi, Presbyteri et Confessóris.
\switchcolumn
\selectlanguage{english}
At Bourges, St. Augustus, priest and 
 confessor.
\switchcolumn*
\selectlanguage{latin}
In pago Rheménsi sancti 
 Heláni Presbyteri.
\switchcolumn
\selectlanguage{english}
In the diocese of Rheims, St. 
 Helanus, priest.
\switchcolumn*
\selectlanguage{latin}
In Suécia Translátio 
 córporis sanctæ Birgíttæ Víduæ.
\switchcolumn
\selectlanguage{english}
In Sweden, the translation of the 
 body of St. Bridget, widow.
\switchcolumn*
\selectlanguage{latin}
\end{paracol}


% ---- martyrology/mart10/mart1008.htm
\needspace{10\baselineskip}
\begin{paracol}{2}
\selectlanguage{latin}
\begin{center}{\color{gregoriocolor} Octávo Idus Octóbris. 
 Luna\dots\ }\end{center}
\switchcolumn
\selectlanguage{english}
\begin{center}{\color{gregoriocolor} The   8\textsuperscript{th} Day of 
 October. The\dots\ Day of the Moon.}\end{center}
\end{paracol}

\noindent\begin{tabularx}{\linewidth}{*{19}{>{\centering\arraybackslash}X}}
 \textcolor{gregoriocolor}{a} & \textcolor{gregoriocolor}{b} & \textcolor{gregoriocolor}{c} & \textcolor{gregoriocolor}{d} & \textcolor{gregoriocolor}{e} & \textcolor{gregoriocolor}{f} & \textcolor{gregoriocolor}{g} & \textcolor{gregoriocolor}{h} & \textcolor{gregoriocolor}{i} & \textcolor{gregoriocolor}{k} & \textcolor{gregoriocolor}{l} & \textcolor{gregoriocolor}{m} & \textcolor{gregoriocolor}{n} & \textcolor{gregoriocolor}{p} & \textcolor{gregoriocolor}{q} & \textcolor{gregoriocolor}{r} & \textcolor{gregoriocolor}{s} & \textcolor{gregoriocolor}{t} & \textcolor{gregoriocolor}{u} \\
 16 & 17 & 18 & 19 & 20 & 21 & 22 & 23 & 24 & 25 & 26 & 27 & 28 & 29 & 1 & 2 & 3 & 4 & 5 \\
\end{tabularx}
\vspace{0.5\baselineskip}
\noindent\begin{tabularx}{\linewidth}{*{12}{>{\centering\arraybackslash}X}}
 \textcolor{gregoriocolor}{A} & \textcolor{gregoriocolor}{B} & \textcolor{gregoriocolor}{C} & \textcolor{gregoriocolor}{D} & \textcolor{gregoriocolor}{E} & F & \textcolor{gregoriocolor}{F} & \textcolor{gregoriocolor}{G} & \textcolor{gregoriocolor}{H} & \textcolor{gregoriocolor}{M} & \textcolor{gregoriocolor}{N} & \textcolor{gregoriocolor}{P} \\
 6 & 7 & 8 & 9 & 10 & 11 & 10 & 11 & 12 & 13 & 14 & 15 \\
\end{tabularx}

\begin{paracol}{2}
\selectlanguage{latin}
\lettrine[lines=2]{S}{anctæ} Birgíttæ Víduæ, 
 cujus dies natális décimo Kaléndas Augústi, ac Translátio Nonis Octóbris 
 recensétur.
\switchcolumn
\selectlanguage{english}
\lettrine[lines=2]{S}{t.} Bridget, widow, whose birthday 
 is observed on the 23rd of July, and the translation of her holy body on the 
 7th of October.
\switchcolumn*
\selectlanguage{latin}
Eódem die natális beáti 
 Simeónis senis, qui in Evangélio Dóminum Jesum, præsentátum in Templo, suis 
 in ulnis accepísse ac de illo prophetásse légitur.
\switchcolumn
\selectlanguage{english}
Also, the birthday of blessed 
 Simeon, an aged man, who as we read in the Gospel, took our Lord Jesus in 
 his arms and prophesied concerning him when he was presented in the Temple.
\switchcolumn*
\selectlanguage{latin}
Laodicéæ, in Phrygia, 
 sancti Artémonis Presbyteri, qui per ignem, sub Diocletiáno, martyrii 
 corónam accépit.
\switchcolumn
\selectlanguage{english}
At Laodicea in Phrygia, during the 
 reign of Diocletian, St. Artemon, a priest, who gained the crown of 
 martyrdom by fire.
\switchcolumn*
\selectlanguage{latin}
Thessalonícæ sancti 
 Demétrii Procónsulis, qui, cum plúrimos ad Christi fidem perdúceret, ideo, 
 Maximiáni Imperatóris jussu lánceis confóssus, martyrium consummávit.
\switchcolumn
\selectlanguage{english}
At Thessalonica, St. Demetrius, a 
 proconsul. For having brought many to the faith of Christ he was 
 pierced with spears by order of Emperor Maximian, and thus completed his 
 martyrdom.
\switchcolumn*
\selectlanguage{latin}
Ibídem sancti Néstoris 
 Mártyris.
\switchcolumn
\selectlanguage{english}
In the same place, St. Nestor, 
 martyr.
\switchcolumn*
\selectlanguage{latin}
Híspali, in Hispánia, 
 sancti Petri Mártyris.
\switchcolumn
\selectlanguage{english}
At Seville in Spain, St. Peter, 
 martyr.
\switchcolumn*
\selectlanguage{latin}
Cæsaréæ, in Palæstína, 
 pássio sanctæ Reparátæ, Vírginis et Mártyris; quæ, cum nollet idólis 
 sacrificáre, sub Décio Imperatóre, váriis tormentórum genéribus cruciátur, 
 ac demum gládio percútitur. Ipsíus autem ánima, in colúmbæ spécie, de 
 córpore égredi cælúmque conscéndere visa est.
\switchcolumn
\selectlanguage{english}
At Caesarea in Palestine, in the 
 reign of Decius, St. Reparata, virgin and martyr. For refusing to 
 sacrifice to idols, she was subjected to various kinds of torments and was 
 finally struck with the sword. Her soul was seen to leave her body in 
 the form of a dove and ascend to heaven.
\switchcolumn*
\selectlanguage{latin}
In território 
 Laudunénsi natális sanctæ Benedíctæ, Vírginis et Mártyris.
\switchcolumn
\selectlanguage{english}
In the country of Laon, St. 
 Benedicta, virgin and martyr.
\switchcolumn*
\selectlanguage{latin}
Ancónæ sanctárum 
 Palatiátis et Lauréntiæ, quæ, in persecutióne Diocletiáni, sub Dióne Præside, 
 in exsílium deportátæ, labóribus et ærúmnis conféctæ sunt.
\switchcolumn
\selectlanguage{english}
At Ancona, Saints Palatius and 
 Laurentia, who were sent into exile during the persecution of Diocletian, 
 under the governor Dion, and were overcome by the weight of toil and misery.
\switchcolumn*
\selectlanguage{latin}
Rotómagi sancti Evódii, 
 Epíscopi et Confessóris.
\switchcolumn
\selectlanguage{english}
At Rouen, St. Evodius, bishop and 
 confessor.
\switchcolumn*
\selectlanguage{latin}
Hierosólymis sanctæ 
 Pelágiæ, cognoménto Pæniténtis.
\switchcolumn
\selectlanguage{english}
At Jerusalem, St. Palagia, surnamed 
 the Penitent.
\switchcolumn*
\selectlanguage{latin}
\end{paracol}


% ---- martyrology/mart10/mart1009.htm
\needspace{10\baselineskip}
\begin{paracol}{2}
\selectlanguage{latin}
\begin{center}{\color{gregoriocolor} Séptimo Idus Octóbris. 
 Luna\dots\ }\end{center}
\switchcolumn
\selectlanguage{english}
\begin{center}{\color{gregoriocolor} The   9\textsuperscript{th} Day of 
 October. The\dots\ Day of the Moon.}\end{center}
\end{paracol}

\noindent\begin{tabularx}{\linewidth}{*{19}{>{\centering\arraybackslash}X}}
 \textcolor{gregoriocolor}{a} & \textcolor{gregoriocolor}{b} & \textcolor{gregoriocolor}{c} & \textcolor{gregoriocolor}{d} & \textcolor{gregoriocolor}{e} & \textcolor{gregoriocolor}{f} & \textcolor{gregoriocolor}{g} & \textcolor{gregoriocolor}{h} & \textcolor{gregoriocolor}{i} & \textcolor{gregoriocolor}{k} & \textcolor{gregoriocolor}{l} & \textcolor{gregoriocolor}{m} & \textcolor{gregoriocolor}{n} & \textcolor{gregoriocolor}{p} & \textcolor{gregoriocolor}{q} & \textcolor{gregoriocolor}{r} & \textcolor{gregoriocolor}{s} & \textcolor{gregoriocolor}{t} & \textcolor{gregoriocolor}{u} \\
 17 & 18 & 19 & 20 & 21 & 22 & 23 & 24 & 25 & 26 & 27 & 28 & 29 & 1 & 2 & 3 & 4 & 5 & 6 \\
\end{tabularx}
\vspace{0.5\baselineskip}
\noindent\begin{tabularx}{\linewidth}{*{12}{>{\centering\arraybackslash}X}}
 \textcolor{gregoriocolor}{A} & \textcolor{gregoriocolor}{B} & \textcolor{gregoriocolor}{C} & \textcolor{gregoriocolor}{D} & \textcolor{gregoriocolor}{E} & F & \textcolor{gregoriocolor}{F} & \textcolor{gregoriocolor}{G} & \textcolor{gregoriocolor}{H} & \textcolor{gregoriocolor}{M} & \textcolor{gregoriocolor}{N} & \textcolor{gregoriocolor}{P} \\
 7 & 8 & 9 & 10 & 11 & 12 & 11 & 12 & 13 & 14 & 15 & 16 \\
\end{tabularx}

\begin{paracol}{2}
\selectlanguage{latin}
\lettrine[lines=2]{R}{omæ} sancti Joánnis 
 Leonárdi, Confessóris, Fundatóris Congregatiónis Clericórum Regulárium a 
 Matre Dei, labóribus et miráculis clari, cujus ópera Missiónes a Propagánda 
 Fide institútæ sunt.
\switchcolumn
\selectlanguage{english}
\lettrine[lines=2]{A}{t} Rome, St. John Leonard, 
 confessor, founder of the Congregation of Clerks Regular of the Mother of 
 God, renowned for his labours and miracles, and by whose zeal were begun 
 missions for the propagation of the faith.
\switchcolumn*
\selectlanguage{latin}
Lutétiæ Parisiórum 
 natális sanctórum Mártyrum Dionysii Areopagítæ Epíscopi, Rústici Presbyteri, 
 et Eleuthérii Diáconi. Ex his Dionysius, ab Apóstolo Paulo baptizátus, 
 primus Atheniénsium Epíscopus ordinátus est; deínde Romam venit, atque inde 
 a beáto Cleménte, Románo Pontífice, in Gállias prædicándi grátia diréctus 
 est, et ad præfátam urbem devénit; ibíque, cum per áliquot annos commíssum 
 sibi opus fidéliter prosecútus esset, tandem, a Præfécto Fescénnio, post 
 gravíssima tormentórum génera, una cum Sóciis, gládio animadvérsus, 
 martyrium complévit.
\switchcolumn
\selectlanguage{english}
At Paris, the birthday of the holy 
 martyrs Denis the Areopagite, a bishop, Rusticus, a priest, and Eleutherius, 
 a deacon. Denis was baptized by the apostle St. Paul, and consecrated 
 first bishop of Athens. Then going to Rome, he was sent to France by 
 the blessed Roman Pontiff Clement to preach the Gospel. He proceeded 
 to Paris, and after having for some years faithfully filled the office 
 entrusted to him, he was subjected to the severest kinds of torments by the 
 prefect Fescennius, and at length was beheaded with his companions, thus 
 completing his martyrdom.
\switchcolumn*
\selectlanguage{latin}
Eódem die memória 
 sancti Abrahæ, Patriárchæ et ómnium credéntium Patris.
\switchcolumn
\selectlanguage{english}
On the same day, the commemoration 
 of the holy patriarch Abraham, father of all believers.
\switchcolumn*
\selectlanguage{latin}
Apud Cassínum sancti 
 Deúsdedit Abbátis, qui, a Sicárdo tyránno in cárcerem trusus, illic, 
 fame et ærúmnis conféctus, réddidit spíritum.
\switchcolumn
\selectlanguage{english}
At Monte Cassino, St. Deusdedit, 
 abbot, who was cast into prison by the tyrant Sicardus, and being there 
 consumed with hunger and misery, yielded up his soul.
\switchcolumn*
\selectlanguage{latin}
Apud Júliam, in 
 território Parménsi, via Cláudia, sancti Domníni Mártyris; qui, sub 
 Maximiáno Imperatóre, cum vellet persecutiónis rábiem declináre, a 
 persequéntibus est comprehénsus, et, gládio transverberátus, glorióse 
 occúbuit.
\switchcolumn
\selectlanguage{english}
At Julia, in the region of Parma, on 
 the Via Claudia, St. Domninus, martyr. Under the Emperor Maximian, in 
 the rage of persecution, he was taken by the persecutors and died gloriously 
 by being pierced with a sword.
\switchcolumn*
\selectlanguage{latin}
In Hannónia sancti 
 Gisléni, Epíscopi et Confessóris; qui, relícto Episcopátu, Mónachi vitam in 
 monastério a se constrúcto exércuit, et multis virtútibus cláruit.
\switchcolumn
\selectlanguage{english}
In Hainault, St. Gislenus, bishop 
 and confessor, who resigning his bishopric, led the monastic life in a 
 monastery built by himself, and was distinguished by many virtues.
\switchcolumn*
\selectlanguage{latin}
Valéntiæ, in Hispánia 
 Tarraconénsi, sancti Ludovíci Bertrándi, ex Ordine Prædicatórum, Confessóris; 
 qui, apostólico spíritu clarus, Evangélium quod Americánis prædicáverat, 
 vitæ innocéntia multísque éditis miráculis confirmávit.
\switchcolumn
\selectlanguage{english}
At Valencia in Spain, St. Louis 
 Bertrand, of the Order of Preachers. Being filled with the apostolic 
 spirit, he confirmed by the innocency of his life and the working of many 
 miracles the Gospel which he had preached in America.
\switchcolumn*
\selectlanguage{latin}
Hierosólymis sanctórum 
 Androníci et Athanásiæ cónjugis.
\switchcolumn
\selectlanguage{english}
At Jerusalem, Saints Andronicus and 
 his wife Athanasia.
\switchcolumn*
\selectlanguage{latin}
Antiochíæ sanctæ Públiæ 
 Abbatíssæ, quæ, transeúnte Juliáno Apóstata, Davídicum illud cum suis 
 Virgínibus canens: « Simulácra Géntium argéntum et aurum », et « Símiles 
 illis fiant qui fáciunt ea », Imperatóris jussu, álapis cæsa est, et 
 gráviter objurgáta.
\switchcolumn
\selectlanguage{english}
At Antioch, St. Publia, abbess. 
 While Julian the Apostate was passing by, she and her religious sang these 
 words of David: ``The idols of the Gentiles are silver and gold,'' and ``Let 
 them that make them be like unto them.'' By the command of the emperor, 
 she was struck on the face and severely rebuked.
\switchcolumn*
\selectlanguage{latin}
\end{paracol}


% ---- martyrology/mart10/mart1010.htm
\needspace{10\baselineskip}
\begin{paracol}{2}
\selectlanguage{latin}
\begin{center}{\color{gregoriocolor} Sexto Idus Octóbris. 
 Luna\dots\ }\end{center}
\switchcolumn
\selectlanguage{english}
\begin{center}{\color{gregoriocolor} The   10\textsuperscript{th} Day of 
 October. The\dots\ Day of the Moon.}\end{center}
\end{paracol}

\noindent\begin{tabularx}{\linewidth}{*{19}{>{\centering\arraybackslash}X}}
 \textcolor{gregoriocolor}{a} & \textcolor{gregoriocolor}{b} & \textcolor{gregoriocolor}{c} & \textcolor{gregoriocolor}{d} & \textcolor{gregoriocolor}{e} & \textcolor{gregoriocolor}{f} & \textcolor{gregoriocolor}{g} & \textcolor{gregoriocolor}{h} & \textcolor{gregoriocolor}{i} & \textcolor{gregoriocolor}{k} & \textcolor{gregoriocolor}{l} & \textcolor{gregoriocolor}{m} & \textcolor{gregoriocolor}{n} & \textcolor{gregoriocolor}{p} & \textcolor{gregoriocolor}{q} & \textcolor{gregoriocolor}{r} & \textcolor{gregoriocolor}{s} & \textcolor{gregoriocolor}{t} & \textcolor{gregoriocolor}{u} \\
 18 & 19 & 20 & 21 & 22 & 23 & 24 & 25 & 26 & 27 & 28 & 29 & 1 & 2 & 3 & 4 & 5 & 6 & 7 \\
\end{tabularx}
\vspace{0.5\baselineskip}
\noindent\begin{tabularx}{\linewidth}{*{12}{>{\centering\arraybackslash}X}}
 \textcolor{gregoriocolor}{A} & \textcolor{gregoriocolor}{B} & \textcolor{gregoriocolor}{C} & \textcolor{gregoriocolor}{D} & \textcolor{gregoriocolor}{E} & F & \textcolor{gregoriocolor}{F} & \textcolor{gregoriocolor}{G} & \textcolor{gregoriocolor}{H} & \textcolor{gregoriocolor}{M} & \textcolor{gregoriocolor}{N} & \textcolor{gregoriocolor}{P} \\
 8 & 9 & 10 & 11 & 12 & 13 & 12 & 13 & 14 & 15 & 16 & 17 \\
\end{tabularx}

\begin{paracol}{2}
\selectlanguage{latin}
\lettrine[lines=2]{S}{ancti} Francísci Bórgiæ, 
 Sacerdótis e Societáte Jesu et Confessóris, cujus dies natális prídie 
 Kaléndas Octóbris recensétur.
\switchcolumn
\selectlanguage{english}
\lettrine[lines=2]{S}{t.} Francis Borgia, confessor, 
 priest of the Society of Jesus, whose birthday is mentioned on the 30th of 
 September.
\switchcolumn*
\selectlanguage{latin}
Apud Septam, in 
 Mauritánia Tingitána, pássio sanctórum septem Mártyrum, ex Ordine Minórum, 
 scílicet Daniélis, Samuélis, Angeli, Leónis, Nicolái, Hugolíni, et Domni; 
 qui, cum essent omnes præter Domnum Sacerdótes, ibídem, ob Evangélii 
 prædicatiónem et Mahuméticæ confutatiónem sectæ, a Saracénis contumélias, 
 víncula et flagélla perpéssi, demum, capítibus abscíssis, martyrii palmam 
 adépti sunt.
\switchcolumn
\selectlanguage{english}
At Ceuta in Morocco, the passion of 
 seven holy martyrs of the Order of Friars Minor: Daniel, Samuel, Angelus, 
 Leo, Nicholas, Ugolino, and Domnus, all of whom were priests except Domnus. 
 Because they had preached the Gospel and put to silence the doctrines of 
 Mohammed, they suffered insults, fetters, and scourgings from the Saracens 
 in that place. They were at last beheaded and thus obtained the palm 
 of martyrdom.
\switchcolumn*
\selectlanguage{latin}
Colóniæ Agrippínæ 
 sancti Gereónis Mártyris, cum áliis trecéntis decem et octo; qui pro vera 
 pietáte, in persecutióne Maximiáni, gládiis patiéntes colla subdidérunt.
\switchcolumn
\selectlanguage{english}
At Cologne, in the persecution of 
 Maximian, St. Gereon and three hundred and eighteen other martyrs who 
 patiently bowed to the sword for the true religion.
\switchcolumn*
\selectlanguage{latin}
In território ejúsdem 
 urbis sanctórum Victóris et Sociórum Mártyrum.
\switchcolumn
\selectlanguage{english}
In the neighbourhood of the same 
 city, the holy martyrs Victor and his companions.
\switchcolumn*
\selectlanguage{latin}
Bonnæ, in Germánia, sanctórum Mártyrum Cássii et Floréntii, cum áliis plúrimis.
\switchcolumn
\selectlanguage{english}
At Bonn in Germany, the holy martyrs 
 Cassius and Florentius, with many others.
\switchcolumn*
\selectlanguage{latin}
Nicomedíæ sanctórum 
 Mártyrum Eulámpii, et soróris Eulámpiæ Vírginis. Hæc porro, cum 
 audísset pro Christo fratrem torquéri, in médiam turbam exsíluit, et, 
 fratrem amplexáta, huic sóciam se adjúnxit; atque ambo, conjécti in ollam 
 fervéntis ólei, sed nulla ex parte læsi, tandem, una cum áliis ducéntis, qui 
 eo miráculo permóti credidérunt in Christum, cápitis obtruncatióne martyrium 
 complevérunt.
\switchcolumn
\selectlanguage{english}
At Nicomedia, the holy martyrs 
 Eulampius, and his sister, the virgin Eulampia. Upon hearing that her 
 brother was tortured for Christ, she rushed through the crowd, embraced him, 
 and became his companion. Both were cast into a cauldron of boiling 
 oil, but being uninjured, their martyrdom was completed by beheading along 
 with two hundred others, who, impressed by the miracle, had believed in 
 Christ.
\switchcolumn*
\selectlanguage{latin}
In Creta ínsula beáti 
 Pinyti, inter Epíscopos nobilíssimi. Hic, Gnósiæ urbis Epíscopus, sub 
 Marco Antoníno Vero et Lúcio Aurélio Cómmodo flóruit, et in scriptis suis, 
 velut in quodam spéculo, vivéntem sui relíquit imáginem.
\switchcolumn
\selectlanguage{english}
On the island of Crete, blessed 
 Pinytus, most noble of bishops. He was bishop of Gnosia, and 
 flourished under Marcus Antoninus Verus and Lucius Aurelius Commodus. 
 He left in his writings, as in a mirror, a vivid picture of himself.
\switchcolumn*
\selectlanguage{latin}
Eboráci, in Anglia, 
 sancti Paulíni Epíscopi, qui fuit beáti Gregórii Papæ discípulus; et, una 
 cum áliis, ad prædicándum Evangélium illuc ab eo missus, Edwínum Regem 
 ejúsque pópulum ad Christi fidem convértit.
\switchcolumn
\selectlanguage{english}
At York in England, the holy bishop 
 Paulinus, disciple of the blessed pope Gregory. He was sent there by 
 that pope along with others to preach the Gospel, and he converted King 
 Edwin and his people to the faith of Christ.
\switchcolumn*
\selectlanguage{latin}
Populónii, in Túscia, 
 sancti Cerbónii, Epíscopi et Confessóris, qui (ut sanctus Gregórius Papa 
 refert) in vita et morte miráculis cláruit.
\switchcolumn
\selectlanguage{english}
At Piombino in Tuscany, St. 
 Cerbonius, bishop and confessor. St. Gregory relates that he was 
 renowned for miracles, both during life and after death.
\switchcolumn*
\selectlanguage{latin}
Verónæ sancti Cerbónii 
 Epíscopi.
\switchcolumn
\selectlanguage{english}
At Verona, another St. Cerbonius, 
 bishop.
\switchcolumn*
\selectlanguage{latin}
Cápuæ sancti Paulíni 
 Epíscopi.
\switchcolumn
\selectlanguage{english}
At Capua, St. Paulinus, bishop.
\switchcolumn*
\selectlanguage{latin}
\end{paracol}


% ---- martyrology/mart10/mart1011.htm
\needspace{10\baselineskip}
\begin{paracol}{2}
\selectlanguage{latin}
\begin{center}{\color{gregoriocolor} Quinto Idus Octóbris. 
 Luna\dots\ }\end{center}
\switchcolumn
\selectlanguage{english}
\begin{center}{\color{gregoriocolor} The   11\textsuperscript{th} Day of 
 October. The\dots\ Day of the Moon.}\end{center}
\end{paracol}

\noindent\begin{tabularx}{\linewidth}{*{19}{>{\centering\arraybackslash}X}}
 \textcolor{gregoriocolor}{a} & \textcolor{gregoriocolor}{b} & \textcolor{gregoriocolor}{c} & \textcolor{gregoriocolor}{d} & \textcolor{gregoriocolor}{e} & \textcolor{gregoriocolor}{f} & \textcolor{gregoriocolor}{g} & \textcolor{gregoriocolor}{h} & \textcolor{gregoriocolor}{i} & \textcolor{gregoriocolor}{k} & \textcolor{gregoriocolor}{l} & \textcolor{gregoriocolor}{m} & \textcolor{gregoriocolor}{n} & \textcolor{gregoriocolor}{p} & \textcolor{gregoriocolor}{q} & \textcolor{gregoriocolor}{r} & \textcolor{gregoriocolor}{s} & \textcolor{gregoriocolor}{t} & \textcolor{gregoriocolor}{u} \\
 19 & 20 & 21 & 22 & 23 & 24 & 25 & 26 & 27 & 28 & 29 & 1 & 2 & 3 & 4 & 5 & 6 & 7 & 8 \\
\end{tabularx}
\vspace{0.5\baselineskip}
\noindent\begin{tabularx}{\linewidth}{*{12}{>{\centering\arraybackslash}X}}
 \textcolor{gregoriocolor}{A} & \textcolor{gregoriocolor}{B} & \textcolor{gregoriocolor}{C} & \textcolor{gregoriocolor}{D} & \textcolor{gregoriocolor}{E} & F & \textcolor{gregoriocolor}{F} & \textcolor{gregoriocolor}{G} & \textcolor{gregoriocolor}{H} & \textcolor{gregoriocolor}{M} & \textcolor{gregoriocolor}{N} & \textcolor{gregoriocolor}{P} \\
 9 & 10 & 11 & 12 & 13 & 14 & 13 & 14 & 15 & 16 & 17 & 18 \\
\end{tabularx}

\begin{paracol}{2}
\selectlanguage{latin}
\lettrine[lines=1]{F}{estum} Maternitátis beátæ Maríæ Vírginis.
\switchcolumn
\selectlanguage{english}
\lettrine[lines=1]{T}{he} Motherhood of the Blessed Virgin Mary.
\switchcolumn*
\selectlanguage{latin}
Tarsi, in Cilícia, sanctárum mulíerum Zenáidis et Philoníllæ sorórum, quæ beáti Pauli Apóstoli 
 consanguíneæ et in fide fuérunt discípulæ.
\switchcolumn
\selectlanguage{english}
At Tarsus in Cilicia, the holy women 
 Zenaides and Philonilla, sisters, who were relatives of the blessed apostle 
 Paul and his disciples in the faith.
\switchcolumn*
\selectlanguage{latin}
In pago Vilcassíno, in 
 Gállia, pássio sanctórum Mártyrum Nicásii, qui erat Rotomagénsis Epíscopus, 
 Quiríni Presbyteri, Scubículi Diáconi, et Piéntiæ Vírginis, sub Præside 
 Fescénnio.
\switchcolumn
\selectlanguage{english}
In the neighbourhood of Vexin in 
 France, in the time of the governor Fescennius, the passion of the holy 
 martyrs Nicasius, bishop of Rouen, the priest Quirinus, the deacon 
 Scubiculus, and Pientia, a virgin.
\switchcolumn*
\selectlanguage{latin}
Vesontióne, in Gálliis, 
 sancti Germáni, Epíscopi et Mártyris.
\switchcolumn
\selectlanguage{english}
At Besançon in France, St. Germanus, 
 bishop and martyr.
\switchcolumn*
\selectlanguage{latin}
Item pássio sanctórum 
 Anastásii Presbyteri, Plácidi, Genésii et Sociórum.
\switchcolumn
\selectlanguage{english}
Also, the martyrdom of the Saints 
 Anastasius, a priest, Placidus, Genesius, and their companions.
\switchcolumn*
\selectlanguage{latin}
Tarsi, in Cilícia, natális sanctórum Mártyrum Tháraci, Probi et Andrónici; qui, in persecutióne 
 Diocletiáni, longo témpore cárceris squalóre afflícti, et tértio divérsis 
 torméntis et supplíciis examináti, tandem, in confessióne Christi, abscíssis 
 cervícibus, triúmphum glóriæ sunt adépti.
\switchcolumn
\selectlanguage{english}
At Tarsus in Cilicia, the birthday 
 of the holy martyrs Tharacus, Probus, and Andronicus, who endured a long and 
 painful imprisonment during the persecution of Diocletian. They were 
 three times subjected to diverse punishments and tortures, and finally 
 obtained a glorious triumph for the confession of Christ by having their 
 heads struck off.
\switchcolumn*
\selectlanguage{latin}
In Thebáide sancti 
 Sarmátæ, qui fuit discípulus beáti Antónii Abbátis, et a Saracénis pro 
 Christo necátus est.
\switchcolumn
\selectlanguage{english}
In Thebais, St. Sarmata, disciple of 
 the blessed abbot Anthony, who was put to death for Christ by the Saracens.
\switchcolumn*
\selectlanguage{latin}
Ucétiæ, in Gállia 
 Narbonénsi, sancti Firmíni, Epíscopi et Confessóris.
\switchcolumn
\selectlanguage{english}
At Uzea in France, St. Firmin, 
 bishop and confessor.
\switchcolumn*
\selectlanguage{latin}
Calótii, in diœcési 
 Asténsi, olim Papiénsi, sancti Alexándri Sauli, e Clericórum Regulárium 
 sancti Pauli Congregatióne, Epíscopi et Confessóris; quem, génere, 
 virtútibus, doctrína et miráculis clarum, Pius Décimus, Póntifex Máximus, 
 Sanctórum fastis adscrípsit.
\switchcolumn
\selectlanguage{english}
At Calozzo, in the diocese of Asti, 
 formerly that of Pavia, St. Alexander Sauli, bishop and confessor of the 
 Clerics Regular of St. Paul. He was of noble birth and renowned for 
 virtues, learning, and miracles. Pope Pius X placed him in the canon 
 of the saints.
\switchcolumn*
\selectlanguage{latin}
In monastério 
 Achadh-boénsi, in Hibérnia, sancti Cánici, Presbyteri et Abbátis.
\switchcolumn
\selectlanguage{english}
In the monastery of Aghaboe in 
 Ireland, St. Kenny, priest and abbot.
\switchcolumn*
\selectlanguage{latin}
Lyræ, in Bélgio, 
 deposítio sancti Gummári Confessóris,
\switchcolumn
\selectlanguage{english}
At Lier in Belgium, the death of St. 
 Gummarus, confessor.
\switchcolumn*
\selectlanguage{latin}
Apud Rhédones, in 
 Gállia, sancti Æmiliáni Confessóris.
\switchcolumn
\selectlanguage{english}
At Rennes in France, St. Emilian, 
 confessor.
\switchcolumn*
\selectlanguage{latin}
Verónæ sanctæ Placídiæ 
 Vírginis.
\switchcolumn
\selectlanguage{english}
At Verona, St. Placidia, virgin.
\switchcolumn*
\selectlanguage{latin}
\end{paracol}


% ---- martyrology/mart10/mart1012.htm
\needspace{10\baselineskip}
\begin{paracol}{2}
\selectlanguage{latin}
\begin{center}{\color{gregoriocolor} Quarto Idus Octóbris. 
 Luna\dots\ }\end{center}
\switchcolumn
\selectlanguage{english}
\begin{center}{\color{gregoriocolor} The   12\textsuperscript{th} Day of 
 October. The\dots\ Day of the Moon.}\end{center}
\end{paracol}

\noindent\begin{tabularx}{\linewidth}{*{19}{>{\centering\arraybackslash}X}}
 \textcolor{gregoriocolor}{a} & \textcolor{gregoriocolor}{b} & \textcolor{gregoriocolor}{c} & \textcolor{gregoriocolor}{d} & \textcolor{gregoriocolor}{e} & \textcolor{gregoriocolor}{f} & \textcolor{gregoriocolor}{g} & \textcolor{gregoriocolor}{h} & \textcolor{gregoriocolor}{i} & \textcolor{gregoriocolor}{k} & \textcolor{gregoriocolor}{l} & \textcolor{gregoriocolor}{m} & \textcolor{gregoriocolor}{n} & \textcolor{gregoriocolor}{p} & \textcolor{gregoriocolor}{q} & \textcolor{gregoriocolor}{r} & \textcolor{gregoriocolor}{s} & \textcolor{gregoriocolor}{t} & \textcolor{gregoriocolor}{u} \\
 20 & 21 & 22 & 23 & 24 & 25 & 26 & 27 & 28 & 29 & 1 & 2 & 3 & 4 & 5 & 6 & 7 & 8 & 9 \\
\end{tabularx}
\vspace{0.5\baselineskip}
\noindent\begin{tabularx}{\linewidth}{*{12}{>{\centering\arraybackslash}X}}
 \textcolor{gregoriocolor}{A} & \textcolor{gregoriocolor}{B} & \textcolor{gregoriocolor}{C} & \textcolor{gregoriocolor}{D} & \textcolor{gregoriocolor}{E} & F & \textcolor{gregoriocolor}{F} & \textcolor{gregoriocolor}{G} & \textcolor{gregoriocolor}{H} & \textcolor{gregoriocolor}{M} & \textcolor{gregoriocolor}{N} & \textcolor{gregoriocolor}{P} \\
 10 & 11 & 12 & 13 & 14 & 15 & 14 & 15 & 16 & 17 & 18 & 19 \\
\end{tabularx}

\begin{paracol}{2}
\selectlanguage{latin}
\lettrine[lines=2]{R}{omæ} sanctórum Mártyrum 
 Evágrii, Prisciáni, et Sociórum.
\switchcolumn
\selectlanguage{english}
\lettrine[lines=2]{A}{t} Rome, the holy martyrs Evagrius, 
 Priscian, and their companions.
\switchcolumn*
\selectlanguage{latin}
In Africa sanctórum 
 Confessórum et Mártyrum quátuor míllium nongentórum sexagínta sex, in 
 persecutióne Wandálica, sub Hunneríco, Rege Ariáno. Hi, cum essent 
 partim Epíscopi Ecclesiárum Dei, partim Presbyteri et Diáconi, associátis 
 sibi turbis fidélium populórum, pro defensióne cathólicæ veritátis in 
 horríbilis erémi exsílium trusi sunt; ex quibus plúrimi, dum crudéliter a 
 Mauris duceréntur, hastílium cuspídibus impúlsi ad curréndum et lapídibus 
 tunsi, álii, ligátis pédibus, velut cadávera per dura et áspera loca tracti 
 et síngulis membris discérpti, ad extrémum, várie excruciáti, martyrium 
 celebrárunt. Erant inter eos præcípui Sacerdótes Dómini, Felix et 
 Cypriánus Epíscopi.
\switchcolumn
\selectlanguage{english}
In Africa, four thousand nine 
 hundred and sixty-six holy confessors and martyrs in the persecution of the 
 Vandals under the Arian king Hunneric. Some of them were bishops of 
 the churches of God, some priests and deacons, and there was a multitude of 
 the faithful who were driven into a frightful wilderness for the defence of 
 the Catholic truth. Many of them were cruelly molested by the Moorish 
 leaders, and with sharp-pointed spears and stones were forced to hasten 
 their march; others, with their feet tied, were dragged like corpses through 
 rough places and were mangled in all their limbs. At the end they were 
 tortured in different manners and won the honours of martyrdom. The 
 principal ones among them were the bishops Felix and Cyprian.
\switchcolumn*
\selectlanguage{latin}
Ravénnæ, via Laurentína, natális sancti Edístii Mártyris.
\switchcolumn
\selectlanguage{english}
At Ravenna, on the Via Laurentina, 
 the birthday of St. Edistus, martyr.
\switchcolumn*
\selectlanguage{latin}
In Lycia sanctæ Domnínæ 
 Mártyris, sub Diocletiáno Imperatóre.
\switchcolumn
\selectlanguage{english}
In Lycia, under Emperor Diocletian, 
 St. Domnina, martyr.
\switchcolumn*
\selectlanguage{latin}
Celénæ, in Pannónia, 
 sancti Maximiliáni, Epíscopi Laureacénsis.
\switchcolumn
\selectlanguage{english}
At Cilli in Austria, St. Maximilian, 
 bishop of Lorsch.
\switchcolumn*
\selectlanguage{latin}
Eboráci, in Anglia, 
 sancti Walfrídi, Epíscopi et Confessóris.
\switchcolumn
\selectlanguage{english}
At York in England, St. Wilfrid, 
 bishop and confessor.
\switchcolumn*
\selectlanguage{latin}
Medioláni sancti Monæ 
 Epíscopi, qui, cum de Epíscopo eligéndo agerétur, cælésti lúmine circumfúsus, 
 eo signo mirabíliter in Pontíficem illíus Ecclésiæ est assúmptus.
\switchcolumn
\selectlanguage{english}
At Milan, St. Monas, bishop. 
 He was chosen as head of that church because a miraculous light from heaven 
 surrounded him when they were deliberating on the choice of a bishop.
\switchcolumn*
\selectlanguage{latin}
Verónæ sancti Salvíni 
 Epíscopi.
\switchcolumn
\selectlanguage{english}
At Verona, St. Salvinus, bishop.
\switchcolumn*
\selectlanguage{latin}
In Syria sancti 
 Eustáchii, Presbyteri et Confessóris.
\switchcolumn
\selectlanguage{english}
In Syria, St. Eustace, priest and 
 confessor.
\switchcolumn*
\selectlanguage{latin}
Asculi, in Picéno, sancti Seraphíni Confessóris, ex Ordine Minórum Capuccinórum, vitæ 
 sanctimónia et humilitáte conspícui; quem Clemens Décimus tértius, Póntifex 
 Máximus, Sanctórum fastis adscrípsit.
\switchcolumn
\selectlanguage{english}
At Ascoli in Piceno, St. Seraphinus, 
 confessor, of the Order of Friars Minor Capuchin, distinguished by his 
 humility and holiness of life. He was enrolled among the saints by the 
 Sovereign Pontiff Clement XIII.
\switchcolumn*
\selectlanguage{latin}
\end{paracol}


% ---- martyrology/mart10/mart1013.htm
\needspace{10\baselineskip}
\begin{paracol}{2}
\selectlanguage{latin}
\begin{center}{\color{gregoriocolor} Tértio Idus Octóbris. 
 Luna\dots\ }\end{center}
\switchcolumn
\selectlanguage{english}
\begin{center}{\color{gregoriocolor} The   13\textsuperscript{th} Day of 
 October. The\dots\ Day of the Moon.}\end{center}
\end{paracol}

\noindent\begin{tabularx}{\linewidth}{*{19}{>{\centering\arraybackslash}X}}
 \textcolor{gregoriocolor}{a} & \textcolor{gregoriocolor}{b} & \textcolor{gregoriocolor}{c} & \textcolor{gregoriocolor}{d} & \textcolor{gregoriocolor}{e} & \textcolor{gregoriocolor}{f} & \textcolor{gregoriocolor}{g} & \textcolor{gregoriocolor}{h} & \textcolor{gregoriocolor}{i} & \textcolor{gregoriocolor}{k} & \textcolor{gregoriocolor}{l} & \textcolor{gregoriocolor}{m} & \textcolor{gregoriocolor}{n} & \textcolor{gregoriocolor}{p} & \textcolor{gregoriocolor}{q} & \textcolor{gregoriocolor}{r} & \textcolor{gregoriocolor}{s} & \textcolor{gregoriocolor}{t} & \textcolor{gregoriocolor}{u} \\
 21 & 22 & 23 & 24 & 25 & 26 & 27 & 28 & 29 & 1 & 2 & 3 & 4 & 5 & 6 & 7 & 8 & 9 & 10 \\
\end{tabularx}
\vspace{0.5\baselineskip}
\noindent\begin{tabularx}{\linewidth}{*{12}{>{\centering\arraybackslash}X}}
 \textcolor{gregoriocolor}{A} & \textcolor{gregoriocolor}{B} & \textcolor{gregoriocolor}{C} & \textcolor{gregoriocolor}{D} & \textcolor{gregoriocolor}{E} & F & \textcolor{gregoriocolor}{F} & \textcolor{gregoriocolor}{G} & \textcolor{gregoriocolor}{H} & \textcolor{gregoriocolor}{M} & \textcolor{gregoriocolor}{N} & \textcolor{gregoriocolor}{P} \\
 11 & 12 & 13 & 14 & 15 & 16 & 15 & 16 & 17 & 18 & 19 & 20 \\
\end{tabularx}

\begin{paracol}{2}
\selectlanguage{latin}
\lettrine[lines=2]{S}{ancti} Eduárdi, Regis 
 Anglórum et Confessóris, qui Nonis Januárii obdormívit in Dómino, sed hac 
 die, ob Translatiónem córporis ejus, potíssimum cólitur.
\switchcolumn
\selectlanguage{english}
\lettrine[lines=2]{S}{t.} Edward, king of England and 
 confessor, who died on the 5th day of January. He is specially 
 honoured on this day because of the translation of his body.
\switchcolumn*
\selectlanguage{latin}
Apud Tróadem, Asiæ 
 minóris urbem, natális sancti Carpi, qui fuit discípulus beáti Pauli 
 Apóstoli.
\switchcolumn
\selectlanguage{english}
At Troas in Asia Minor, the birthday 
 of St. Carpus, a disciple of the blessed apostle Paul.
\switchcolumn*
\selectlanguage{latin}
Córdubæ, in Hispánia, 
 item natális sanctórum Mártyrum Fausti, Januárii et Martiális; qui, primo 
 equúlei pœna cruciáti, deínde, supercíliis rasis, déntibus evúlsis, áuribus 
 quoque et náribus præcísis, ignis passióne martyrium consummárunt.
\switchcolumn
\selectlanguage{english}
At Cordova in Spain, the birthday of 
 the holy martyrs Faustus, Januarius, and Martial. They were first 
 tortured on the rack, their eyebrows were then shaven, their teeth torn out, 
 their ears and noses cut off, and the martyrdom was completed by fire.
\switchcolumn*
\selectlanguage{latin}
Thessalonícæ sancti 
 Floréntii Mártyris, qui, post vária torménta, igne combústus est.
\switchcolumn
\selectlanguage{english}
At Thessalonica, St. Florentius, a 
 martyr, who, after enduring various torments, was burned alive.
\switchcolumn*
\selectlanguage{latin}
Apud Stokeráviam, in 
 Austria, sancti Colmánni Mártyris.
\switchcolumn
\selectlanguage{english}
At Stockerau in Austria, St. Colman, 
 martyr.
\switchcolumn*
\selectlanguage{latin}
Antiochíæ sancti 
 Theóphili Epíscopi, qui, sextus post beátum Petrum Apóstolum, ejúsdem 
 Ecclésiæ Pontificátum ténuit.
\switchcolumn
\selectlanguage{english}
At Antioch, St. Theophilus, the 
 bishop who held the pontificate in that church, the sixth after the blessed 
 apostle Peter.
\switchcolumn*
\selectlanguage{latin}
Turónis, in Gállia, 
 sancti Venántii, Abbátis et Confessóris.
\switchcolumn
\selectlanguage{english}
At Tours in France, St. Venantius, 
 abbot and confessor.
\switchcolumn*
\selectlanguage{latin}
Apud Sublácum, in Látio, 
 sanctæ Chelidóniæ Vírginis.
\switchcolumn
\selectlanguage{english}
At Subiaco in Italy, St. Chelidonia, 
 virgin.
\switchcolumn*
\selectlanguage{latin}
\end{paracol}


% ---- martyrology/mart10/mart1014.htm
\needspace{10\baselineskip}
\begin{paracol}{2}
\selectlanguage{latin}
\begin{center}{\color{gregoriocolor} Prídie Idus Octóbris. 
 Luna\dots\ }\end{center}
\switchcolumn
\selectlanguage{english}
\begin{center}{\color{gregoriocolor} The   14\textsuperscript{th} Day of 
 October. The\dots\ Day of the Moon.}\end{center}
\end{paracol}

\noindent\begin{tabularx}{\linewidth}{*{19}{>{\centering\arraybackslash}X}}
 \textcolor{gregoriocolor}{a} & \textcolor{gregoriocolor}{b} & \textcolor{gregoriocolor}{c} & \textcolor{gregoriocolor}{d} & \textcolor{gregoriocolor}{e} & \textcolor{gregoriocolor}{f} & \textcolor{gregoriocolor}{g} & \textcolor{gregoriocolor}{h} & \textcolor{gregoriocolor}{i} & \textcolor{gregoriocolor}{k} & \textcolor{gregoriocolor}{l} & \textcolor{gregoriocolor}{m} & \textcolor{gregoriocolor}{n} & \textcolor{gregoriocolor}{p} & \textcolor{gregoriocolor}{q} & \textcolor{gregoriocolor}{r} & \textcolor{gregoriocolor}{s} & \textcolor{gregoriocolor}{t} & \textcolor{gregoriocolor}{u} \\
 22 & 23 & 24 & 25 & 26 & 27 & 28 & 29 & 1 & 2 & 3 & 4 & 5 & 6 & 7 & 8 & 9 & 10 & 11 \\
\end{tabularx}
\vspace{0.5\baselineskip}
\noindent\begin{tabularx}{\linewidth}{*{12}{>{\centering\arraybackslash}X}}
 \textcolor{gregoriocolor}{A} & \textcolor{gregoriocolor}{B} & \textcolor{gregoriocolor}{C} & \textcolor{gregoriocolor}{D} & \textcolor{gregoriocolor}{E} & F & \textcolor{gregoriocolor}{F} & \textcolor{gregoriocolor}{G} & \textcolor{gregoriocolor}{H} & \textcolor{gregoriocolor}{M} & \textcolor{gregoriocolor}{N} & \textcolor{gregoriocolor}{P} \\
 12 & 13 & 14 & 15 & 16 & 17 & 16 & 17 & 18 & 19 & 20 & 21 \\
\end{tabularx}

\begin{paracol}{2}
\selectlanguage{latin}
\lettrine[lines=2]{R}{omæ,} via Aurélia, natális beáti Callísti Primi, Papæ et Mártyris; qui, Alexándri Imperatóris 
 jussu, diútius fame in cárcere cruciátus, et quotídie fústibus cæsus, 
 tandem, præcipitátus e fenéstra domus in qua custodiebátur, atque in púteum 
 demérsus, victóriæ triúmphum proméruit.
\switchcolumn
\selectlanguage{english}
\lettrine[lines=2]{A}{t} Rome, on the Aurelian Way, the 
 birthday of blessed Callistus I, pope and martyr. By order of Emperor 
 Alexander, he was kept in prison for a long time without food, and was daily 
 scourged with rods. He was finally hurled from a window of the house 
 in which he had been shut up, and was cast into a well, and thus merited the 
 triumph of victory.
\switchcolumn*
\selectlanguage{latin}
Arímini sancti 
 Gaudéntii, Epíscopi et Mártyris.
\switchcolumn
\selectlanguage{english}
At Rimini, St. Gaudentius, bishop 
 and martyr.
\switchcolumn*
\selectlanguage{latin}
Cæsaréæ, in Palæstína, 
 sanctórum Carpónii, Evarísti et Prisciáni, fratrum beátæ Fortunátæ, qui, 
 gládio juguláti, páriter martyrii corónam percepérunt.
\switchcolumn
\selectlanguage{english}
At Caesarea in Palestine, the Saints 
 Carponius, Evaristus, and Priscian, brothers of blessed Fortunata, who 
 obtained the crown of martyrdom together, their throats being cut with the 
 sword.
\switchcolumn*
\selectlanguage{latin}
Item sanctórum 
 Saturníni et Lupi.
\switchcolumn
\selectlanguage{english}
Also, the Saints Saturninus and 
 Lupus.
\switchcolumn*
\selectlanguage{latin}
Cæsaréæ, in Palæstína, 
 sanctæ Fortunátæ, Vírginis et Mártyris, ac prædicatórum Mártyrum Carpónii, 
 Evarísti et Prisciáni soróris; quæ, in persecutióne Diocletiáni, post 
 equúleum, ignes, béstias et ália torménta superáta, spíritum Deo réddidit. 
 Ipsíus corpus póstea Neápolim, in Campánia, delátum fuit.
\switchcolumn
\selectlanguage{english}
At Caesarea in Palestine, St. 
 Fortunata, virgin and martyr, the sister of the martyrs Carponius, Evaristus, 
 and Priscian. After having been subjected to the rack, to fire, to the 
 teeth of beasts, and other torments during the persecution of Diocletian, she 
 gave up her soul to God. Her body was afterwards conveyed to Naples in 
 Campania.
\switchcolumn*
\selectlanguage{latin}
Tudérti, in Umbria, 
 sancti Fortunáti Epíscopi, qui (ut beátus Gregórius Papa refert) in immúndis 
 spirítibus effugándis imménsæ grátia virtútis emícuit.
\switchcolumn
\selectlanguage{english}
At Todi in Umbria, St. Fortunatus, 
 bishop, who, as is mentioned by blessed Gregory, was endowed with an 
 extraordinary gift for casting out unclean spirits.
\switchcolumn*
\selectlanguage{latin}
Herbípoli, in Germánia, sancti Burchárdi, qui fuit primus illíus civitátis Epíscopus.
\switchcolumn
\selectlanguage{english}
At Wurzburg in Germany, St. Burchard, 
 first bishop of that city.
\switchcolumn*
\selectlanguage{latin}
Brugis Flandrórum 
 sancti Donatiáni, Rheménsis Epíscopi.
\switchcolumn
\selectlanguage{english}
At Bruges in Belgium, St. Donatian, 
 bishop of Rheims.
\switchcolumn*
\selectlanguage{latin}
Tréviris sancti Rústici 
 Epíscopi.
\switchcolumn
\selectlanguage{english}
At Treves, St. Rusticus, bishop.
\switchcolumn*
\selectlanguage{latin}
Lugdúni, in Gállia, 
 sancti Justi, Epíscopi et Confessóris, miræ sanctitátis et prophétici 
 spíritus viri; qui, Episcopátu demísso, in erémum Ægypti, una cum Lectóre 
 suo Viatóre, secéssit, ibíque, cum áliquot annos próximam Angelis egísset 
 vitam, et dignus suórum labórum finis advenísset, corónam justítiæ 
 perceptúrus migrávit ad Dóminum. Ipsíus sanctum corpus, una cum 
 óssibus beáti Viatóris, qui ejúsdem Epíscopi fúerat miníster, Lugdúnum 
 póstea quarto Nonas Septémbris delátum fuit.
\switchcolumn
\selectlanguage{english}
At Lyons in France, St. Justus, 
 bishop and confessor, a man of extraordinary sanctity and endowed with the 
 spirit of prophecy. He resigned his bishopric and retired into a 
 desert in Egypt with his lector Viator. When he had for some years led 
 an almost angelic life, and the end of his meritorious labours had come, he 
 went to our Lord to receive the crown of justice. His holy body and 
 the relics of his lector, blessed Viator, were afterwards taken to Lyons on 
 the 2nd of September.
\switchcolumn*
\selectlanguage{latin}
Eódem die deposítio 
 beáti Domínici Loricáti.
\switchcolumn
\selectlanguage{english}
On the same day, the death of 
 blessed Dominic Loricatus.
\switchcolumn*
\selectlanguage{latin}
Arpíni, in Látio, 
 sancti Bernárdi Confessóris.
\switchcolumn
\selectlanguage{english}
At Arpino in Italy, St. Bernard, 
 confessor.
\switchcolumn*
\selectlanguage{latin}
\end{paracol}


% ---- martyrology/mart10/mart1015.htm
\needspace{10\baselineskip}
\begin{paracol}{2}
\selectlanguage{latin}
\begin{center}{\color{gregoriocolor} Idibus Octóbris. 
 Luna\dots\ }\end{center}
\switchcolumn
\selectlanguage{english}
\begin{center}{\color{gregoriocolor} The   15\textsuperscript{th} Day of 
 October. The\dots\ Day of the Moon.}\end{center}
\end{paracol}

\noindent\begin{tabularx}{\linewidth}{*{19}{>{\centering\arraybackslash}X}}
 \textcolor{gregoriocolor}{a} & \textcolor{gregoriocolor}{b} & \textcolor{gregoriocolor}{c} & \textcolor{gregoriocolor}{d} & \textcolor{gregoriocolor}{e} & \textcolor{gregoriocolor}{f} & \textcolor{gregoriocolor}{g} & \textcolor{gregoriocolor}{h} & \textcolor{gregoriocolor}{i} & \textcolor{gregoriocolor}{k} & \textcolor{gregoriocolor}{l} & \textcolor{gregoriocolor}{m} & \textcolor{gregoriocolor}{n} & \textcolor{gregoriocolor}{p} & \textcolor{gregoriocolor}{q} & \textcolor{gregoriocolor}{r} & \textcolor{gregoriocolor}{s} & \textcolor{gregoriocolor}{t} & \textcolor{gregoriocolor}{u} \\
 23 & 24 & 25 & 26 & 27 & 28 & 29 & 1 & 2 & 3 & 4 & 5 & 6 & 7 & 8 & 9 & 10 & 11 & 12 \\
\end{tabularx}
\vspace{0.5\baselineskip}
\noindent\begin{tabularx}{\linewidth}{*{12}{>{\centering\arraybackslash}X}}
 \textcolor{gregoriocolor}{A} & \textcolor{gregoriocolor}{B} & \textcolor{gregoriocolor}{C} & \textcolor{gregoriocolor}{D} & \textcolor{gregoriocolor}{E} & F & \textcolor{gregoriocolor}{F} & \textcolor{gregoriocolor}{G} & \textcolor{gregoriocolor}{H} & \textcolor{gregoriocolor}{M} & \textcolor{gregoriocolor}{N} & \textcolor{gregoriocolor}{P} \\
 13 & 14 & 15 & 16 & 17 & 18 & 17 & 18 & 19 & 20 & 21 & 22 \\
\end{tabularx}

\begin{paracol}{2}
\selectlanguage{latin}
\lettrine[lines=2]{A}{lbæ,} in Hispánia, 
 sanctæ Terésiæ Vírginis, quæ Fratrum ac Sorórum Ordinis Carmelitárum 
 arctióris observántiæ mater éxstitit et magístra.
\switchcolumn
\selectlanguage{english}
\lettrine[lines=2]{A}{t} Avila in Spain, the virgin St. 
 Teresa, mother and mistress of the Brothers and Sisters of the Carmelite 
 Order of the Strict Observance.
\switchcolumn*
\selectlanguage{latin}
Cracóviæ, in Polónia, 
 natális sanctæ Hedwígis Víduæ, Polonórum Ducíssæ, quæ, páuperum obséquio 
 dédita, étiam miráculis cláruit; et a Cleménte Quarto, Pontífice Máximo, 
 Sanctórum número adscrípta est. Ipsíus autem festívitas sequénti die 
 celebrátur.
\switchcolumn
\selectlanguage{english}
At Cracow in Poland, St. Hedwig, 
 duchess of Poland, who devoted herself to the service of the poor, and was 
 renowned for miracles. She was inscribed among the saints by Pope 
 Clement IV. Her feast is celebrated on the following day.
\switchcolumn*
\selectlanguage{latin}
Romæ, via Aurélia, sancti Fortunáti Mártyris.
\switchcolumn
\selectlanguage{english}
At Rome, on the Aurelian Way, St. 
 Fortunatus, martyr.
\switchcolumn*
\selectlanguage{latin}
In Borússia sancti 
 Brunónis, Epíscopi Ruthenórum et Mártyris; qui, Evangélium in ea regióne 
 prædicans, ab ímpiis tentus est, ac, mánibus pedibúsque præcísis, cápite 
 truncátus.
\switchcolumn
\selectlanguage{english}
In Prussia, St. Bruno, bishop of the 
 Ruthenians and martyr. While preaching the Gospel in that region he 
 was arrested by impious men, his hands and feet were cut off, and he was 
 then beheaded.
\switchcolumn*
\selectlanguage{latin}
Apud Colóniam 
 Agrippínam natális sanctórum trecentórum Mártyrum, qui, in Maximiáni 
 persecutióne, cursum sui agónis complevérunt.
\switchcolumn
\selectlanguage{english}
At Cologne, the birthday of three 
 hundred holy martyrs, who met their trials in the persecution of Maximian.
\switchcolumn*
\selectlanguage{latin}
Carthágine sancti 
 Agiléi Mártyris, in cujus die natáli sanctus Augustínus tractátum de ipso ad 
 pópulum hábuit.
\switchcolumn
\selectlanguage{english}
At Carthage, St. Agileus, martyr, on 
 whose birthday St. Augustine delivered a discourse to the people concerning 
 him.
\switchcolumn*
\selectlanguage{latin}
Lugdúni, in Gállia, 
 sancti Antíochi Epíscopi, qui, strénue administráto Pontifícii cúlmine ad 
 quod assúmptus fúerat, regnum cæléste adéptus est.
\switchcolumn
\selectlanguage{english}
At Lyons in France, St. Antiochus, 
 bishop, who entered the heavenly kingdom after having courageously 
 fulfilled the duties of the high station to which he had been called.
\switchcolumn*
\selectlanguage{latin}
Tréviris sancti Sevéri, 
 Epíscopi et Confessóris.
\switchcolumn
\selectlanguage{english}
At Treves, St. Severus, bishop and 
 confessor.
\switchcolumn*
\selectlanguage{latin}
Argentoráti sanctæ 
 Auréliæ Vírginis.
\switchcolumn
\selectlanguage{english}
At Strasbourg, St. Aurelia, virgin.
\switchcolumn*
\selectlanguage{latin}
In Germánia sanctæ 
 Theclæ, Abbatíssæ et Vírginis, quæ, monastériis Kitzíngæ et Ochsenfúrti 
 præpósita, multis cumuláta méritis in cælum migrávit.
\switchcolumn
\selectlanguage{english}
In Germany, St. Thecla, abbess and 
 virgin. She governed the convents of Kitzingen and Ochsenfurt, and 
 departed to heaven filled with merits.
\switchcolumn*
\selectlanguage{latin}
\end{paracol}


% ---- martyrology/mart10/mart1016.htm
\needspace{10\baselineskip}
\begin{paracol}{2}
\selectlanguage{latin}
\begin{center}{\color{gregoriocolor} Décimo séptimo Kaléndas Novémbris. 
 Luna\dots\ }\end{center}
\switchcolumn
\selectlanguage{english}
\begin{center}{\color{gregoriocolor} The   16\textsuperscript{th} Day of 
 October. The\dots\ Day of the Moon.}\end{center}
\end{paracol}

\noindent\begin{tabularx}{\linewidth}{*{19}{>{\centering\arraybackslash}X}}
 \textcolor{gregoriocolor}{a} & \textcolor{gregoriocolor}{b} & \textcolor{gregoriocolor}{c} & \textcolor{gregoriocolor}{d} & \textcolor{gregoriocolor}{e} & \textcolor{gregoriocolor}{f} & \textcolor{gregoriocolor}{g} & \textcolor{gregoriocolor}{h} & \textcolor{gregoriocolor}{i} & \textcolor{gregoriocolor}{k} & \textcolor{gregoriocolor}{l} & \textcolor{gregoriocolor}{m} & \textcolor{gregoriocolor}{n} & \textcolor{gregoriocolor}{p} & \textcolor{gregoriocolor}{q} & \textcolor{gregoriocolor}{r} & \textcolor{gregoriocolor}{s} & \textcolor{gregoriocolor}{t} & \textcolor{gregoriocolor}{u} \\
 24 & 25 & 26 & 27 & 28 & 29 & 1 & 2 & 3 & 4 & 5 & 6 & 7 & 8 & 9 & 10 & 11 & 12 & 13 \\
\end{tabularx}
\vspace{0.5\baselineskip}
\noindent\begin{tabularx}{\linewidth}{*{12}{>{\centering\arraybackslash}X}}
 \textcolor{gregoriocolor}{A} & \textcolor{gregoriocolor}{B} & \textcolor{gregoriocolor}{C} & \textcolor{gregoriocolor}{D} & \textcolor{gregoriocolor}{E} & F & \textcolor{gregoriocolor}{F} & \textcolor{gregoriocolor}{G} & \textcolor{gregoriocolor}{H} & \textcolor{gregoriocolor}{M} & \textcolor{gregoriocolor}{N} & \textcolor{gregoriocolor}{P} \\
 14 & 15 & 16 & 17 & 18 & 19 & 18 & 19 & 20 & 21 & 22 & 23 \\
\end{tabularx}

\begin{paracol}{2}
\selectlanguage{latin}
\lettrine[lines=2]{S}{anctæ} Hedwígis Víduæ, 
 Polonórum Ducíssæ, quæ prídie hujus diéi obdormívit in Dómino.
\switchcolumn
\selectlanguage{english}
\lettrine[lines=2]{S}{t.} Hedwig, widow, duchess of 
 Poland, who went to her rest in the Lord on the day previous.
\switchcolumn*
\selectlanguage{latin}
In monastério Dervénsi, 
 in Gállia, sancti Berchárii, Abbátis et Mártyris.
\switchcolumn
\selectlanguage{english}
In the monastery of Moutier-en-Der, 
 in France, St. Bercharius, abbot and martyr.
\switchcolumn*
\selectlanguage{latin}
In Africa sanctórum 
 Mártyrum ducentórum septuagínta, páriter coronatórum.
\switchcolumn
\selectlanguage{english}
In Africa, two hundred and seventy 
 holy martyrs, crowned together.
\switchcolumn*
\selectlanguage{latin}
Ibídem sanctórum 
 Martiniáni et Saturiáni, cum duóbus eórum frátribus; qui, témpore Wandálicæ 
 persecutiónis, sub Rege Ariáno Genseríco, cum servi essent cujúsdam Wándali, 
 et a sancta Máxima Vírgine, ipsórum consérva, ad Christi fidem convérsi 
 fuíssent, omnes ab hærético illórum dómino, pro constántia fídei cathólicæ, 
 primum nodósis fústibus cæsi sunt et usque ad ossa laniáti. Sed, cum 
 tália multo témpore pateréntur, et sequénti die nihilóminus redderéntur 
 semper incólumes, exsílio tandem relegántur; ubi, cum multos barbarórum ad 
 Christi convertíssent fidem, et a Románo Pontífice Presbyterum aliósque 
 minístros, qui eos baptizárent, obtinuíssent, novíssime, vinctis pédibus 
 post terga curréntium quadrigárum, inter spinósa loca silvárum jussi sunt 
 páriter interíre. Máxima vero, post multos superátos agónes divínitus 
 liberáta, in monastério, multárum Vírginum Mater, sancto fine quiévit.
\switchcolumn
\selectlanguage{english}
Likewise, the Saints Martinian and 
 Saturnian, with their two brothers. While the persecution of the 
 Vandals was raging in the reign of the Arian king Genseric, they were slaves 
 to a man of that race. They were converted to the faith of Christ by 
 Maxima, a slave like themselves, and they manifested their attachment to the 
 truth with such courage that they were beaten with rough clubs and lacerated 
 in all parts of their bodies to the very bones. Although this 
 barbarous treatment was continued for a considerable period, their wounds 
 were each time healed overnight. They were at length sent into exile 
 where they converted many barbarians to the faith, and obtained from the 
 Roman Pontiff a priest and other ministers to baptize them. Finally 
 there were condemned to die by having their feet tied behind running 
 chariots and being dragged through thorns. Maxima, after enduring many 
 tribulations, was miraculously delivered and became the superior of a large 
 monastery of virgins, where she ended her days in peace.
\switchcolumn*
\selectlanguage{latin}
Item sanctórum 
 Saturníni, Nérei et aliórum trecentórum sexagínta quinque Mártyrum.
\switchcolumn
\selectlanguage{english}
Also, the Saints Saturninus, Nereus, 
 and three hundred and sixty-five other martyrs.
\switchcolumn*
\selectlanguage{latin}
Colóniæ Agrippínæ 
 sancti Elíphii Mártyris, sub Juliáno Apóstata.
\switchcolumn
\selectlanguage{english}
At Cologne, under Julian the 
 Apostate, the martyr St. Eliphius.
\switchcolumn*
\selectlanguage{latin}
In território Bituricénsi sancti Ambrósii, Epíscopi Caturcénsis.
\switchcolumn
\selectlanguage{english}
Near Bourges, St. Ambrose, bishop of 
 Cahors.
\switchcolumn*
\selectlanguage{latin}
Mogúntiæ sancti Lulli, 
 Epíscopi et Confessóris.
\switchcolumn
\selectlanguage{english}
At Mainz, St. Lullus, bishop and 
 confessor.
\switchcolumn*
\selectlanguage{latin}
Tréviris sancti 
 Florentíni Epíscopi.
\switchcolumn
\selectlanguage{english}
At Treves, St. Florentinus, bishop.
\switchcolumn*
\selectlanguage{latin}
Apud Arbónam, in 
 Germánia, sancti Galli Abbátis, qui fuit discípulus beáti Columbáni.
\switchcolumn
\selectlanguage{english}
At Arbon in Germany, St. Gall, 
 abbot, a disciple of blessed Columban.
\switchcolumn*
\selectlanguage{latin}
Muri, in Lucánia, 
 sancti Gerárdi Majélla, Confessóris, Láici proféssi Congregatiónis a 
 sanctíssimo Redemptóre nuncupátæ, quem, miráculis clarum, Pius Décimus, 
 Póntifex Máximus, Sanctórum albo accénsuit.
\switchcolumn
\selectlanguage{english}
At Muro in Italy, St. Gerard Majella, 
 confessor and professed lay brother of the Congregation of the Most Holy 
 Redeemer. Renowned for miracles, he was added to the list of the 
 saints by Pope Pius X.
\switchcolumn*
\selectlanguage{latin}
\end{paracol}


% ---- martyrology/mart10/mart1017.htm
\needspace{10\baselineskip}
\begin{paracol}{2}
\selectlanguage{latin}
\begin{center}{\color{gregoriocolor} Sextodécimo Kaléndas Novémbris. 
 Luna\dots\ }\end{center}
\switchcolumn
\selectlanguage{english}
\begin{center}{\color{gregoriocolor} The   17\textsuperscript{th} Day of 
 October. The\dots\ Day of the Moon.}\end{center}
\end{paracol}

\noindent\begin{tabularx}{\linewidth}{*{19}{>{\centering\arraybackslash}X}}
 \textcolor{gregoriocolor}{a} & \textcolor{gregoriocolor}{b} & \textcolor{gregoriocolor}{c} & \textcolor{gregoriocolor}{d} & \textcolor{gregoriocolor}{e} & \textcolor{gregoriocolor}{f} & \textcolor{gregoriocolor}{g} & \textcolor{gregoriocolor}{h} & \textcolor{gregoriocolor}{i} & \textcolor{gregoriocolor}{k} & \textcolor{gregoriocolor}{l} & \textcolor{gregoriocolor}{m} & \textcolor{gregoriocolor}{n} & \textcolor{gregoriocolor}{p} & \textcolor{gregoriocolor}{q} & \textcolor{gregoriocolor}{r} & \textcolor{gregoriocolor}{s} & \textcolor{gregoriocolor}{t} & \textcolor{gregoriocolor}{u} \\
 25 & 26 & 27 & 28 & 29 & 1 & 2 & 3 & 4 & 5 & 6 & 7 & 8 & 9 & 10 & 11 & 12 & 13 & 14 \\
\end{tabularx}
\vspace{0.5\baselineskip}
\noindent\begin{tabularx}{\linewidth}{*{12}{>{\centering\arraybackslash}X}}
 \textcolor{gregoriocolor}{A} & \textcolor{gregoriocolor}{B} & \textcolor{gregoriocolor}{C} & \textcolor{gregoriocolor}{D} & \textcolor{gregoriocolor}{E} & F & \textcolor{gregoriocolor}{F} & \textcolor{gregoriocolor}{G} & \textcolor{gregoriocolor}{H} & \textcolor{gregoriocolor}{M} & \textcolor{gregoriocolor}{N} & \textcolor{gregoriocolor}{P} \\
 15 & 16 & 17 & 18 & 19 & 20 & 19 & 20 & 21 & 22 & 23 & 24 \\
\end{tabularx}

\begin{paracol}{2}
\selectlanguage{latin}
\lettrine[lines=2]{P}{arédii,} in diœcési 
 Augustodunénsi, sanctæ Margarítæ Maríæ Alacóque, quæ, Ordinem Visitatiónis 
 beátæ Maríæ Vírginis proféssa, exímiis in devotióne erga sacratíssimum Cor 
 Jesu propagánda et público ejúsdem cultu provehéndo méritis excélluit; 
 atque in sanctárum Vírginum album a Benedícto Papa Décimo quinto reláta fuit.
\switchcolumn
\selectlanguage{english}
\lettrine[lines=2]{A}{t} Paray, in the diocese of Autun, 
 St. Margaret Mary Alacoque. She made her profession in the Order of 
 the Visitation of Blessed Mary the Virgin, and she excelled with great merit 
 in spreading devotion to the Most Sacred Heart of Jesus and in furthering 
 its public veneration. Pope Benedict XV added her name to the list of 
 holy virgins.
\switchcolumn*
\selectlanguage{latin}
Antiochíæ natális 
 sancti Herónis, qui fuit discípulus beáti Ignátii; atque, post eum factus 
 Epíscopus, viam magístri pius imitátor est secútus, et pro commendáto sibi 
 grege amátor Christi occúbuit.
\switchcolumn
\selectlanguage{english}
At Antioch, the birthday of St. 
 Heron, a disciple of blessed Ignatius. Being made bishop after him, he 
 religiously followed his master's footsteps, and, as a true lover of Christ, 
 died for the flock entrusted to his keeping.
\switchcolumn*
\selectlanguage{latin}
Eódem die pássio 
 sanctórum Victóris, Alexándri et Mariáni.
\switchcolumn
\selectlanguage{english}
The same day, the martyrdom of the 
 Saints Victor, Alexander, and Marian.
\switchcolumn*
\selectlanguage{latin}
In Pérside sanctæ 
 Maméltæ Mártyris, quæ, a cultu idolórum ad fidem Angélico mónitu convérsa, a 
 Gentílibus lapidáta est et in profúndum lacum demérsa.
\switchcolumn
\selectlanguage{english}
In Persia, St. Mamelta, martyr. 
 He was converted from idolatry to the faith by the warning of an angel, and 
 was later stoned by heathens and cast into a deep lake.
\switchcolumn*
\selectlanguage{latin}
Aráusicæ, in Gálliis, 
 sancti Floréntii Epíscopi, qui, multis clarus virtútibus, quiévit in pace.
\switchcolumn
\selectlanguage{english}
At Orange in France, St. Florentinus, 
 bishop, who died leaving a reputation for many virtues.
\switchcolumn*
\selectlanguage{latin}
\end{paracol}


% ---- martyrology/mart10/mart1018.htm
\needspace{10\baselineskip}
\begin{paracol}{2}
\selectlanguage{latin}
\begin{center}{\color{gregoriocolor} Quintodécimo Kaléndas Novémbris. 
 Luna\dots\ }\end{center}
\switchcolumn
\selectlanguage{english}
\begin{center}{\color{gregoriocolor} The   18\textsuperscript{th} Day of 
 October. The\dots\ Day of the Moon.}\end{center}
\end{paracol}

\noindent\begin{tabularx}{\linewidth}{*{19}{>{\centering\arraybackslash}X}}
 \textcolor{gregoriocolor}{a} & \textcolor{gregoriocolor}{b} & \textcolor{gregoriocolor}{c} & \textcolor{gregoriocolor}{d} & \textcolor{gregoriocolor}{e} & \textcolor{gregoriocolor}{f} & \textcolor{gregoriocolor}{g} & \textcolor{gregoriocolor}{h} & \textcolor{gregoriocolor}{i} & \textcolor{gregoriocolor}{k} & \textcolor{gregoriocolor}{l} & \textcolor{gregoriocolor}{m} & \textcolor{gregoriocolor}{n} & \textcolor{gregoriocolor}{p} & \textcolor{gregoriocolor}{q} & \textcolor{gregoriocolor}{r} & \textcolor{gregoriocolor}{s} & \textcolor{gregoriocolor}{t} & \textcolor{gregoriocolor}{u} \\
 26 & 27 & 28 & 29 & 1 & 2 & 3 & 4 & 5 & 6 & 7 & 8 & 9 & 10 & 11 & 12 & 13 & 14 & 15 \\
\end{tabularx}
\vspace{0.5\baselineskip}
\noindent\begin{tabularx}{\linewidth}{*{12}{>{\centering\arraybackslash}X}}
 \textcolor{gregoriocolor}{A} & \textcolor{gregoriocolor}{B} & \textcolor{gregoriocolor}{C} & \textcolor{gregoriocolor}{D} & \textcolor{gregoriocolor}{E} & F & \textcolor{gregoriocolor}{F} & \textcolor{gregoriocolor}{G} & \textcolor{gregoriocolor}{H} & \textcolor{gregoriocolor}{M} & \textcolor{gregoriocolor}{N} & \textcolor{gregoriocolor}{P} \\
 16 & 17 & 18 & 19 & 20 & 21 & 20 & 21 & 22 & 23 & 24 & 25 \\
\end{tabularx}

\begin{paracol}{2}
\selectlanguage{latin}
\lettrine[lines=2]{I}{n} Bithynia natális 
 beáti Lucæ Evangelístæ, qui, multa passus pro Christi nómine, obiit Spíritu 
 Sancto plenus. Ipsíus autem ossa póstea Constantinópolim transláta 
 sunt, et inde Patávium deláta.
\switchcolumn
\selectlanguage{english}
\lettrine[lines=2]{I}{n} Bithynia, the birthday of St. 
 Luke the Evangelist. He died, filled with the Holy Ghost, after having 
 suffered much for the Name of Christ. His relics were translated to 
 Constantinople, and thence taken to Pavia.
\switchcolumn*
\selectlanguage{latin}
Romæ item natális 
 sancti Pauli a Cruce, Presbyteri et Confessóris; qui Congregatiónis a Cruce 
 et Passióne Dómini nostri Jesu Christi nuncupátæ Institútor fuit. 
 Ipsum vero, mira innocéntia ac pæniténtia conspícuum et singulári in 
 Christum crucifíxum caritáte incénsum, Pius Papa Nonus fastis Sanctórum 
 adjúnxit, et ejúsdem festivitátem quarto Kaléndas Maji recoléndam indíxit.
\switchcolumn
\selectlanguage{english}
At Rome, the birthday of St. Paul of 
 the Cross, priest, confessor, and founder of the Congregation of the Cross 
 and Passion of our Lord Jesus Christ. Known for his remarkable 
 innocency of life and his penitential spirit, and aflame with love for 
 Christ crucified, he was canonized by Pope Pius IX, and the 28th of April 
 was assigned as his feast day.
\switchcolumn*
\selectlanguage{latin}
Arénis, in Hispánia, 
 natális páriter sancti Petri de Alcántara, Sacerdótis ex Ordine Minórum et 
 Confessóris; quem, propter admirábilem pæniténtiam múltaque mirácula, 
 Clemens Nonus, Póntifex Máximus, Sanctórum número adscrípsit. Ejus 
 autem festum sequénti die celebrátur.
\switchcolumn
\selectlanguage{english}
At Arenas in Spain, the birthday of 
 St. Peter of Alcantara, confessor and priest of the Order of Friars Minor. 
 He was canonized by Pope Clement IX because of his admirable penance and 
 many miracles, and his feast is observed on the day following.
\switchcolumn*
\selectlanguage{latin}
Antiochíæ sancti 
 Asclepíadis Epíscopi, qui fuit unus ex præcláro illórum Mártyrum número, qui 
 glorióse sub Macríno passi sunt.
\switchcolumn
\selectlanguage{english}
At Antioch, the bishop St. 
 Asclepiades, who was one of the celebrated band of martyrs who suffered so 
 gloriously under Macrinus.
\switchcolumn*
\selectlanguage{latin}
Neocæsaréæ, in Ponto, 
 sancti Athenodóri Epíscopi, qui fuit frater sancti Gregórii Thaumatúrgi; et, 
 doctrína clarus, in persecutióne Aureliáni, martyrium consummávit.
\switchcolumn
\selectlanguage{english}
At Neocaesarea in Pontus, the holy 
 and learned Bishop Athenodorus, brother of St. Gregory Thaumaturgus, who 
 underwent martyrdom in the persecution of Aurelian.
\switchcolumn*
\selectlanguage{latin}
Sinomovíci, in 
 território Bellovacénsi, sancti Justi Mártyris, qui adhuc puer, in 
 persecutióne Diocletiáni Imperatóris, sub Rictiováro Præside, cápite 
 amputátus est.
\switchcolumn
\selectlanguage{english}
At Louvres, in the diocese of 
 Beauvais, St. Justus, martyr, who, being but a boy, was put to death in the 
 persecution of Diocletian, under the governor Rictiovarus.
\switchcolumn*
\selectlanguage{latin}
Romæ sanctæ Tryphóniæ, 
 quæ Décii Cæsaris quondam uxor ac sanctæ Vírginis et Mártyris Cyrillæ mater 
 éxstitit; cujus corpus in crypta, juxta sanctum Hippólytum, sepúltum est.
\switchcolumn
\selectlanguage{english}
At Rome, St. Tryphonia, at one time 
 the wife of Caesar Decius, the mother of St. Cyrilla, virgin and martyr. 
 She was buried in a crypt, near that of St. Hippolytus.
\switchcolumn*

% FIX: There should be special wording for the USA
\selectlanguage{latin}
Apud Auriesville, in 
 statu Neo-Eboracénsi, sanctórum Mártyrum e Societáte Jesu, Isaáci Jogues, 
 Sacerdotis, et Joánnis de La Lande, Coadjutóris temporális, qui hac et 
 sequénti die ab Iroquénsibus dire necáti sunt, eódem loco ubi, paucis ante 
 annis, Renátus Goupil, et ipse Coadjútor temporális, martyrii palmam 
 consecútus fúerat.
\switchcolumn
\selectlanguage{english}
At Auriesville, in the state of New 
 York, the birthday of the holy martyrs Isaac Jogues, priest of the Society 
 of Jesus, and John de la Lande, a temporary helper to the same Society, who 
 came from France to teach the faith. On this and the following day 
 they were cruelly tortured and killed by the Iroquois in the same place 
 where, a few years before, one of the companions, René Goupil, also a 
 temporary assistant, had received the palm of martyrdom.
\switchcolumn*
\selectlanguage{latin}
In fínibus Edessénæ 
 regiónis, in Mesopotámia, commemorátio sancti Juliáni Eremítæ, cognoménto Sabæ, de quo 
 ágitur étiam sextodécimo Kaléndas Februárii.
\switchcolumn
\selectlanguage{english}
In Mesopotamia, in the neighbourhood 
 of Edessa, the commemoration of St. Julian the Hermit, surnamed Sabas, who 
 is mentioned also on the 17th of January.
\switchcolumn*
\selectlanguage{latin}
\end{paracol}


% ---- martyrology/mart10/mart1019.htm
\needspace{10\baselineskip}
\begin{paracol}{2}
\selectlanguage{latin}
\begin{center}{\color{gregoriocolor} Quartodécimo Kaléndas Novémbris. 
 Luna\dots\ }\end{center}
\switchcolumn
\selectlanguage{english}
\begin{center}{\color{gregoriocolor} The   19\textsuperscript{th} Day of 
 October. The\dots\ Day of the Moon.}\end{center}
\end{paracol}

\noindent\begin{tabularx}{\linewidth}{*{19}{>{\centering\arraybackslash}X}}
 \textcolor{gregoriocolor}{a} & \textcolor{gregoriocolor}{b} & \textcolor{gregoriocolor}{c} & \textcolor{gregoriocolor}{d} & \textcolor{gregoriocolor}{e} & \textcolor{gregoriocolor}{f} & \textcolor{gregoriocolor}{g} & \textcolor{gregoriocolor}{h} & \textcolor{gregoriocolor}{i} & \textcolor{gregoriocolor}{k} & \textcolor{gregoriocolor}{l} & \textcolor{gregoriocolor}{m} & \textcolor{gregoriocolor}{n} & \textcolor{gregoriocolor}{p} & \textcolor{gregoriocolor}{q} & \textcolor{gregoriocolor}{r} & \textcolor{gregoriocolor}{s} & \textcolor{gregoriocolor}{t} & \textcolor{gregoriocolor}{u} \\
 27 & 28 & 29 & 1 & 2 & 3 & 4 & 5 & 6 & 7 & 8 & 9 & 10 & 11 & 12 & 13 & 14 & 15 & 16 \\
\end{tabularx}
\vspace{0.5\baselineskip}
\noindent\begin{tabularx}{\linewidth}{*{12}{>{\centering\arraybackslash}X}}
 \textcolor{gregoriocolor}{A} & \textcolor{gregoriocolor}{B} & \textcolor{gregoriocolor}{C} & \textcolor{gregoriocolor}{D} & \textcolor{gregoriocolor}{E} & F & \textcolor{gregoriocolor}{F} & \textcolor{gregoriocolor}{G} & \textcolor{gregoriocolor}{H} & \textcolor{gregoriocolor}{M} & \textcolor{gregoriocolor}{N} & \textcolor{gregoriocolor}{P} \\
 17 & 18 & 19 & 20 & 21 & 22 & 21 & 22 & 23 & 24 & 25 & 26 \\
\end{tabularx}

\begin{paracol}{2}
\selectlanguage{latin}
\lettrine[lines=2]{S}{ancti} Petri de 
 Alcántara, Sacerdótis ex Ordine Minórum et Confessóris, qui migrávit in 
 cælum prídie hujus diéi.
\switchcolumn
\selectlanguage{english}
\lettrine[lines=2]{S}{t.} Peter of Alcantara, priest of 
 the Order of Friars Minor and confessor, whose birthday was mentioned in the 
 day previous to this.
\switchcolumn*
\selectlanguage{latin}
Romæ natális sanctórum 
 Mártyrum Ptolomǽi et Lúcii, sub Marco Antoníno. Horum prior (ut 
	scribit Justínus Martyr), cum impud\'icam mulíerem ad Christi convertísset 
 fidem, et castitátem cólere docuísset, ídeo, ab impúro viro apud Præféctum 
 Urbícium accusátus, multo témpore squalóre cárceris macerátus est, et ad 
 últimum, cum de Christi magistério pública confessióne testarétur, jussus 
 est duci ad mortem; Lúcius quoque, cum Urbícii senténtiam improbáret et se 
 Christiánum líbere faterétur, símilem senténtiam excépit; quibus et álius 
 tértius adjúnctus est, qui étiam eódem supplício damnátus fuit.
\switchcolumn
\selectlanguage{english}
At Rome, the birthday of the holy 
 martyrs Ptolemy and Lucius, in the time of Marcus Antoninus. The 
 former, as we learn from the martyr Justin, converted a certain immodest 
 woman to the faith of Christ and induced her to practice chastity. He 
 was accused by an evil man before the prefect Urbicius and made to undergo a 
 long imprisonment in a foul dungeon. At length, because he declared by 
 a public confession that Christ was his master, he was led to execution. 
 Lucius protested against the sentence of Urbicius, and freely proclaimed 
 himself to be a Christian, whereby he received the same sentence. To 
 them was added still a third martyr, who was condemned to suffer a like 
 punishment.
\switchcolumn*
\selectlanguage{latin}
Antiochíæ sanctórum 
 Mártyrum Beroníci, Pelágiæ Vírginis, et aliórum quadragínta novem.
\switchcolumn
\selectlanguage{english}
At Antioch, the holy martyrs 
 Beronicus, the virgin Pelagia, and forty-nine others.
\switchcolumn*
\selectlanguage{latin}
In Ægypto sancti Vari 
 mílitis, qui, sub Maximíno Imperatóre, dum sanctos septem Mónachos in cárcere deténtos visitáret atque refíceret, vóluit, uno ex ipsis defúncto, 
 in ejus locum subrogári; atque ita cum illis, sævíssima passus, martyrii 
 palmam adéptus est.
\switchcolumn
\selectlanguage{english}
In Egypt, St. Varus, a soldier, who, 
 under Emperor Maximian, visited and comforted seven holy monks who were kept 
 in prison. When one of them died he wished to be accepted in his 
 place, and after suffering most cruel torments with them he obtained the 
 palm of martyrdom.
\switchcolumn*
\selectlanguage{latin}
Ebróicis, in Gállia, 
 sancti Aquilíni, Epíscopi et Confessóris.
\switchcolumn
\selectlanguage{english}
At Evreux in France, St. Aquilinus, 
 bishop and confessor.
\switchcolumn*
\selectlanguage{latin}
In território 
 Aurelianénsi deposítio sancti Veráni Epíscopi.
\switchcolumn
\selectlanguage{english}
In the diocese of Orleans, the death 
 of St. Veranus, bishop.
\switchcolumn*
\selectlanguage{latin}
Apud Salérnum sancti 
 Eustérii Epíscopi.
\switchcolumn
\selectlanguage{english}
At Salerno, St. Eusterius, bishop.
\switchcolumn*
\selectlanguage{latin}
In monastério Silvæ 
 Necténsis, in Hibérnia, sancti Ethbíni Abbátis.
\switchcolumn
\selectlanguage{english}
In Ireland, in the monastery of the 
 Forest of Kildare, St. Ethbin, abbot.
\switchcolumn*
\selectlanguage{latin}
Oxónii, in Anglia, 
 sanctæ Fredeswíndæ Vírginis.
\switchcolumn
\selectlanguage{english}
At Oxford in England, St. Frideswide, 
 virgin.
\switchcolumn*
\selectlanguage{latin}
\end{paracol}


% ---- martyrology/mart10/mart1020.htm
\needspace{10\baselineskip}
\begin{paracol}{2}
\selectlanguage{latin}
\begin{center}{\color{gregoriocolor} Tertiodécimo Kaléndas Novémbris. 
 Luna\dots\ }\end{center}
\switchcolumn
\selectlanguage{english}
\begin{center}{\color{gregoriocolor} The   20\textsuperscript{th} Day of 
 October. The\dots\ Day of the Moon.}\end{center}
\end{paracol}

\noindent\begin{tabularx}{\linewidth}{*{19}{>{\centering\arraybackslash}X}}
 \textcolor{gregoriocolor}{a} & \textcolor{gregoriocolor}{b} & \textcolor{gregoriocolor}{c} & \textcolor{gregoriocolor}{d} & \textcolor{gregoriocolor}{e} & \textcolor{gregoriocolor}{f} & \textcolor{gregoriocolor}{g} & \textcolor{gregoriocolor}{h} & \textcolor{gregoriocolor}{i} & \textcolor{gregoriocolor}{k} & \textcolor{gregoriocolor}{l} & \textcolor{gregoriocolor}{m} & \textcolor{gregoriocolor}{n} & \textcolor{gregoriocolor}{p} & \textcolor{gregoriocolor}{q} & \textcolor{gregoriocolor}{r} & \textcolor{gregoriocolor}{s} & \textcolor{gregoriocolor}{t} & \textcolor{gregoriocolor}{u} \\
 28 & 29 & 1 & 2 & 3 & 4 & 5 & 6 & 7 & 8 & 9 & 10 & 11 & 12 & 13 & 14 & 15 & 16 & 17 \\
\end{tabularx}
\vspace{0.5\baselineskip}
\noindent\begin{tabularx}{\linewidth}{*{12}{>{\centering\arraybackslash}X}}
 \textcolor{gregoriocolor}{A} & \textcolor{gregoriocolor}{B} & \textcolor{gregoriocolor}{C} & \textcolor{gregoriocolor}{D} & \textcolor{gregoriocolor}{E} & F & \textcolor{gregoriocolor}{F} & \textcolor{gregoriocolor}{G} & \textcolor{gregoriocolor}{H} & \textcolor{gregoriocolor}{M} & \textcolor{gregoriocolor}{N} & \textcolor{gregoriocolor}{P} \\
 18 & 19 & 20 & 21 & 22 & 23 & 22 & 23 & 24 & 25 & 26 & 27 \\
\end{tabularx}

\begin{paracol}{2}
\selectlanguage{latin}
\lettrine[lines=2]{S}{ancti} Joánnis Cántii, 
 Presbyteri et Confessóris, qui nono Kaléndas Januárii obdormívit in Dómino.
\switchcolumn
\selectlanguage{english}
\lettrine[lines=2]{S}{t.} John Cantius, priest and 
 confessor, who fell asleep in the Lord on the 24th of December.
\switchcolumn*
\selectlanguage{latin}
In Aviénsi civitáte, 
 prope Aquilam, in Vestínis, natális beáti Máximi, Levítæ et Mártyris; qui, 
 patiéndi desidério, inquiréntibus se persecutóribus palam osténdit, et, post 
 responsiónis constántiam, equúleo suspénsus ac tortus, deínde fústibus cæsus, 
 ad últimum, e sublími loco præcipitátus, occúbuit.
\switchcolumn
\selectlanguage{english}
At Abia, near Aquila in Abruzzo, the 
 birthday of blessed Maximus, deacon and martyr. Because of his desire 
 to suffer he shewed himself to the persecutors of his own accord. 
 After answering with great constancy, he was racked and tortured, then 
 beaten with rods, and he finally died by being cast headlong from a high 
 place.
\switchcolumn*
\selectlanguage{latin}
Agénni, in Gállia, 
 sancti Caprásii Mártyris, qui, cum rábiem persecutiónis declínans latéret in 
 spelúnca, tandem, áudiens quáliter beáta Fides Virgo pro Christo agonizáret, 
 índeque animátus ad tolerántiam passiónum, orávit ad Dóminum, ut, si eum 
 glória martyrii dignum judicáret, ex lápide spelúncæ limpidíssima aqua 
 manáret; quod cum Dóminum præstitísset, secúrus ad áream certáminis 
 properávit, et palmam martyrii, sub Maximiáno Imperatóre, fórtiter dimicándo 
 proméruit.
\switchcolumn
\selectlanguage{english}
At Agen in France, St. Caprasius, 
 martyr. He was hiding in a cavern to avoid the violence of the 
 persecution when the report of the blessed virgin Faith's courage in 
 suffering for Christ roused him to endure the torments. He prayed to 
 God that, if he were deemed worthy of the glory of martyrdom, clear water 
 might flow from the rock of his cave. God granted his prayer, and he 
 went with confidence to the scene of the trial, where, after a valiant 
 struggle, he merited the palm of martyrdom under Maximian.
\switchcolumn*
\selectlanguage{latin}
Antiochíæ sancti 
 Artémii, Ducis Augustális, qui, sub Constantíno Magno præcláris milítiæ 
 honóribus functus, a Juliáno Apóstata, quem sævítiæ in Christiános argúerat, 
 fústibus cædi, aliísque torméntis afflígi, ac demum cápite truncári jubétur.
\switchcolumn
\selectlanguage{english}
At Antioch, St. Artemius, an 
 imperial officer who had filled high positions in the army under Constantine 
 the Great. Julian the Apostate, however, whom he rebuked for his 
 cruelty towards Christians, ordered him to be beaten with rods, subjected to 
 other torments, and finally beheaded.
\switchcolumn*
\selectlanguage{latin}
Constantinópoli sancti 
 Andréæ Creténsis Mónachi, qui ob cultum sacrárum Imáginum, sub Constantíno 
 Coprónymo, sæpius verberátus, tandem, amputáto áltero pede, réddidit 
 spíritum.
\switchcolumn
\selectlanguage{english}
At Constantinople, St. Andrew of 
 Crete, a monk who had often been scourged by Constantine Copronymus for his 
 veneration of holy images. After one of his feet had been cut off he 
 rendered up his soul.
\switchcolumn*
\selectlanguage{latin}
Colóniæ Agrippínæ 
 pássio sanctárum Vírginum Marthæ et Saulæ, cum áliis plúribus.
\switchcolumn
\selectlanguage{english}
At Cologne, the martyrdom of the 
 holy virgins Martha and Saula, with many others.
\switchcolumn*
\selectlanguage{latin}
Apud Nabántiam, in 
 Lusitánia, sanctæ Irénes, Vírginis et Mártyris; cujus corpus honorífice 
 sepúltum fuit in óppido Scálabi, quod ipsíus Sanctæ nómine insign\'itum inde 
 permánsit.
\switchcolumn
\selectlanguage{english}
In Portugal, St. Irene, virgin and 
 martyr. Her body was honourably buried in the town of Scalabris. 
 Since that time the town has been named Santarem, which is derived from her 
 name.
\switchcolumn*
\selectlanguage{latin}
Alsóntiæ, in território 
 Rheménsi, sancti Sindúlphi, Presbyteri et Confessóris.
\switchcolumn
\selectlanguage{english}
At Aussonce, in the diocese of 
 Rheims, St. Sindulphus, priest and confessor.
\switchcolumn*
\selectlanguage{latin}
Apud Mindam, in 
 Germánia, Translátio sancti Feliciáni, Epíscopi Fulginátis et Mártyris; 
 cujus ibi depósita fuit sacrárum pars reliquiárum, quæ in Germániam advéctæ 
 sunt ex Umbriæ civitáte Fulgíneo, ubi ipse nono Kaléndas Februárii passus 
 quondam fúerat.
\switchcolumn
\selectlanguage{english}
At Minden in Germany, the 
 translation of St. Felician, bishop of Foligno and martyr. From his 
 holy relics a portion was placed in an urn and brought to Germany from the 
 city of Foligno in Umbria, where he had died on the 24th of January.
\switchcolumn*
\selectlanguage{latin}
Lutétiæ Parisiórum item 
 Translátio sanctórum Mártyrum Geórgii Diáconi, et Aurélii, ex urbe Hispániæ 
 Córduba, in qua olim, una cum áliis tribus Sóciis, ambo martyrium 
 compléverant sexto Kaléndas Augústi.
\switchcolumn
\selectlanguage{english}
At Paris, the translation of the 
 holy martyrs George, a deacon, and Aurelius from Cordova, a city of Spain, 
 where they had died with three companions on the 27th of July.
\switchcolumn*
\selectlanguage{latin}
\end{paracol}


% ---- martyrology/mart10/mart1021.htm
\needspace{10\baselineskip}
\begin{paracol}{2}
\selectlanguage{latin}
\begin{center}{\color{gregoriocolor} Duodécimo Kaléndas Novémbris. 
 Luna\dots\ }\end{center}
\switchcolumn
\selectlanguage{english}
\begin{center}{\color{gregoriocolor} The 
 21\textsuperscript{st} Day of 
 October. The\dots\ Day of the Moon.}\end{center}
\end{paracol}

\noindent\begin{tabularx}{\linewidth}{*{19}{>{\centering\arraybackslash}X}}
 \textcolor{gregoriocolor}{a} & \textcolor{gregoriocolor}{b} & \textcolor{gregoriocolor}{c} & \textcolor{gregoriocolor}{d} & \textcolor{gregoriocolor}{e} & \textcolor{gregoriocolor}{f} & \textcolor{gregoriocolor}{g} & \textcolor{gregoriocolor}{h} & \textcolor{gregoriocolor}{i} & \textcolor{gregoriocolor}{k} & \textcolor{gregoriocolor}{l} & \textcolor{gregoriocolor}{m} & \textcolor{gregoriocolor}{n} & \textcolor{gregoriocolor}{p} & \textcolor{gregoriocolor}{q} & \textcolor{gregoriocolor}{r} & \textcolor{gregoriocolor}{s} & \textcolor{gregoriocolor}{t} & \textcolor{gregoriocolor}{u} \\
 29 & 1 & 2 & 3 & 4 & 5 & 6 & 7 & 8 & 9 & 10 & 11 & 12 & 13 & 14 & 15 & 16 & 17 & 18 \\
\end{tabularx}
\vspace{0.5\baselineskip}
\noindent\begin{tabularx}{\linewidth}{*{12}{>{\centering\arraybackslash}X}}
 \textcolor{gregoriocolor}{A} & \textcolor{gregoriocolor}{B} & \textcolor{gregoriocolor}{C} & \textcolor{gregoriocolor}{D} & \textcolor{gregoriocolor}{E} & F & \textcolor{gregoriocolor}{F} & \textcolor{gregoriocolor}{G} & \textcolor{gregoriocolor}{H} & \textcolor{gregoriocolor}{M} & \textcolor{gregoriocolor}{N} & \textcolor{gregoriocolor}{P} \\
 19 & 20 & 21 & 22 & 23 & 24 & 23 & 24 & 25 & 26 & 27 & 28 \\
\end{tabularx}

\begin{paracol}{2}
\selectlanguage{latin}
\lettrine[lines=2]{I}{n} Cypro natális sancti 
 Hilariónis Abbátis, cujus vitam, virtútibus atque miráculis plenam, sanctus 
 Hierónymus scripsit.
\switchcolumn
\selectlanguage{english}
\lettrine[lines=2]{I}{n} Cyprus, the birthday of the holy 
 abbot Hilarion. His life, full of virtues and miracles, was written by 
 St. Jerome.
\switchcolumn*
\selectlanguage{latin}
Apud Colóniam 
 Aggripínam item natális sanctárum Ursulæ et Sociárum ejus; quæ, pro 
 Christiána religióne et virginitátis constántia ab Hunnis interféctæ, 
 martyrio vitam consummárunt, earúmque plúrima córpora fuérunt Colóniæ 
 cóndita.
\switchcolumn
\selectlanguage{english}
At Cologne, the birthday of St. 
 Ursula and her companions, who gained the martyr's crown by being slain by 
 the Huns for the Christian religion and their constancy in keeping their 
 virginity. Many of their bodies are buried in Cologne.
\switchcolumn*
\selectlanguage{latin}
Apud Ostia Tiberína 
 sancti Astérii, Presbyteri et Mártyris; qui (ut in passióne beáti Callísti 
 Papæ légitur) sub Alexándro Imperatóre passus est.
\switchcolumn
\selectlanguage{english}
At Ostia, St. Asterius, priest and 
 martyr, who suffered under Emperor Alexander, as we read in the Acts of 
 blessed Pope Callistus.
\switchcolumn*
\selectlanguage{latin}
Nicomedíæ natális 
 sanctórum Dásii, Zótici, Caji et aliórum duódecim mílitum; qui, post divérsa 
 torménta, in mare demérsi sunt.
\switchcolumn
\selectlanguage{english}
At Nicomedia, the birthday of Saints 
 Dasius, Zoticus, Caius, and twelve other soldiers, who, after suffering 
 various torments, were drowned in the sea.
\switchcolumn*
\selectlanguage{latin}
Lugdúni, in Gállia, 
 sancti Viatóris, qui éxstitit miníster beáti Justi, Lugdunénsis Epíscopi.
\switchcolumn
\selectlanguage{english}
At Lyons in France, St. Viator, 
 deacon of blessed Justus, bishop of that city.
\switchcolumn*
\selectlanguage{latin}
Maróniæ, prope 
 Antiochíam, in Syria, sancti Malchi Mónachi.
\switchcolumn
\selectlanguage{english}
At Maronia, near Antioch in Syria, 
 St. Malchus, a monk.
\switchcolumn*
\selectlanguage{latin}
In castro Laudunénsi 
 sanctæ Cilíniæ, matris beáti Remígii, Epíscopi Rheménsis.
\switchcolumn
\selectlanguage{english}
At Laon, St. Cilinia, mother of 
 blessed Remigius, bishop of Rheims.
\switchcolumn*
\selectlanguage{latin}
\end{paracol}


% ---- martyrology/mart10/mart1022.htm
\needspace{10\baselineskip}
\begin{paracol}{2}
\selectlanguage{latin}
\begin{center}{\color{gregoriocolor} Undécimo Kaléndas Novémbris. 
 Luna\dots\ }\end{center}
\switchcolumn
\selectlanguage{english}
\begin{center}{\color{gregoriocolor} The 
 22\textsuperscript{nd} Day of 
 October. The\dots\ Day of the Moon.}\end{center}
\end{paracol}

\noindent\begin{tabularx}{\linewidth}{*{19}{>{\centering\arraybackslash}X}}
 \textcolor{gregoriocolor}{a} & \textcolor{gregoriocolor}{b} & \textcolor{gregoriocolor}{c} & \textcolor{gregoriocolor}{d} & \textcolor{gregoriocolor}{e} & \textcolor{gregoriocolor}{f} & \textcolor{gregoriocolor}{g} & \textcolor{gregoriocolor}{h} & \textcolor{gregoriocolor}{i} & \textcolor{gregoriocolor}{k} & \textcolor{gregoriocolor}{l} & \textcolor{gregoriocolor}{m} & \textcolor{gregoriocolor}{n} & \textcolor{gregoriocolor}{p} & \textcolor{gregoriocolor}{q} & \textcolor{gregoriocolor}{r} & \textcolor{gregoriocolor}{s} & \textcolor{gregoriocolor}{t} & \textcolor{gregoriocolor}{u} \\
 1 & 2 & 3 & 4 & 5 & 6 & 7 & 8 & 9 & 10 & 11 & 12 & 13 & 14 & 15 & 16 & 17 & 18 & 19 \\
\end{tabularx}
\vspace{0.5\baselineskip}
\noindent\begin{tabularx}{\linewidth}{*{12}{>{\centering\arraybackslash}X}}
 \textcolor{gregoriocolor}{A} & \textcolor{gregoriocolor}{B} & \textcolor{gregoriocolor}{C} & \textcolor{gregoriocolor}{D} & \textcolor{gregoriocolor}{E} & F & \textcolor{gregoriocolor}{F} & \textcolor{gregoriocolor}{G} & \textcolor{gregoriocolor}{H} & \textcolor{gregoriocolor}{M} & \textcolor{gregoriocolor}{N} & \textcolor{gregoriocolor}{P} \\
 20 & 21 & 22 & 23 & 24 & 25 & 24 & 25 & 26 & 27 & 28 & 29 \\
\end{tabularx}

\begin{paracol}{2}
\selectlanguage{latin}
\lettrine[lines=2]{H}{ierosólymis} sanctæ 
 Maríæ Salóme, matris sanctórum Jacóbi et Joánnis Apostolórum, quæ in 
 Evangelio légitur circa Dómini sepultúram sollícita.
\switchcolumn
\selectlanguage{english}
\lettrine[lines=2]{A}{t} Jerusalem, St. Mary Salome, the 
 mother of the apostles James and John, who is referred to in the Gospel as 
 having cared for the burial of our Lord.
\switchcolumn*
\selectlanguage{latin}
Item Hierosólymis beáti 
 Marci Epíscopi, claríssimi et doctíssimi viri, qui primus ex Géntibus 
 Ecclésiam Hierosolymórum suscépit regéndam, ac, non multo post, sub Antoníno 
 Imperatóre, martyrii méruit palmam.
\switchcolumn
\selectlanguage{english}
At Jerusalem, blessed Bishop Mark, a 
 noble and learned man, who was the first Gentile to govern the Church of 
 Jerusalem. His brief episcopate was rewarded by the palm of martyrdom 
 under Emperor Antoninus.
\switchcolumn*
\selectlanguage{latin}
Hadrianópoli, in 
 Thrácia, natális sanctórum Mártyrum Philíppi Epíscopi, Sevéri Presbyteri, 
 Eusébii et Hermétis; qui, sub Juliáno Apóstata, post cárceres et flagélla, 
 incéndio cremáti sunt.
\switchcolumn
\selectlanguage{english}
At Adrianople in Thrace, the 
 birthday of the holy martyrs Philip, a bishop, Severus, a priest, Eusebius, 
 and Hermes. After being imprisoned and scourged, they were burned 
 alive in the time of Julian the Apostate.
\switchcolumn*
\selectlanguage{latin}
Item sanctórum Mártyrum 
 Alexándri Epíscopi, Heraclíi mílitis, et Sociórum.
\switchcolumn
\selectlanguage{english}
Also, the holy martyrs Alexander, a 
 bishop, Heraclius, a soldier, and their companions.
\switchcolumn*
\selectlanguage{latin}
Apud Firmum, in Picéno, natális sancti Philíppi, Epíscopi et Mártyris.
\switchcolumn
\selectlanguage{english}
At Fermo in Piceno, the birthday of 
 St. Philip, bishop and martyr.
\switchcolumn*
\selectlanguage{latin}
Apud Colóniam 
 Agrippínam sanctæ Córdulæ, quæ, cum esset una ex sodálibus sanctæ Ursulæ, 
 atque aliárum supplíciis et cædibus pertérrita se occultásset, postrídie, 
 ejus rei pænitens, se ultro patefécit Hunnis, et, novíssima ómnium, martyrii 
 corónam accépit.
\switchcolumn
\selectlanguage{english}
At Cologne, St. Cordula, who was one 
 of the companions of St. Ursula. Being terrified by the punishments 
 and slaughter of the others, she hid herself, but repenting her deed, on the 
 next day she declared herself to the Huns of her own accord, and thus was 
 the last of them all to receive the crown of martyrdom.
\switchcolumn*
\selectlanguage{latin}
Oscæ, in Hispánia, 
 sanctárum Vírginum Nunilónis et Alódiæ sorórum, quæ, a Saracénis ob fídei 
 confessiónem capitáli senténtia punítæ, martyrium consummárunt.
\switchcolumn
\selectlanguage{english}
At Huesca in Spain, the holy virgins 
 Nunilo and Alodia, sisters, who endured martyrdom by being condemned to 
 capital punishment by the Saracens for the confession of the faith.
\switchcolumn*
\selectlanguage{latin}
Hierápoli, in Phrygia, 
 sancti Abércii Epíscopi, qui sub Marco Antoníno Imperatóre cláruit.
\switchcolumn
\selectlanguage{english}
At Hieropolis in Phrygia, St. 
 Abercius, bishop, who flourished under Emperor Marcus Antoninus.
\switchcolumn*
\selectlanguage{latin}
Rotómagi sancti Melánii 
 Epíscopi, qui, a sancto Stéphano Papa ordinátus, illuc ad prædicándum 
 Evangélium missus est.
\switchcolumn
\selectlanguage{english}
At Rouen, St. Melanius, bishop, who 
 was ordained by Pope St. Stephen and sent there to preach the Gospel.
\switchcolumn*
\selectlanguage{latin}
In Túscia sancti Donáti 
 Scoti, Epíscopi Fæsuláni.
\switchcolumn
\selectlanguage{english}
In Tuscany, St. Donatus of Scotland, 
 bishop of Fiesole.
\switchcolumn*
\selectlanguage{latin}
Verónæ sancti Verecúndi, 
 Epíscopi et Confessóris.
\switchcolumn
\selectlanguage{english}
At Verona, St. Verecundius, bishop 
 and confessor.
\switchcolumn*
\selectlanguage{latin}
\end{paracol}


% ---- martyrology/mart10/mart1023.htm
\needspace{10\baselineskip}
\begin{paracol}{2}
\selectlanguage{latin}
\begin{center}{\color{gregoriocolor} Décimo Kaléndas Novémbris. 
 Luna\dots\ }\end{center}
\switchcolumn
\selectlanguage{english}
\begin{center}{\color{gregoriocolor} The 
 23\textsuperscript{rd} Day of 
 October. The\dots\ Day of the Moon.}\end{center}
\end{paracol}

\noindent\begin{tabularx}{\linewidth}{*{19}{>{\centering\arraybackslash}X}}
 \textcolor{gregoriocolor}{a} & \textcolor{gregoriocolor}{b} & \textcolor{gregoriocolor}{c} & \textcolor{gregoriocolor}{d} & \textcolor{gregoriocolor}{e} & \textcolor{gregoriocolor}{f} & \textcolor{gregoriocolor}{g} & \textcolor{gregoriocolor}{h} & \textcolor{gregoriocolor}{i} & \textcolor{gregoriocolor}{k} & \textcolor{gregoriocolor}{l} & \textcolor{gregoriocolor}{m} & \textcolor{gregoriocolor}{n} & \textcolor{gregoriocolor}{p} & \textcolor{gregoriocolor}{q} & \textcolor{gregoriocolor}{r} & \textcolor{gregoriocolor}{s} & \textcolor{gregoriocolor}{t} & \textcolor{gregoriocolor}{u} \\
 2 & 3 & 4 & 5 & 6 & 7 & 8 & 9 & 10 & 11 & 12 & 13 & 14 & 15 & 16 & 17 & 18 & 19 & 20 \\
\end{tabularx}
\vspace{0.5\baselineskip}
\noindent\begin{tabularx}{\linewidth}{*{12}{>{\centering\arraybackslash}X}}
 \textcolor{gregoriocolor}{A} & \textcolor{gregoriocolor}{B} & \textcolor{gregoriocolor}{C} & \textcolor{gregoriocolor}{D} & \textcolor{gregoriocolor}{E} & F & \textcolor{gregoriocolor}{F} & \textcolor{gregoriocolor}{G} & \textcolor{gregoriocolor}{H} & \textcolor{gregoriocolor}{M} & \textcolor{gregoriocolor}{N} & \textcolor{gregoriocolor}{P} \\
 21 & 22 & 23 & 24 & 25 & 26 & 25 & 26 & 27 & 28 & 29 & 1 \\
\end{tabularx}

\begin{paracol}{2}
\selectlanguage{latin}
\lettrine[lines=2]{A}{pud} Villáckum, in 
 Pannónia, natális sancti Joánnis de Capistráno, Sacerdótis ex Ordine Minórum 
 et Confessóris, vitæ sanctitáte ac fídei cathólicæ propagándæ zelo illústris; 
 qui Taurunénsem arcem, validíssimo Turcárum exércitu profligáto, suis 
 précibus et miráculis ab obsidióne liberávit. Ejus tamen festívitas 
 quinto Kaléndas Aprílis recólitur.
\switchcolumn
\selectlanguage{english}
\lettrine[lines=2]{A}{t} Vilak in Hungary, the birthday of 
 St. John Capistran, priest and confessor of the Order of Friars Minor, 
 illustrious for the sanctity of his life and his zeal for the propagation of 
 the Catholic faith. By his prayers and miracles, he routed a powerful 
 army of Turks, and forced them to quit the siege of Tornau. His 
 feastday, however, is celebrated on the 28th of March.
\switchcolumn*
\selectlanguage{latin}
Antiochíæ item natális 
 sancti Theodóri Presbyteri, qui, in persecutióne ímpii Juliáni comprehénsus, 
 et, post equúlei pœnam et multos ac duríssimos cruciátus, lampádibus étiam 
 circa látera appósitis adústus, tandem, cum in confessióne Christi 
 persísteret, gládii occisióne martyrium consummávit.
\switchcolumn
\selectlanguage{english}
At Antioch, the birthday of the holy 
 priest Theodore, who was arrested in the persecution of the impious Julian. 
 After the torment of the rack and many severe tortures, including the 
 burning of his sides with torches, he persisted in the confession of Christ, 
 and so his martyrdom was completed by death with the sword.
\switchcolumn*
\selectlanguage{latin}
Ad fundum Ursoniánum 
 prope Gades, in Hispánia, sanctórum Mártyrum Servándi et Germáni, qui, in 
 persecutióne Diocletiáni, sub Viatóre Vicário, post vérbera, squalórem cárceris, famis ac sitis injúriam, et longíssimi itíneris labórem, quem 
 pertulérunt ferro onústi, novíssime martyrii sui cursum, cæsis cervícibus, 
 implevérunt; ex quibus Germánus Eméritæ, Servándus autem Híspali cónditus 
 est.
\switchcolumn
\selectlanguage{english}
At Osuma, near Cadiz in Spain, in 
 the persecution of Diocletian, under the subgovernor Viator, the holy 
 martyrs Servandus and Germanus. They were subjected to scourging, 
 imprisonment in a foul dungeon, want of food and drink, and the fatigue of a 
 long journey while loaded with fetters, and at length reached the end of 
 their martyrdom by having their heads stricken off. Germanus was 
 buried at Merida, and Servandus at Seville.
\switchcolumn*
\selectlanguage{latin}
Constantinópoli sancti 
 Ignátii Epíscopi, qui, cum Bardam Cæsarem ob repudiátam uxórem arguísset, ab 
 eo multis injúriis afféctus est, et in exsílium pulsus; sed, a sancto 
 Nicoláo, Románo Pontífice, restitútus, tandem in pace quiévit.
\switchcolumn
\selectlanguage{english}
At Constantinople, St. Ignatius, 
 bishop, who rebuked Bardas Caesar for putting away his wife, for which he 
 was subjected to many insults and driven into banishment. He was, 
 however, restored to his See by the Roman Pontiff Nicholas, and there died 
 in peace.
\switchcolumn*
\selectlanguage{latin}
Burdígalæ sancti 
 Severíni, Epíscopi Coloniénsis et Confessóris.
\switchcolumn
\selectlanguage{english}
At Bordeaux, St. Severin, bishop of 
 Cologne and confessor.
\switchcolumn*
\selectlanguage{latin}
Rotómagi sancti Románi 
 Epíscopi.
\switchcolumn
\selectlanguage{english}
At Rouen, Bishop St. Romanus.
\switchcolumn*
\selectlanguage{latin}
Apud Salérnum sancti 
 Veri Epíscopi.
\switchcolumn
\selectlanguage{english}
At Salerno, Bishop St. Verus.
\switchcolumn*
\selectlanguage{latin}
In território Ambianénsi sancti Domítii Presbyteri.
\switchcolumn
\selectlanguage{english}
In the district of Amiens, St. 
 Domitius, a priest.
\switchcolumn*
\selectlanguage{latin}
In pago Pictaviénsi 
 sancti Benedícti Confessóris.
\switchcolumn
\selectlanguage{english}
In the country of Poitiers, St. 
 Benedict, confessor.
\switchcolumn*
\selectlanguage{latin}
Mántuæ beáti Joánnis 
 Boni, ex Eremitárum sancti Augustíni Ordine, Confessóris; cujus præcláram 
 vitam sanctus Antonínus conscrípsit.
\switchcolumn
\selectlanguage{english}
At Mantua, blessed John the Good, of 
 the Order of Hermits of St. Augustine, whose celebrated life was written by 
 St. Antoninus.
\switchcolumn*
\selectlanguage{latin}
\end{paracol}


% ---- martyrology/mart10/mart1024.htm
\needspace{10\baselineskip}
\begin{paracol}{2}
\selectlanguage{latin}
\begin{center}{\color{gregoriocolor} Nono Kaléndas Novémbris. 
 Luna\dots\ }\end{center}
\switchcolumn
\selectlanguage{english}
\begin{center}{\color{gregoriocolor} The 
 24\textsuperscript{th} Day of 
 October. The\dots\ Day of the Moon.}\end{center}
\end{paracol}

\noindent\begin{tabularx}{\linewidth}{*{19}{>{\centering\arraybackslash}X}}
 \textcolor{gregoriocolor}{a} & \textcolor{gregoriocolor}{b} & \textcolor{gregoriocolor}{c} & \textcolor{gregoriocolor}{d} & \textcolor{gregoriocolor}{e} & \textcolor{gregoriocolor}{f} & \textcolor{gregoriocolor}{g} & \textcolor{gregoriocolor}{h} & \textcolor{gregoriocolor}{i} & \textcolor{gregoriocolor}{k} & \textcolor{gregoriocolor}{l} & \textcolor{gregoriocolor}{m} & \textcolor{gregoriocolor}{n} & \textcolor{gregoriocolor}{p} & \textcolor{gregoriocolor}{q} & \textcolor{gregoriocolor}{r} & \textcolor{gregoriocolor}{s} & \textcolor{gregoriocolor}{t} & \textcolor{gregoriocolor}{u} \\
 3 & 4 & 5 & 6 & 7 & 8 & 9 & 10 & 11 & 12 & 13 & 14 & 15 & 16 & 17 & 18 & 19 & 20 & 21 \\
\end{tabularx}
\vspace{0.5\baselineskip}
\noindent\begin{tabularx}{\linewidth}{*{12}{>{\centering\arraybackslash}X}}
 \textcolor{gregoriocolor}{A} & \textcolor{gregoriocolor}{B} & \textcolor{gregoriocolor}{C} & \textcolor{gregoriocolor}{D} & \textcolor{gregoriocolor}{E} & F & \textcolor{gregoriocolor}{F} & \textcolor{gregoriocolor}{G} & \textcolor{gregoriocolor}{H} & \textcolor{gregoriocolor}{M} & \textcolor{gregoriocolor}{N} & \textcolor{gregoriocolor}{P} \\
 22 & 23 & 24 & 25 & 26 & 27 & 26 & 27 & 28 & 29 & 1 & 2 \\
\end{tabularx}

\begin{paracol}{2}
\selectlanguage{latin}
\lettrine[lines=2]{F}{estum} sancti Raphaélis 
 Archángeli, cujus dígnitas ac benefícia in sacro Tobíæ libro celebrántur.
\switchcolumn
\selectlanguage{english}
\lettrine[lines=2]{T}{he} Feast of St. Raphael the 
 Archangel, whose dignity and benefits to mankind are set forth in the holy 
 book of Tobias.
\switchcolumn*
\selectlanguage{latin}
Venúsiæ, in Apúlia, natális sanctórum Mártyrum Felícis, Epíscopi Africáni; Audácti et Januárii, 
 Presbyterórum; Fortunáti et Séptimi, Lectórum. Hi omnes, témpore 
 Diocletiáni, a Magdelliáno Procuratóre multis diu vínculis et carcéribus in 
 Africa et Sicília maceráti sunt, et, cum Felix sacros Libros juxta ipsíus 
 Imperatóris edíctum trádere nullátenus voluísset, gládii tandem occisióne 
 consummáti sunt.
\switchcolumn
\selectlanguage{english}
At Venosa in Apulia, the birthday of 
 the holy martyrs Felix, an African bishop, Audactus and Januarius, priests, 
 and the lectors Fortunatus and Septimus. In the time of Diocletian, 
 under the governor Magdellian, they were loaded with fetters and imprisoned 
 for a long time in Africa and Sicily. Because Felix refused to deliver 
 the sacred books, they were at last slain with the sword.
\switchcolumn*
\selectlanguage{latin}
Tungris, in Bélgio, 
 sancti Evergísli, Epíscopi Coloniénsis et Mártyris; qui, ob pastorális 
 offícii curam illuc proféctus, ibídem, dum nocte solus ad monastérium 
 sanctíssimæ Genitrícis Dei Maríæ oratúrus pérgeret, a latrónibus sagítta 
 percússus occúbuit.
\switchcolumn
\selectlanguage{english}
At Tongres in Belgium, St. 
 Evergislus, bishop of Cologne and martyr. Because of his duties in the 
 pastoral office he journeyed there, and on the way stopped to pray alone at 
 the monastery of the Blessed Virgin Mary where he was killed by robbers who 
 struck him with an arrow.
\switchcolumn*
\selectlanguage{latin}
In civitáte Nagran, 
 apud Homerítas, in Arábia, pássio sanctórum Arétæ, et Sociórum ejus 
 trecentórum et quadragínta, témpore Justíni Imperatóris, sub Dúnaan, Judæo 
 tyránno. Post eos Christiána múlier incéndio trádita est; cujus fílius 
 annórum quinque, cum Christum balbutiéndo confiterétur, nec blandítiis nec 
 minis retinéri posset, in ignem ubi mater ardébat, se præcípitem dedit.
\switchcolumn
\selectlanguage{english}
In the city of Nagran in Arabia 
 Felix, the passion of St. Aretas and his companions, to the number of three 
 hundred and forty, in the time of Emperor Justin, under the Jewish tyrant 
 Dunaan. After them, a Christian woman was burned alive, whose 
 five-year-old son confessed Christ in a lisping voice and could not be 
 prevented by caresses or threats from rushing into the fire in which his 
 mother was burning
\switchcolumn*
\selectlanguage{latin}
Constantinópoli sancti 
 Procli Epíscopi.
\switchcolumn
\selectlanguage{english}
At Constantinople, St. Proclus, 
 bishop.
\switchcolumn*
\selectlanguage{latin}
In ínsula Sargiénsi sancti Maglórii Epíscopi, qui ibídem, dimísso Episcopáli offício, quod erga 
 sparsos in Armórica Británnos per triénnium exercúerat, monastérium 
 constrúxit, in quo sancte réliquum vitæ tempus exégit; cujus corpus Lutétiam 
 Parisiórum póstea translátum fuit.
\switchcolumn
\selectlanguage{english}
On the island of Jersey, St. 
 Maglorius, bishop, who laid down the Episcopal office after exercising it 
 for three years towards a few scattered people in Brittany. He built a 
 monastery on that island, and there spent the remainder of his life in holy 
 conversation. His body was later translated to Paris.
\switchcolumn*
\selectlanguage{latin}
In monastério Montis 
 Frígidi, diœcésis Carcassonénsis, in Gállia, sancti Antónii Maríæ Claret, 
 olim Archiepíscopi Cubáni, Fundatóris Missionariórum Filiórum Immaculáti 
 Cordis beátæ Maríæ Vírginis, animárum zelo et mansuetúdine præclári; quem 
 Pius Duodécimus, Póntifex Máximus, Sanctórum fastis adscrípsit.
\switchcolumn
\selectlanguage{english}
In the monastery of Fontfroide in 
 the diocese of Carcassonne in France, St. Anthony Mary Claret, formerly 
 Archbishop of Cuba, and founder of the Missionary Sons of the Immaculate 
 Heart of the Blessed Virgin Mary. He was renowned for his meekness and 
 zeal for souls, and was canonized by the Supreme Pontiff, Pius XII.
\switchcolumn*
\selectlanguage{latin}
In monastério Duríni, 
 in Gállia, sancti Martíni, Diáconi et Abbátis, cujus corpus inde ad Vertávum 
 monastérium delátum est.
\switchcolumn
\selectlanguage{english}
In the monastery of Durin in France, 
 St. Martin, abbot and deacon. His body was translated to the monastery 
 of Vertou.
\switchcolumn*
\selectlanguage{latin}
In Campánia sancti 
 Marci Solitárii, cujus præclára ópera sanctus Gregórius Papa descrípsit.
\switchcolumn
\selectlanguage{english}
In Campania, St. Mark, a solitary, 
 whose noble accomplishments have been recorded by St. Gregory.
\switchcolumn*
\selectlanguage{latin}
\end{paracol}


% ---- martyrology/mart10/mart1025.htm
\needspace{10\baselineskip}
\begin{paracol}{2}
\selectlanguage{latin}
\begin{center}{\color{gregoriocolor} Octávo Kaléndas Novémbris. 
 Luna\dots\ }\end{center}
\switchcolumn
\selectlanguage{english}
\begin{center}{\color{gregoriocolor} The 
 25\textsuperscript{th} Day of 
 October. The\dots\ Day of the Moon.}\end{center}
\end{paracol}

\noindent\begin{tabularx}{\linewidth}{*{19}{>{\centering\arraybackslash}X}}
 \textcolor{gregoriocolor}{a} & \textcolor{gregoriocolor}{b} & \textcolor{gregoriocolor}{c} & \textcolor{gregoriocolor}{d} & \textcolor{gregoriocolor}{e} & \textcolor{gregoriocolor}{f} & \textcolor{gregoriocolor}{g} & \textcolor{gregoriocolor}{h} & \textcolor{gregoriocolor}{i} & \textcolor{gregoriocolor}{k} & \textcolor{gregoriocolor}{l} & \textcolor{gregoriocolor}{m} & \textcolor{gregoriocolor}{n} & \textcolor{gregoriocolor}{p} & \textcolor{gregoriocolor}{q} & \textcolor{gregoriocolor}{r} & \textcolor{gregoriocolor}{s} & \textcolor{gregoriocolor}{t} & \textcolor{gregoriocolor}{u} \\
 4 & 5 & 6 & 7 & 8 & 9 & 10 & 11 & 12 & 13 & 14 & 15 & 16 & 17 & 18 & 19 & 20 & 21 & 22 \\
\end{tabularx}
\vspace{0.5\baselineskip}
\noindent\begin{tabularx}{\linewidth}{*{12}{>{\centering\arraybackslash}X}}
 \textcolor{gregoriocolor}{A} & \textcolor{gregoriocolor}{B} & \textcolor{gregoriocolor}{C} & \textcolor{gregoriocolor}{D} & \textcolor{gregoriocolor}{E} & F & \textcolor{gregoriocolor}{F} & \textcolor{gregoriocolor}{G} & \textcolor{gregoriocolor}{H} & \textcolor{gregoriocolor}{M} & \textcolor{gregoriocolor}{N} & \textcolor{gregoriocolor}{P} \\
 23 & 24 & 25 & 26 & 27 & 28 & 27 & 28 & 29 & 1 & 2 & 3 \\
\end{tabularx}

\begin{paracol}{2}
\selectlanguage{latin}
\lettrine[lines=2]{R}{omæ} sanctórum Mártyrum 
 Chrysánthi et Daríæ uxóris, qui, post multas, quas sub Celeríno Præfécto pro 
 Christo sustinuérunt, passiónes, a Numeriáno Imperatóre jussi sunt via 
 Salária in Arenário depóni, atque vivéntes illic terra et lapídibus óbrui.
\switchcolumn
\selectlanguage{english}
\lettrine[lines=2]{A}{t} Rome, the holy martyrs 
 Chrysanthus and his wife Daria. After many sufferings endured for 
 Christ under the prefect Celerinus, they were ordered by Emperor Numerian to 
 be thrown into a sandpit on the Salarian Way, where, being still alive, were 
 covered with earth and stones.
\switchcolumn*
\selectlanguage{latin}
Ibídem natális sancti 
 Marcellíni, Papæ et Mártyris; qui, sub Maximiáno, pro fide Christi, una cum 
 Cláudio, Cyríno et Antoníno, cápite truncátus est. Quo témpore ita 
 magna fuit persecútio, ut decem et septem míllia Christianórum, intra unum 
 mensem, martyrio coronaréntur. Ipsíus tamen sancti Marcellíni festum, 
 una cum festo sancti Cleti, Papæ et Mártyris, sexto Kaléndas Maji celebrátur.
\switchcolumn
\selectlanguage{english}
Also, the birthday of St. 
 Marcellinus, pope and martyr, who was beheaded for the faith of Christ in 
 the reign of Maximian along with Claudius Cyrinus and Antoninus. So 
 great was the persecution then that seventeen thousand Christians received 
 the crown of martyrdom in the space of one month. The feast of St. 
 Marcellinus is celebrated with that of St. Cletus, pope and martyr, on the 
 26th of April.
\switchcolumn*
\selectlanguage{latin}
Petragóricis, in 
 Gállia, sancti Frontónis, qui, a beáto Petro Apóstolo Epíscopus ordinátus, 
 cum Geórgio Presbytero magnam illíus gentis multitúdinem convértit ad 
 Christum, et miráculis clarus, in pace quiévit.
\switchcolumn
\selectlanguage{english}
At Perigueux in France, St. Fronto, 
 who was made bishop by the blessed apostle Peter. Along with a priest 
 named George, he converted to Christ a large number of people of that place, 
 and, renowned for miracles, rested in peace.
\switchcolumn*
\selectlanguage{latin}
Romæ natális sanctórum 
 quadragínta sex mílitum, qui, simul baptizáti a sancto Dionysio Papa, mox 
 Cláudii Imperatóris jussu decolláti sunt, ac via Salária sepúlti; ubi et 
 álii Mártyres centum vigínti et unus pósiti sunt, inter quos fuérunt quátuor 
 mílites Christi, scílicet Theodósius, Lúcius, Marcus et Petrus.
\switchcolumn
\selectlanguage{english}
Also at Rome, the birthday of 
 forty-six holy soldiers, who were baptized at the same time by Pope Denis, 
 and soon after beheaded by order of Emperor Claudius. They were buried 
 on the Salarian Way with one hundred and twenty-one other martyrs. 
 Among them are named four soldiers of Christ: Theodosius, Lucius, Mark, and 
 Peter.
\switchcolumn*
\selectlanguage{latin}
Túrribus, in Sardínia, sanctórum Mártyrum Proti Presbyteri, et Januárii Diáconi, qui, a sancto Cajo 
 Papa ad eam ínsulam missi, ibídem, témpore Diocletiáni, sub Bárbaro Præside, 
 consummáti sunt.
\switchcolumn
\selectlanguage{english}
At Sassari in Sardinia, the holy 
 martyrs Protus, a priest, and Januarius, a deacon, who were sent to that 
 island Pope St. Caius, and were martyred in the time of Diocletian under the 
 governor Barbarus.
\switchcolumn*
\selectlanguage{latin}
Constantinópoli pássio 
 sanctórum Martyrii Subdiáconi, et Marciáni Cantóris, qui ab hæréticis, sub 
 Constántio Imperatóre, necáti sunt.
\switchcolumn
\selectlanguage{english}
At Constantinople, the martyrdom of 
 the Saints Martyrius, subdeacon, and Marcian, a cantor, who were slain by 
 the heretics during the reign of Emperor Constantius.
\switchcolumn*
\selectlanguage{latin}
Suessióne, in Gálliis, 
 sanctórum Mártyrum Crispíni et Crispiniáni, nobílium Romanórum, qui, in 
 persecutióne Diocletiáni, sub Rictiováro Præside, post immánia torménta 
 gládio trucidáti, corónam martyrii sunt consecúti; quorum córpora póstea 
 Romam deláta fuérunt, atque in Ecclésia sancti Lauréntii in Pane et Perna 
 honorífice tumuláta.
\switchcolumn
\selectlanguage{english}
At Soissons in France, in the 
 persecution of Diocletian, the holy martyrs Crispin and Crispinian, noble 
 Romans. Under Governor Rictiovarus, after horrible torments, they were 
 put to the sword, and thus obtained the crown of martyrdom. Their 
 bodies were afterwards conveyed to Rome and entombed with due honours in the 
 church of St. Lawrence in Panisperna.
\switchcolumn*
\selectlanguage{latin}
Floréntiæ pássio beáti 
 Miniátis mílitis, qui, sub Décio Príncipe, pro fide Christi egrégie certans, 
 nóbili martyrio coronátur.
\switchcolumn
\selectlanguage{english}
At Florence, St. Minias, a soldier, 
 who fought valorously for the faith of Christ and was gloriously crowned 
 with martyrdom during the reign of Decius.
\switchcolumn*
\selectlanguage{latin}
Bríxiæ natális sancti 
 Gaudéntii Epíscopi, eruditióne et sanctitáte conspícui.
\switchcolumn
\selectlanguage{english}
At Brescia, the birthday of St. 
 Gaudentius, bishop, distinguished for his learning and holiness.
\switchcolumn*
\selectlanguage{latin}
Gavális, in Gállia, 
 sancti Hilárii Epíscopi.
\switchcolumn
\selectlanguage{english}
At Javoux in France, St. Hilary, 
 bishop.
\switchcolumn*
\selectlanguage{latin}
\end{paracol}


% ---- martyrology/mart10/mart1026.htm
\needspace{10\baselineskip}
\begin{paracol}{2}
\selectlanguage{latin}
\begin{center}{\color{gregoriocolor} Séptimo Kaléndas Novémbris. 
 Luna\dots\ }\end{center}
\switchcolumn
\selectlanguage{english}
\begin{center}{\color{gregoriocolor} The 
 26\textsuperscript{th} Day of 
 October. The\dots\ Day of the Moon.}\end{center}
\end{paracol}

\noindent\begin{tabularx}{\linewidth}{*{19}{>{\centering\arraybackslash}X}}
 \textcolor{gregoriocolor}{a} & \textcolor{gregoriocolor}{b} & \textcolor{gregoriocolor}{c} & \textcolor{gregoriocolor}{d} & \textcolor{gregoriocolor}{e} & \textcolor{gregoriocolor}{f} & \textcolor{gregoriocolor}{g} & \textcolor{gregoriocolor}{h} & \textcolor{gregoriocolor}{i} & \textcolor{gregoriocolor}{k} & \textcolor{gregoriocolor}{l} & \textcolor{gregoriocolor}{m} & \textcolor{gregoriocolor}{n} & \textcolor{gregoriocolor}{p} & \textcolor{gregoriocolor}{q} & \textcolor{gregoriocolor}{r} & \textcolor{gregoriocolor}{s} & \textcolor{gregoriocolor}{t} & \textcolor{gregoriocolor}{u} \\
 5 & 6 & 7 & 8 & 9 & 10 & 11 & 12 & 13 & 14 & 15 & 16 & 17 & 18 & 19 & 20 & 21 & 22 & 23 \\
\end{tabularx}
\vspace{0.5\baselineskip}
\noindent\begin{tabularx}{\linewidth}{*{12}{>{\centering\arraybackslash}X}}
 \textcolor{gregoriocolor}{A} & \textcolor{gregoriocolor}{B} & \textcolor{gregoriocolor}{C} & \textcolor{gregoriocolor}{D} & \textcolor{gregoriocolor}{E} & F & \textcolor{gregoriocolor}{F} & \textcolor{gregoriocolor}{G} & \textcolor{gregoriocolor}{H} & \textcolor{gregoriocolor}{M} & \textcolor{gregoriocolor}{N} & \textcolor{gregoriocolor}{P} \\
 24 & 25 & 26 & 27 & 28 & 29 & 28 & 29 & 1 & 2 & 3 & 4 \\
\end{tabularx}

\begin{paracol}{2}
\selectlanguage{latin}
\lettrine[lines=2]{R}{omæ} sancti Evarísti, 
 Papæ et Mártyris, qui Dei Ecclésiam, sub Hadriáno Imperatóre, suo sánguine 
 purpurávit.
\switchcolumn
\selectlanguage{english}
\lettrine[lines=2]{A}{t} Rome, St. Evaristus, pope and 
 martyr, who enriched the Church of God with his blood under Emperor Hadrian.
\switchcolumn*
\selectlanguage{latin}
In Africa sanctórum 
 Mártyrum Rogatiáni Presbyteri, et Felicíssimi, qui, in persecutióne 
 Valeriáni et Galliéni, illústri martyrio coronáti sunt; de quibus étiam 
 scribit sanctus Cypriánus in epístola ad Confessóres.
\switchcolumn
\selectlanguage{english}
In Africa, the holy martyrs 
 Felicissimus and the priest Rogatian, who received the bright crown of 
 martyrs in the persecution of Valerian and Gallienus. They are 
 mentioned by St. Cyprian in his Epistle to the Confessors.
\switchcolumn*
\selectlanguage{latin}
Nicomedíæ sanctórum 
 Mártyrum Luciáni, Flórii et Sociórum.
\switchcolumn
\selectlanguage{english}
At Nicomedia, the holy martyrs 
 Lucian, Florius, and their companions.
\switchcolumn*
\selectlanguage{latin}
Narbóne, in Gállia, 
 sancti Rústici, Epíscopi et Confessóris; qui cláruit tempóribus Valentiniáni 
 et Leónis Imperatórum.
\switchcolumn
\selectlanguage{english}
At Narbonne, St. Rusticus, bishop 
 and confessor, who flourished in the reigns of Emperors Leo and Valentian.
\switchcolumn*
\selectlanguage{latin}
Apud Salérnum sancti 
 Gaudiósi Epíscopi.
\switchcolumn
\selectlanguage{english}
At Salerno, St. Gaudiosus, bishop.
\switchcolumn*
\selectlanguage{latin}
Papíæ sancti Fulci 
 Epíscopi.
\switchcolumn
\selectlanguage{english}
At Pavia, Bishop St. Fulk.
\switchcolumn*
\selectlanguage{latin}
Item sancti 
 Quadragésimi Subdiáconi, qui et mórtuum resuscitávit.
\switchcolumn
\selectlanguage{english}
Also St. Quadragesimus, subdeacon, 
 who raised a dead man to life.
\switchcolumn*
\selectlanguage{latin}
\end{paracol}


% ---- martyrology/mart10/mart1027.htm
\needspace{10\baselineskip}
\begin{paracol}{2}
\selectlanguage{latin}
\begin{center}{\color{gregoriocolor} Sexto Kaléndas Novémbris. 
 Luna\dots\ }\end{center}
\switchcolumn
\selectlanguage{english}
\begin{center}{\color{gregoriocolor} The 
 27\textsuperscript{th} Day of 
 October. The\dots\ Day of the Moon.}\end{center}
\end{paracol}

\noindent\begin{tabularx}{\linewidth}{*{19}{>{\centering\arraybackslash}X}}
 \textcolor{gregoriocolor}{a} & \textcolor{gregoriocolor}{b} & \textcolor{gregoriocolor}{c} & \textcolor{gregoriocolor}{d} & \textcolor{gregoriocolor}{e} & \textcolor{gregoriocolor}{f} & \textcolor{gregoriocolor}{g} & \textcolor{gregoriocolor}{h} & \textcolor{gregoriocolor}{i} & \textcolor{gregoriocolor}{k} & \textcolor{gregoriocolor}{l} & \textcolor{gregoriocolor}{m} & \textcolor{gregoriocolor}{n} & \textcolor{gregoriocolor}{p} & \textcolor{gregoriocolor}{q} & \textcolor{gregoriocolor}{r} & \textcolor{gregoriocolor}{s} & \textcolor{gregoriocolor}{t} & \textcolor{gregoriocolor}{u} \\
 6 & 7 & 8 & 9 & 10 & 11 & 12 & 13 & 14 & 15 & 16 & 17 & 18 & 19 & 20 & 21 & 22 & 23 & 24 \\
\end{tabularx}
\vspace{0.5\baselineskip}
\noindent\begin{tabularx}{\linewidth}{*{12}{>{\centering\arraybackslash}X}}
 \textcolor{gregoriocolor}{A} & \textcolor{gregoriocolor}{B} & \textcolor{gregoriocolor}{C} & \textcolor{gregoriocolor}{D} & \textcolor{gregoriocolor}{E} & F & \textcolor{gregoriocolor}{F} & \textcolor{gregoriocolor}{G} & \textcolor{gregoriocolor}{H} & \textcolor{gregoriocolor}{M} & \textcolor{gregoriocolor}{N} & \textcolor{gregoriocolor}{P} \\
 25 & 26 & 27 & 28 & 29 & 30 & 29 & 1 & 2 & 3 & 4 & 5 \\
\end{tabularx}

\begin{paracol}{2}
\selectlanguage{latin}
\lettrine[lines=1]{V}{igília} sanctórum Apostolórum Simónis et Judæ.
\switchcolumn
\selectlanguage{english}
\lettrine[lines=1]{T}{he} vigil of the holy apostles Simon and Jude.
\switchcolumn*
\selectlanguage{latin}
Abulæ, in Hispánia, 
 pássio sanctórum Vincéntii, Sabínæ et Christétæ. Hi primum in equúleo 
 ádeo sunt exténti, ut omnes membrórum compáges laxaréntur; deínde cápita 
 eórum, lapídibus superpósita, usque ad cérebri excussiónem válidis véctibus 
 sunt contúsa, atque ita ipsi martyrium complevérunt, sub Præside Daciáno.
\switchcolumn
\selectlanguage{english}
At Avila in Spain, under the 
 governor Dacian, the Saints Vincent, Sabina, and Christeta. They were 
 first stretched on the rack in such a manner that all their limbs were 
 dislocated; then stones being laid on their heads, and their brains beaten 
 out with heavy bars, their martyrdom was fulfilled.
\switchcolumn*
\selectlanguage{latin}
Apud castrum Tyle, in 
 Gállia, sancti Floréntii Mártyris.
\switchcolumn
\selectlanguage{english}
At Tilchatel in France, St. 
 Florentius, martyr.
\switchcolumn*
\selectlanguage{latin}
In Cappadócia sanctárum 
 Mártyrum Capitolínæ, ejúsque ancíllæ Erothéidis, quæ sub Diocletiáno sunt 
 passæ.
\switchcolumn
\selectlanguage{english}
In Cappadocia, the holy martyrs 
 Capitolina, and Erotheides, her handmaid, who suffered under Diocletian.
\switchcolumn*
\selectlanguage{latin}
Apud Indos sancti 
 Fruménti Epíscopi, qui, primum ibi captívus, deínde, Epíscopus a sancto 
 Athanásio ordinátus, Evangélium in ea província propagávit.
\switchcolumn
\selectlanguage{english}
In India, St. Frumentius, bishop. 
 While he was a captive there he was consecrated bishop by St. Athanasius, 
 and propagated the Gospel in that country.
\switchcolumn*
\selectlanguage{latin}
Neápoli, in Campánia, sancti Gaudiósi, Epíscopi Africáni, qui ob Wandalórum persecutiónem venit in 
 Campániam, et, in monastério apud eam urbem, sancto fine quiévit.
\switchcolumn
\selectlanguage{english}
At Naples, St. Gaudiosus, an African 
 bishop who came to Campania because of the Vandal persecution, and died a 
 holy death in a monastery in that city.
\switchcolumn*
\selectlanguage{latin}
In Æthiópia sancti 
 Elésbaan Regis, qui, Christi hóstibus expugnátis, ac, témpore Justíni 
 Imperatóris, misso régio diadémate Hierosólymam, monásticam vitam, ut 
 vóverat agens, migrávit ad Dóminum.
\switchcolumn
\selectlanguage{english}
In Ethiopia, in the time of Emperor 
 Justin, St. Elesbaan, king. After having defeated the enemies of 
 Christ and sent his royal diadem to Jerusalem, he led a monastic life, as he 
 had vowed, and went to his reward.
\switchcolumn*
\selectlanguage{latin}
\end{paracol}


% ---- martyrology/mart10/mart1028.htm
\needspace{10\baselineskip}
\begin{paracol}{2}
\selectlanguage{latin}
\begin{center}{\color{gregoriocolor} Quinto Kaléndas Novémbris. 
 Luna\dots\ }\end{center}
\switchcolumn
\selectlanguage{english}
\begin{center}{\color{gregoriocolor} The 
 28\textsuperscript{th} Day of 
 October. The\dots\ Day of the Moon.}\end{center}
\end{paracol}

\noindent\begin{tabularx}{\linewidth}{*{19}{>{\centering\arraybackslash}X}}
 \textcolor{gregoriocolor}{a} & \textcolor{gregoriocolor}{b} & \textcolor{gregoriocolor}{c} & \textcolor{gregoriocolor}{d} & \textcolor{gregoriocolor}{e} & \textcolor{gregoriocolor}{f} & \textcolor{gregoriocolor}{g} & \textcolor{gregoriocolor}{h} & \textcolor{gregoriocolor}{i} & \textcolor{gregoriocolor}{k} & \textcolor{gregoriocolor}{l} & \textcolor{gregoriocolor}{m} & \textcolor{gregoriocolor}{n} & \textcolor{gregoriocolor}{p} & \textcolor{gregoriocolor}{q} & \textcolor{gregoriocolor}{r} & \textcolor{gregoriocolor}{s} & \textcolor{gregoriocolor}{t} & \textcolor{gregoriocolor}{u} \\
 7 & 8 & 9 & 10 & 11 & 12 & 13 & 14 & 15 & 16 & 17 & 18 & 19 & 20 & 21 & 22 & 23 & 24 & 25 \\
\end{tabularx}
\vspace{0.5\baselineskip}
\noindent\begin{tabularx}{\linewidth}{*{12}{>{\centering\arraybackslash}X}}
 \textcolor{gregoriocolor}{A} & \textcolor{gregoriocolor}{B} & \textcolor{gregoriocolor}{C} & \textcolor{gregoriocolor}{D} & \textcolor{gregoriocolor}{E} & F & \textcolor{gregoriocolor}{F} & \textcolor{gregoriocolor}{G} & \textcolor{gregoriocolor}{H} & \textcolor{gregoriocolor}{M} & \textcolor{gregoriocolor}{N} & \textcolor{gregoriocolor}{P} \\
 26 & 27 & 28 & 29 & 30 & 1 & 1 & 2 & 3 & 4 & 5 & 6 \\
\end{tabularx}

\begin{paracol}{2}
\selectlanguage{latin}
\lettrine[lines=2]{I}{n} Pérside natális 
 beatórum Apostolórum Simónis Chananæi, et Thaddǽi, qui et Judas dícitur. 
 Ex ipsis autem Simon in Ægypto, Thaddǽus in Mesopotámia Evangélium 
 prædicávit; deínde, in Pérsidem simul ingréssi, ibi, cum innúmeram gentis illíus multitúdinem Christo subdidíssent, martyrium consummárunt.
\switchcolumn
\selectlanguage{english}
\lettrine[lines=2]{I}{n} Persia, the birthday of the 
 blessed apostles Simon the Canaanite and Thaddeus, who is also called Jude. 
 Simon preached the Gospel in Egypt, Thaddeus in Mesopotamia. 
 Afterwards, entering Persia together, they converted to Christ a numberless 
 multitude of the inhabitants, then underwent martyrdom.
\switchcolumn*
\selectlanguage{latin}
Romæ sanctórum Mártyrum 
 Anastásiæ senióris Vírginis, et Cyrílli. Ipsa Virgo, in persecutióne 
 Valeriáni, sub Probo Præfécto, vínculis constrícta, cólaphis cæsa, igne et 
 verbéribus est cruciáta, et, cum in confessióne Christi permanéret immóbilis, 
 tandem, abscíssis mamíllis, evúlsis únguibus, déntibus comminútis, mánibus 
 pedibúsque præcísis, truncáta cápite, tot passiónum ornáta monílibus 
 migrávit ad Sponsum; Cyríllus autem, ei peténti aquam propínans, martyrium 
 pro mercéde accépit.
\switchcolumn
\selectlanguage{english}
At Rome, the holy martyrs Cyril and 
 Anastasia the Elder, virgin. In the persecution of Valerian, under the 
 prefect Probus, Anastasia was bound with chains, buffeted, subjected to fire 
 and scourging, and, as she remained immovable in the confession of Christ, 
 her breasts were cut away, her nails plucked out, her teeth broken, and her 
 hands, feet, and head severed from her body. Adorned with her 
 sufferings as with so many jewels, she went to her Spouse. At her 
 request, Cyril gave her some water to drink, and for his reward became a 
 martyr.
\switchcolumn*
\selectlanguage{latin}
Item Romæ sanctæ 
 Cyríllæ Vírginis, quæ fília éxstitit sanctæ Tryphóniæ, et, sub Cláudio 
 Príncipe, pro Christo juguláta est.
\switchcolumn
\selectlanguage{english}
In the same city, during the reign 
 of Claudius, St. Cyrilla, virgin, daughter of St. Tryphonia, who was pierced 
 through the throat for the faith of Christ.
\switchcolumn*
\selectlanguage{latin}
Apud Comum sancti 
 Fidélis Mártyris, sub Maximiáno Imperatóre.
\switchcolumn
\selectlanguage{english}
At Como, under Emperor Maximian, St. 
 Fidelis, martyr.
\switchcolumn*
\selectlanguage{latin}
Mogúntiæ sancti 
 Ferrútii Mártyris.
\switchcolumn
\selectlanguage{english}
At Mainz, St. Ferrutius, martyr.
\switchcolumn*
\selectlanguage{latin}
Meldis, in Gállia, 
 sancti Farónis, Epíscopi et Confessóris.
\switchcolumn
\selectlanguage{english}
At Meaux, in France, St. Faro, 
 bishop and confessor.
\switchcolumn*
\selectlanguage{latin}
Vercéllis sancti 
 Honoráti Epíscopi.
\switchcolumn
\selectlanguage{english}
At Vercelli, St. Honoratus, bishop.
\switchcolumn*
\selectlanguage{latin}
\end{paracol}


% ---- martyrology/mart10/mart1029.htm
\needspace{10\baselineskip}
\begin{paracol}{2}
\selectlanguage{latin}
\begin{center}{\color{gregoriocolor} Quarto Kaléndas Novémbris. 
 Luna\dots\ }\end{center}
\switchcolumn
\selectlanguage{english}
\begin{center}{\color{gregoriocolor} The 
 29\textsuperscript{th} Day of 
 October. The\dots\ Day of the Moon.}\end{center}
\end{paracol}

\noindent\begin{tabularx}{\linewidth}{*{19}{>{\centering\arraybackslash}X}}
 \textcolor{gregoriocolor}{a} & \textcolor{gregoriocolor}{b} & \textcolor{gregoriocolor}{c} & \textcolor{gregoriocolor}{d} & \textcolor{gregoriocolor}{e} & \textcolor{gregoriocolor}{f} & \textcolor{gregoriocolor}{g} & \textcolor{gregoriocolor}{h} & \textcolor{gregoriocolor}{i} & \textcolor{gregoriocolor}{k} & \textcolor{gregoriocolor}{l} & \textcolor{gregoriocolor}{m} & \textcolor{gregoriocolor}{n} & \textcolor{gregoriocolor}{p} & \textcolor{gregoriocolor}{q} & \textcolor{gregoriocolor}{r} & \textcolor{gregoriocolor}{s} & \textcolor{gregoriocolor}{t} & \textcolor{gregoriocolor}{u} \\
 8 & 9 & 10 & 11 & 12 & 13 & 14 & 15 & 16 & 17 & 18 & 19 & 20 & 21 & 22 & 23 & 24 & 25 & 26 \\
\end{tabularx}
\vspace{0.5\baselineskip}
\noindent\begin{tabularx}{\linewidth}{*{12}{>{\centering\arraybackslash}X}}
 \textcolor{gregoriocolor}{A} & \textcolor{gregoriocolor}{B} & \textcolor{gregoriocolor}{C} & \textcolor{gregoriocolor}{D} & \textcolor{gregoriocolor}{E} & F & \textcolor{gregoriocolor}{F} & \textcolor{gregoriocolor}{G} & \textcolor{gregoriocolor}{H} & \textcolor{gregoriocolor}{M} & \textcolor{gregoriocolor}{N} & \textcolor{gregoriocolor}{P} \\
 27 & 28 & 29 & 30 & 1 & 2 & 2 & 3 & 4 & 5 & 6 & 7 \\
\end{tabularx}

\begin{paracol}{2}
\selectlanguage{latin}
\lettrine[lines=2]{S}{anctórum} Episcopórum 
 Maximiliáni Mártyris, et Valentíni Confessóris.
\switchcolumn
\selectlanguage{english}
\lettrine[lines=2]{T}{he} holy bishop Maximian, martyr, 
 and Valentine, confessor.
\switchcolumn*
\selectlanguage{latin}
Sidóne, in Phœnícia, 
 sancti Zenóbii Presbyteri, qui, sub novíssimæ persecutiónis acerbitáte, ad 
 martyrium álios exhórtans, martyrio et ipse dignátus est.
\switchcolumn
\selectlanguage{english}
At Sidon in Phoenicia, St. Zenobius, 
 a priest. When the last persecution was raging, by exhorting others to 
 martyrdom, he himself was deemed worthy of it.
\switchcolumn*
\selectlanguage{latin}
In Lucánia sanctórum 
 Mártyrum Hyacínthi, Quincti, Feliciáni, et Lúcii.
\switchcolumn
\selectlanguage{english}
In Lucania, the holy martyrs 
 Hyacinth, Quinctus, Felician, and Lucius.
\switchcolumn*
\selectlanguage{latin}
Bérgomi sanctæ Eusébiæ, 
 Vírginis et Mártyris.
\switchcolumn
\selectlanguage{english}
At Bergamo, St. Eusebia, virgin and 
 martyr.
\switchcolumn*
\selectlanguage{latin}
Hierosólymis natális 
 beáti Narcíssi Epíscopi, sanctitáte, patiéntia ac fide laudábilis, qui, 
 centum et séxdecim annórum senex, felíciter migrávit ad Dóminum.
\switchcolumn
\selectlanguage{english}
At Jerusalem, the birthday of 
 blessed Narcissus, a bishop distinguished for holiness, patience, and faith, 
 who went to the kingdom of God at the age of one hundred and sixteen years.
\switchcolumn*
\selectlanguage{latin}
Augustodúni sancti 
 Joánnis, Epíscopi et Confessóris.
\switchcolumn
\selectlanguage{english}
At Autun, St. John, bishop and 
 confessor.
\switchcolumn*
\selectlanguage{latin}
Cassíope, in ínsula 
 Corcyra, sancti Donáti Epíscopi, de quo scribit beátus Gregórius Papa.
\switchcolumn
\selectlanguage{english}
At Cassiope, in the island of Corfu, 
 Bishop St. Donatus, mentioned by blessed Pope Gregory.
\switchcolumn*
\selectlanguage{latin}
Viénnæ, in Gállia, 
 deposítio beáti Theodóri Abbátis.
\switchcolumn
\selectlanguage{english}
At Vienne in France, the death of 
 blessed Theodore, abbot.
\switchcolumn*
\selectlanguage{latin}
\end{paracol}


% ---- martyrology/mart10/mart1030.htm
\needspace{10\baselineskip}
\begin{paracol}{2}
\selectlanguage{latin}
\begin{center}{\color{gregoriocolor} Tértio Kaléndas Novémbris. 
 Luna\dots\ }\end{center}
\switchcolumn
\selectlanguage{english}
\begin{center}{\color{gregoriocolor} The 
 30\textsuperscript{th} Day of 
 October. The\dots\ Day of the Moon.}\end{center}
\end{paracol}

\noindent\begin{tabularx}{\linewidth}{*{19}{>{\centering\arraybackslash}X}}
 \textcolor{gregoriocolor}{a} & \textcolor{gregoriocolor}{b} & \textcolor{gregoriocolor}{c} & \textcolor{gregoriocolor}{d} & \textcolor{gregoriocolor}{e} & \textcolor{gregoriocolor}{f} & \textcolor{gregoriocolor}{g} & \textcolor{gregoriocolor}{h} & \textcolor{gregoriocolor}{i} & \textcolor{gregoriocolor}{k} & \textcolor{gregoriocolor}{l} & \textcolor{gregoriocolor}{m} & \textcolor{gregoriocolor}{n} & \textcolor{gregoriocolor}{p} & \textcolor{gregoriocolor}{q} & \textcolor{gregoriocolor}{r} & \textcolor{gregoriocolor}{s} & \textcolor{gregoriocolor}{t} & \textcolor{gregoriocolor}{u} \\
 9 & 10 & 11 & 12 & 13 & 14 & 15 & 16 & 17 & 18 & 19 & 20 & 21 & 22 & 23 & 24 & 25 & 26 & 27 \\
\end{tabularx}
\vspace{0.5\baselineskip}
\noindent\begin{tabularx}{\linewidth}{*{12}{>{\centering\arraybackslash}X}}
 \textcolor{gregoriocolor}{A} & \textcolor{gregoriocolor}{B} & \textcolor{gregoriocolor}{C} & \textcolor{gregoriocolor}{D} & \textcolor{gregoriocolor}{E} & F & \textcolor{gregoriocolor}{F} & \textcolor{gregoriocolor}{G} & \textcolor{gregoriocolor}{H} & \textcolor{gregoriocolor}{M} & \textcolor{gregoriocolor}{N} & \textcolor{gregoriocolor}{P} \\
 28 & 29 & 30 & 1 & 2 & 3 & 3 & 4 & 5 & 6 & 7 & 8 \\
\end{tabularx}

\begin{paracol}{2}
\selectlanguage{latin}
\lettrine[lines=2]{I}{n} Sardínia natális 
 sancti Pontiáni, Papæ et Mártyris, qui, ab Alexándro Imperatóre, una cum 
 Hippólyto Presbytero, in eam ínsulam deportátus, ibídem, mactátus fústibus, 
 martyrium consummávit. Ejus corpus a beáto Fabiáno Papa Romam delátum 
 est, atque in cœmetério Callísti sepúltum. Ipsíus tamen festum 
 recólitur tertiodécimo Kaléndas Decémbris.
\switchcolumn
\selectlanguage{english}
\lettrine[lines=2]{I}{n} Sardinia, the birthday of St. 
 Pontianus, pope and martyr. In the company of the priest Hippolytus, he 
 was exiled by Emperor Alexander, and achieved martyrdom by being scourged. 
 His body was brought to Rome by blessed Pope Fabian and buried in the 
 cemetery of Callistus. His feast, however, is celebrated on the 19th 
 of November.
\switchcolumn*
\selectlanguage{latin}
Ægéæ, in Cilícia, pássio sanctórum Zenóbii Epíscopi, et Zenóbiæ soróris, sub Diocletiáno 
 Imperatóre et Lysia Præside.
\switchcolumn
\selectlanguage{english}
At Aegea in Cilicia, in the reign of 
 Diocletian, under the governor Lysias, the martyrdom of Saints Zenobius, 
 bishop, and his sister Zenobia.
\switchcolumn*
\selectlanguage{latin}
Altíni, in Venetórum 
 fínibus, sancti Theonésti, Epíscopi et Mártyris, qui ab Ariánis occísus est.
\switchcolumn
\selectlanguage{english}
At Altino, in the neighbourhood of 
 Venice, St. Theonestus, bishop and martyr, who was slain by the Arians.
\switchcolumn*
\selectlanguage{latin}
In Africa natális 
 sanctórum Mártyrum ducentórum vigínti.
\switchcolumn
\selectlanguage{english}
In Africa, the birthday of two 
 hundred and twenty holy martyrs.
\switchcolumn*
\selectlanguage{latin}
Tingi, in Mauritánia, pássio sancti Marcélli Centuriónis, qui, sanctórum Cláudii ac Lupérci et 
 Victórii Mártyrum pater, cápitis abscissióne martyrium complévit sub 
 Agricoláo, agénte vices Præfécti prætório.
\switchcolumn
\selectlanguage{english}
At Tangier in Morocco, St. 
 Marcellus, a centurion, the father of Saints Claudius, Lupercus, and 
 Victorius. He achieved martyrdom by beheading under Agricola, deputy 
 praetor for Praefectus.
\switchcolumn*
\selectlanguage{latin}
Alexandríæ sanctórum 
 trédecim Mártyrum, qui, cum sanctis Juliáno, Euno et Macário, passi sunt sub 
 Décio Imperatóre.
\switchcolumn
\selectlanguage{english}
At Alexandria, in the reign of 
 Decius, thirteen holy martyrs who suffered with Saints Julian, Eunus, and 
 Macarius.
\switchcolumn*
\selectlanguage{latin}
Cárali, in Sardínia, sancti Saturníni Mártyris, qui, in persecutióne Diocletiáni, sub Bárbaro 
 Præside, cápite truncátus est.
\switchcolumn
\selectlanguage{english}
At Cagliari in Sardinia, St. 
 Saturninus, martyr, who was beheaded under the governor Barbarus, during the 
 persecution of Diocletian.
\switchcolumn*
\selectlanguage{latin}
Apaméæ, in Phrygia, 
 sancti Máximi Mártyris, sub eódem Diocletiáno
\switchcolumn
\selectlanguage{english}
At Apamea in Phrygia, St. Maximus, 
 martyr, under the same Diocletian.
\switchcolumn*
\selectlanguage{latin}
Legióne, in Hispánia, 
 sanctórum Mártyrum Cláudii, Lupérci et Victórii, filiórum sancti Marcélli 
 Centuriónis; qui, in persecutióne Diocletiáni et Maximiáni, sub Diogeniáno 
 Præside, jussi sunt decollári.
\switchcolumn
\selectlanguage{english}
At Leon in Spain, the holy martyrs 
 Claudius, Lupercus, and Victorius, the sons of St. Marcellus the centurion. 
 They were condemned to be beheaded by Diogenian, the governor, in the 
 persecution of Diocletian and Maximian.
\switchcolumn*
\selectlanguage{latin}
Lutétiæ Parisiórum 
 sancti Lucáni Mártyris.
\switchcolumn
\selectlanguage{english}
At Paris, St. Lucanus, martyr.
\switchcolumn*
\selectlanguage{latin}
Alexandríæ sanctæ 
 Eutrópiæ Mártyris, quæ, Mártyres vísitans, apprehénsa est, et, cum illis 
 sævíssime cruciáta, réddidit spíritum.
\switchcolumn
\selectlanguage{english}
At Alexandria, the martyr St. 
 Eutropia, who was arrested while visiting the martyrs, and rendered up her 
 soul after being cruelly tortured with them.
\switchcolumn*
\selectlanguage{latin}
Antiochíæ sancti 
 Serapiónis Epíscopi, eruditióne claríssimi.
\switchcolumn
\selectlanguage{english}
At Antioch, St. Serapion, a bishop 
 very celebrated for his learning.
\switchcolumn*
\selectlanguage{latin}
Cápuæ sancti Germáni, 
 Epíscopi et Confessóris, magnæ sanctitátis viri; cujus ánimam, in hora 
 óbitus ejus, ab Angelis in cælum deférri sanctus Benedíctus aspéxit.
\switchcolumn
\selectlanguage{english}
At Capua, St. Germanus, bishop and 
 confessor, a man of great sanctity, whose soul, at the very hour of death, 
 was seen by St. Benedict taken to heaven by angels.
\switchcolumn*
\selectlanguage{latin}
Poténtiæ, in Lucánia, 
 sancti Gerárdi Epíscopi.
\switchcolumn
\selectlanguage{english}
At Potenza in Lucania, St. Gerard, 
 bishop.
\switchcolumn*
\selectlanguage{latin}
\end{paracol}


% ---- martyrology/mart10/mart1031.htm
\needspace{10\baselineskip}
\begin{paracol}{2}
\selectlanguage{latin}
\begin{center}{\color{gregoriocolor} Prídie Kaléndas Novémbris. 
 Luna\dots\ }\end{center}
\switchcolumn
\selectlanguage{english}
\begin{center}{\color{gregoriocolor} The 
 31\textsuperscript{st} Day of 
 October. The\dots\ Day of the Moon.}\end{center}
\end{paracol}

\noindent\begin{tabularx}{\linewidth}{*{19}{>{\centering\arraybackslash}X}}
 \textcolor{gregoriocolor}{a} & \textcolor{gregoriocolor}{b} & \textcolor{gregoriocolor}{c} & \textcolor{gregoriocolor}{d} & \textcolor{gregoriocolor}{e} & \textcolor{gregoriocolor}{f} & \textcolor{gregoriocolor}{g} & \textcolor{gregoriocolor}{h} & \textcolor{gregoriocolor}{i} & \textcolor{gregoriocolor}{k} & \textcolor{gregoriocolor}{l} & \textcolor{gregoriocolor}{m} & \textcolor{gregoriocolor}{n} & \textcolor{gregoriocolor}{p} & \textcolor{gregoriocolor}{q} & \textcolor{gregoriocolor}{r} & \textcolor{gregoriocolor}{s} & \textcolor{gregoriocolor}{t} & \textcolor{gregoriocolor}{u} \\
 10 & 11 & 12 & 13 & 14 & 15 & 16 & 17 & 18 & 19 & 20 & 21 & 22 & 23 & 24 & 25 & 26 & 27 & 28 \\
\end{tabularx}
\vspace{0.5\baselineskip}
\noindent\begin{tabularx}{\linewidth}{*{12}{>{\centering\arraybackslash}X}}
 \textcolor{gregoriocolor}{A} & \textcolor{gregoriocolor}{B} & \textcolor{gregoriocolor}{C} & \textcolor{gregoriocolor}{D} & \textcolor{gregoriocolor}{E} & F & \textcolor{gregoriocolor}{F} & \textcolor{gregoriocolor}{G} & \textcolor{gregoriocolor}{H} & \textcolor{gregoriocolor}{M} & \textcolor{gregoriocolor}{N} & \textcolor{gregoriocolor}{P} \\
 29 & 30 & 1 & 2 & 3 & 4 & 4 & 5 & 6 & 7 & 8 & 9 \\
\end{tabularx}

\begin{paracol}{2}
\selectlanguage{latin}
\lettrine[lines=1]{V}{igília} ómnium 
 Sanctórum.
\switchcolumn
\selectlanguage{english}
\lettrine[lines=1]{T}{he} Vigil of All Saints.
\switchcolumn*
\selectlanguage{latin}
Romæ sanctórum Ampliáti, 
 Urbáni et Narcíssi, quorum méminit sanctus Paulus ad Romános scribens, qui, 
 ob Evangélium Christi, a Judæis et Gentílibus cæsi sunt.
\switchcolumn
\selectlanguage{english}
At Rome, the Saints Ampliatus, 
 Urbanus, and Narcissus, who are mentioned by St. Paul in his Epistle to the 
 Romans. They were put to death by the Jews and Gentiles for the Gospel 
 of Christ.
\switchcolumn*
\selectlanguage{latin}
Constantinópoli sancti 
 Stachis Epíscopi, qui a beáto Andréa Apóstolo primus ejúsdem civitátis 
 Epíscopus ordinátus est.
\switchcolumn
\selectlanguage{english}
At Constantinople, St. Stachis, 
 bishop who was consecrated first bishop of that city by the blessed apostle 
 Andrew.
\switchcolumn*
\selectlanguage{latin}
Apud Augústam 
 Veromanduórum, in Gállia, sancti Quinctíni, civis Románi et ex órdine 
 Senatório viri, qui sub Maximiáno Imperatóre martyrium passus est; cujus 
 corpus post annos quinquagínta quinque, revelánte Angelo, invéntum est 
 incorrúptum.
\switchcolumn
\selectlanguage{english}
At Saint Quentin in France, the 
 martyr St. Quentin, a Roman citizen and senator, who suffered under Emperor 
 Maximian. By the revelation of an angel, his body was found incorrupt 
 after a lapse of fifty-five years.
\switchcolumn*
\selectlanguage{latin}
Medioláni sancti 
 Antoníni, Epíscopi et Confessóris.
\switchcolumn
\selectlanguage{english}
At Milan, St. Antoninus, bishop and 
 confessor.
\switchcolumn*
\selectlanguage{latin}
Ratisbónæ, in Bavária, sancti Wolfgángi Epíscopi.
\switchcolumn
\selectlanguage{english}
At Ratisbon in Bavaria, St. 
 Wolfgang, bishop.
\switchcolumn*
\selectlanguage{latin}
Palmæ, in Majórica ínsula, sancti Alfónsi Rodríguez, Coadjutóris temporális
formáti e Societáte Jesu et Confessóris, humilitáte ac jugi mortificatiónis
stúdio insígnis; quem Leo Décimus tértius, Póntifex Máximus, Sanctórum
fastis adscrípsit.
\switchcolumn
\selectlanguage{english}
At Palma, in the island of Majorca, 
 St. Alphonsus Rodriguez, a lay brother of the Society of Jesus, whom Leo 
 XIII canonized because of his remarkable humility and constant love of 
 mortification.
\switchcolumn*
\selectlanguage{latin}
Romæ Translátio beáti 
 Nemésii Diáconi, et fíliæ Lucíllæ Vírginis, qui octávo Kaléndas Septémbris 
 decolláti sunt.
\switchcolumn
\selectlanguage{english}
At Rome, the translation of blessed 
 Nemesius, deacon, and his daughter, the virgin Lucilla, who were beheaded on 
 the 25th of August.
\switchcolumn*
\selectlanguage{latin}
\end{paracol}

